% !TEX TS-program = lualatex
% !BIB TS-program = biber
% This file must be compiled with LuaLaTeX + biber.
% Background highlighting in the citation report relies on the lua-ul package,
% which only works under LuaLaTeX.
% Run:  lualatex citation_report.tex && biber citation_report && lualatex citation_report.tex && lualatex citation_report.tex


\documentclass{article}

% XeLaTeX + Unicode + font handling
\usepackage{fontspec}
\usepackage{polyglossia}
\usepackage[letterpaper,top=2cm,bottom=2cm,left=3cm,right=3cm,marginparwidth=1.75cm]{geometry}

% Languages: English
\setdefaultlanguage{english}

% Font setup
\setmainfont{Linux Libertine O}
\setsansfont{Linux Libertine O}
\setmonofont{DejaVu Sans Mono}

% Bibliography
\usepackage[backend=biber, style=apa]{biblatex}
\addbibresource{references.bib}

% Math packages
\usepackage{amsmath,amssymb,amsthm,amsfonts,mathtools,bm}
\usepackage{siunitx,physics,chemformula,tensor,hepnames,braket}
\usepackage[version=4]{mhchem}

% Tables and figures
\usepackage{adjustbox,tabularx,graphicx}
\usepackage[ruled,vlined]{algorithm2e}

% Misc
\usepackage{caption,csquotes}

% Layout and links
\usepackage[colorlinks=true,allcolors=blue]{hyperref}


% Custom length
\newlength{\sixtypercentpage}
\setlength{\sixtypercentpage}{0.6\paperheight}


% Packages for citation annotation (color highlighting and PDF comments)
% Citation reports compile with LuaLaTeX so we can use lua-ul, which highlights
% at the node level and works with math, CJK, and fragile commands natively.
\usepackage{xcolor}
\usepackage[color={0.5 0.5 0.5}]{pdfcomment}
\usepackage{graphicx}
\usepackage{luacolor}  % required by lua-ul for colored highlights
\usepackage{lua-ul}
\usepackage{tikz}

% Define colors for highlighting
% hldarkgreen (dark green - strong support)
\definecolor{hldarkgreen}{RGB}{89,171,91}
% hllightgreen (light green - partial support)
\definecolor{hllightgreen}{RGB}{191,255,207}
% hlorange (orange - unsupported)
\definecolor{hlorange}{RGB}{255,189,89}
% hlred (red - contradictory)
\definecolor{hlred}{RGB}{214,105,89}

% Highlighting commands backed by lua-ul's \highLight (node-level, math-safe)
\newcommand{\hldarkgreen}[1]{\highLight[hldarkgreen]{#1}}
\newcommand{\hllightgreen}[1]{\highLight[hllightgreen]{#1}}
\newcommand{\hlorange}[1]{\highLight[hlorange]{#1}}
\newcommand{\hlred}[1]{\highLight[hlred]{#1}}

% Display citekey in bibliography entries (prepended in gray)
\renewbibmacro*{begentry}{%
  \textcolor{gray}{[\thefield{entrykey}]}\space%
}

\begin{document}


% Add ThesisAI logo at the top center
\vspace{-1em}
\begin{center}
\includegraphics[width=0.3\textwidth]{logo.png}
\end{center}

\vspace{0.5em}

\section*{Citation Report}

\noindent\textbf{Title:} MindfulFeed: A Stanford d.School Design Thinking Record for Ethical Social Media Design to Replace Dark Patterns and Protect Cognitive Well Being

\noindent\textbf{Author:} nirwanmohit4@gmail.com

\noindent\textbf{Date:} September 04, 2026

\noindent\textbf{Number of Papers:} 51

\noindent\textbf{Citation Coverage:} 70\%

\noindent\textbf{Strict Evidence:} Yes

\noindent\textbf{Citation Style:} apa

\noindent\textbf{Citation Quality Score:} 97.5\%

\vspace{1em}


\noindent
This document is an automated citation quality report generated by ThesisAI. Each sentence in your document has been analyzed and color-coded based on how well it is supported by the source materials.

\vspace{0.5em}

\noindent
\begin{center}
\begin{tikzpicture}[scale=1.0]

 \fill[hldarkgreen] (0,0) rectangle (7.867572156196943,0.6);

 \node[black] at (3.9337860780984717,0.3) {56.2\%};

 \fill[hllightgreen] (7.867572156196943,0) rectangle (13.64346349745331,0.6);

 \node[black] at (10.755517826825127,0.3) {41.3\%};

 \fill[hlorange] (13.64346349745331,0) rectangle (13.999999999999998,0.6);

\end{tikzpicture}
\end{center}

\vspace{0.5em}

\noindent
The bar shows the distribution of citation quality across all 589 sentences in your document. Green colors (dark and light) indicate well-supported content, representing \textbf{97.5\%} of your document.

\vspace{1em}



\vspace{1em}

\noindent\begin{tabular}{@{}>{\centering\arraybackslash}m{0.08\textwidth}|>{\centering\arraybackslash}m{0.17\textwidth}|m{0.45\textwidth}|>{\centering\arraybackslash}m{0.2\textwidth}@{}}
\textbf{Color} & \textbf{Classification} & \textbf{Explanation} & \textbf{Action Required} \\
\hline
\tikz\fill[hldarkgreen] (0,0) circle (1.2ex); & Fully supported & The sentence is well-supported by the provided sources with clear, direct evidence. & No action needed \\
\hline
\tikz\fill[hllightgreen] (0,0) circle (1.2ex); & Academic synthesis & Valid academic synthesis combining information from multiple sources, or logical inferences drawn from provided evidence. This represents good academic writing and is acceptable practice. Light green indicates the claim is supported through synthesis rather than direct quotation. Common in comparative statements, conclusions, and summary sentences. & No action needed \\
\hline
\tikz\fill[hlorange] (0,0) circle (1.2ex); & Not supported & The sentence lacks support from the provided sources, contains unsupported claims, or is completely unrelated to the sources. & Revision needed \\
\hline
\tikz\fill[hlred] (0,0) circle (1.2ex); & Contradictory & The sentence directly contradicts the provided sources or makes claims that are opposite to what the sources state. & Must correct \\
\end{tabular}

\newpage

\noindent\textbf{Understanding the PDF Comments:} Click on the \textbf{comment icons} in the margins to see detailed explanations for each color rating. These comments reference specific sources by their \textbf{citekey} and \textbf{page number}. Each entry in the \textbf{References} section at the end of this document shows the citekey in gray brackets (e.g., \textcolor{gray}{[5663359]}) at the beginning of the reference, allowing you to easily match the citekeys mentioned in the PDF comments to their corresponding full references. \textbf{Important:} The page numbers shown (e.g., ``p.13'') refer to the actual page number in the PDF file, not the page number printed within the document itself.

\vspace{1em}

\noindent\textbf{How to Improve Orange or Red Sections:} If your document contains orange or red highlighted sentences, consider the following approaches to improve citation quality:

\begin{itemize}
 \item \textbf{Review and Justify:} Carefully review the highlighted sentences and determine whether they genuinely require editing. Sometimes the AI may flag content that is actually acceptable or necessary for your argument.

 \item \textbf{Rewrite or Remove:} Use follow-up prompts in ThesisAI to remove or rewrite specific text parts that are unsupported or contradictory. You can ask the AI to adjust problematic statements to better align with your sources.

 \item \textbf{Add More Literature (Recommended):} The best approach is to find additional literature sources that provide evidence for sections lacking support. Upload these new sources to ThesisAI and regenerate the document. This will improve citation coverage and strengthen your academic work with more comprehensive source material. \textbf{Rule of thumb:} Provide at least 2 papers per page of output. For example, when generating a 35-page document, upload at least 70 papers to ensure sufficient source material for comprehensive citation coverage.

 \item \textbf{Verify Contradictions:} For red-highlighted contradictory statements, carefully review the source material and either correct the statement or reconsider whether the source appropriately supports your argument.
\end{itemize}

\noindent For any questions or feedback about this citation report, please contact us via the support chat at \url{www.thesisai.io}.

\newpage

% Table of Contents
\tableofcontents

\newpage


\section{Contextual Foundations of Ethical Social Media Design}

\subsection{Digital Engagement, Burnout, and Well\-Being}
\label{sec:digital_engagement_burnout} \hldarkgreen{Intensive participation in social media environments exposes users, particularly students and young adults, to multiple psychological risks that are closely related to burnout and diminished well\-being.}\pdfcomment[icon=Note]{Multiple sources explicitly state that intensive social media participation poses psychological risks for students and young adults, linking use to burnout and reduced well-being (6644697 p.11; 9300898 p.6).} \hldarkgreen{Cyberbullying represents a salient form of harmful digital engagement: anonymity in social platforms facilitates abusive behaviour without accountability, and victims frequently report depression, anxiety, social withdrawal, and in severe cases self\-harm as consequences of sustained harassment.}\pdfcomment[icon=Note]{The role of anonymity in enabling cyberbullying and the reported victim outcomes (depression, anxiety, withdrawal, self-harm) are directly described in the provided sources (6644697 p.11).} \hldarkgreen{From an ethical perspective, such hostile interactions breach principles of respect, dignity, and non\-maleficence, while failures in platform moderation and redress mechanisms aggravate harm by leaving affected users without effective remedies.}\pdfcomment[icon=Note]{Sources discuss the ethical breach of cyberbullying (respect, dignity, non-maleficence) and note failures in moderation/redress that exacerbate harm (6644697 p.11; 9300898 p.7).} \hllightgreen{Platform design choices further shape the relationship between engagement and burnout.}\pdfcomment[icon=Note]{The claim that platform design shapes the engagement--burnout relationship is supported by multiple reviews describing design effects on user behaviour and well-being (2404919 p.6; 9300898 p.6).} \hldarkgreen{Many social media systems incorporate infinite scrolling, push notifications, and reward\-based feedback elements such as likes, shares, and comments, intentionally engineered to prolong usage and capture attention through psychological reinforcement mechanisms \parencites[11]{6644697}.}\pdfcomment[icon=Note]{Infinite scrolling, push notifications, and reward-feedback (likes/shares/comments) engineered to prolong usage and exploit reinforcement are explicitly discussed in the sources (6644697 p.11; 2404919 p.16).} \hldarkgreen{Evidence on AI\-enhanced environments shows that machine learning algorithms personalize notifications and content recommendations to maximize engagement, creating compulsive usage patterns and digital dependency that are especially prominent among adolescents and young adults \parencites[6]{9300898}.}\pdfcomment[icon=Note]{Evidence that machine learning personalizes notifications and recommendations to maximize engagement and can produce compulsive usage, especially among adolescents/young adults, is directly reported (9300898 p.6; 2404919 p.8).} \hldarkgreen{Excessive engagement of this type disrupts sleep, reduces academic productivity, and impairs offline social interactions, aligning with burnout\-related symptoms such as cognitive fatigue and social withdrawal \parencites[11]{6644697}[6]{9300898}.}\pdfcomment[icon=Note]{Disruption of sleep, reduced academic productivity, impaired offline interactions and alignment with burnout-related symptoms are noted in the literature (6644697 p.11; 2404919 p.6; 9300898 p.6).} \hldarkgreen{Research on digital well\-being conceptualizes these outcomes within a multidimensional framework that links engagement intensity with psychological, cognitive, social, and physical states.}\pdfcomment[icon=Note]{The multidimensional framework linking engagement intensity to psychological, cognitive, social, and physical states is central to the cited digital well-being reviews (9300898 p.4; 2404919 p.6).} \hldarkgreen{Psychological consequences of AI\-mediated environments include both positive and negative pathways: AI\-based mood tracking and digital therapeutics can support individuals experiencing stress, anxiety, and depressive symptoms, yet algorithmically curated social media feeds that foreground idealized content and social comparison can heighten anxiety, loneliness, and emotional distress.}\pdfcomment[icon=Note]{Sources document both positive AI-mediated supports (mood tracking, digital therapeutics) and negative effects from curated feeds/social comparison that heighten anxiety and distress (9300898 p.6; 9300898 p.3.3.1).} \hllightgreen{Continuous exposure to such curated streams, particularly during formative years, undermines emotional stability and may accelerate burnout trajectories for heavy users \parencites[4]{9300898}.}\pdfcomment[icon=Note]{Reviews link continuous exposure to curated streams with undermined emotional stability and increased risk of burnout among heavy users, a reasonable inference supported by the sources (2404919 p.15; 9300898 p.6).} \hllightgreen{Social and cognitive dimensions compound these pressures.}\pdfcomment[icon=Note]{The notion that social and cognitive dimensions compound pressures from digital engagement is a synthesis supported across the literature (9300898 p.4; 2404919 p.6).} \hldarkgreen{Recommendation algorithms help users form online communities that deliver emotional support and belonging, but they also restructure communication dynamics by prioritizing engagement metrics over meaningful interaction and may reduce face\-to\-face communication \parencites[6]{9300898}.}\pdfcomment[icon=Note]{Recommendation algorithms forming supportive communities while also prioritizing engagement metrics and potentially reducing face-to-face communication is directly described (9300898 p.6; 9300898 p.3.3.2).} \hldarkgreen{Studies of excessive digital engagement associate high usage with reduced attention control, emotional dysregulation, anxiety, depression, sleep disruption, diminished academic performance, workplace distraction, loneliness, and reduced subjective well\-being, while still acknowledging countervailing benefits such as social connectedness and educational access \parencites[6]{2404919}.}\pdfcomment[icon=Note]{Systematic reviews list reduced attention control, emotional dysregulation, anxiety, depression, sleep disruption, diminished academic/work performance, loneliness, and lower subjective well-being alongside benefits like connectedness (2404919 p.6; 2404919 p.12).} \hldarkgreen{This pattern supports a nonlinear view in which engagement is not inherently harmful, but becomes problematic when algorithmic reinforcement erodes autonomy and substitutes essential offline experiences \parencites[6]{2404919}[6]{9300898}.}\pdfcomment[icon=Note]{The nonlinear view, that engagement can be beneficial but becomes problematic when algorithmic reinforcement erodes autonomy and substitutes offline experiences, is explicitly argued in the sources (2404919 p.15; 9300898 p.4).} \hldarkgreen{Organizational contexts illustrate how continuous digital connectivity can culminate in burnout.}\pdfcomment[icon=Note]{Organizational literature in the provided sources shows continuous connectivity in workplaces can culminate in burnout (2404919 p.6; 9300898 p.6).} \hldarkgreen{AI\-supported collaboration tools, virtual meeting systems, and constant notifications blur boundaries between professional and personal life, producing cognitive fatigue, decision exhaustion, reduced creativity, and work\-life imbalance \parencites[6]{2404919}.}\pdfcomment[icon=Note]{AI-supported collaboration tools, virtual meetings, and constant notifications blurring work--life boundaries and producing cognitive fatigue, decision exhaustion, reduced creativity, and imbalance are documented (2404919 p.6; 9300898 p.6).} \hldarkgreen{Similar mechanisms operate in educational settings, where dependence on AI\-assisted learning systems and recommendation algorithms may weaken attention regulation and independent problem\-solving, thereby contributing to academic strain and student burnout \parencites[16]{2404919}.}\pdfcomment[icon=Note]{Similar mechanisms in education, dependence on AI-assisted learning weakening attention regulation and independent problem-solving and contributing to academic strain and student burnout, are discussed (2404919 p.15; 9300898 p.6).} \hllightgreen{These findings underscore that digital engagement intensity, design patterns, and algorithmic incentives jointly determine whether social media participation supports or undermines well\-being.}\pdfcomment[icon=Note]{The synthesis that engagement intensity, design patterns, and algorithmic incentives jointly determine whether participation supports or undermines well-being is supported across multiple sources (9300898 p.4; 2404919 p.12).} \hldarkgreen{The concept of sustainable digital well\-being provides an integrative lens for product decisions in pro\-social platforms.}\pdfcomment[icon=Note]{The concept of sustainable digital well-being as an integrative lens for product decisions on pro-social platforms is a central claim in the reviews (2404919 p.16; 9300898 p.7).} \hldarkgreen{Sustainable approaches aim to optimize human flourishing within AI\-enabled environments by balancing productive technology use with psychological health, recognizing that problems arise when engagement\-maximizing designs compromise autonomy, digital literacy, and ethical safeguards \parencites[6]{2404919}[12]{2404919}.}\pdfcomment[icon=Note]{Sources describe sustainable approaches balancing productive technology use with psychological health and warn that engagement-maximizing designs can compromise autonomy and other safeguards (2404919 p.15; 9300898 p.4).} \hldarkgreen{In practical product development for a pro\-social social app, these insights justify replacing dark patterns and exploitative reinforcement strategies with features that support healthy emotional regulation, attention control, and genuine social connectedness, thereby directly targeting the burnout risks documented in current research \parencites[6]{2404919}[6]{9300898}[11]{6644697}.}\pdfcomment[icon=Note]{Practical product recommendations, replacing dark patterns with features supporting emotional regulation, attention control, and connectedness, are advocated in the sources and linked to reducing burnout risks (9300898 p.7; 6644697 p.11; 2404919 p.16).}

\subsection{Regulatory and Governance Landscape}
\label{sec:regulatory_governance_landscape} \hldarkgreen{Regulatory and governance arrangements form a critical boundary condition for any attempt to embed ethical safeguards and intentional user controls into social media products oriented toward digital well\-being.}\pdfcomment[icon=Note]{The governance-as-boundary thesis is explicitly argued in the policy and governance literature (see 7808872, p.2 and p.4), which frames regulatory and governance arrangements as critical constraints on embedding ethical safeguards in digital platforms.} \hldarkgreen{Existing AI governance initiatives in healthcare and education predominantly focus on data protection, cybersecurity, and algorithmic fairness, while rarely addressing psychological autonomy and emotional stability as explicit regulatory objectives \parencites[2]{7808872}.}\pdfcomment[icon=Note]{Several sources note a technocratic focus in AI governance on data protection, cybersecurity, and algorithmic fairness while psychological autonomy is less attended to (7808872, p.2; 9300898, p.7).} \hldarkgreen{This technocratic emphasis creates a structural gap: the governance of AI is increasingly institutionalized, yet the governance of digital well\-being as a public good remains largely undefined, leaving pro\-social platforms without clear external benchmarks for cognitive\-health protection \parencites[4]{7808872}.}\pdfcomment[icon=Note]{The claim that a technocratic emphasis creates a structural gap, AI governance institutionalized whereas digital well‑being governance remains undefined, is directly stated in the analysis (7808872, p.2; p.4).} \hldarkgreen{Within AI\-enabled occupational health monitoring, privacy, informed consent, and data minimization are treated as foundational ethical requirements.}\pdfcomment[icon=Note]{The treatment of privacy, informed consent, and data minimization as foundational in AI‑enabled occupational health monitoring is explicitly discussed (8760542, p.118).} \hldarkgreen{Computational systems that analyse wearable sensor data, digital communication logs, and behavioural patterns can expose intimate dimensions of emotional well\-being, social interaction, and cognitive functioning, which increases the risk of misuse, unauthorized access, and discrimination without stringent safeguards.}\pdfcomment[icon=Note]{8760542 (p.118) explicitly states that wearable sensors, communication logs, and behavioural analytics can reveal intimate aspects of emotional well‑being and pose risks of misuse, unauthorized access, and discrimination without safeguards.} \hldarkgreen{Ethical governance in such environments demands transparent, revocable consent protocols, strict limits on data retention to what is necessary for legitimate health objectives, and algorithmic transparency that makes predictive processes explainable and auditable.}\pdfcomment[icon=Note]{The need for transparent, revocable consent, strict data‑retention limits, and algorithmic transparency (explainable and auditable predictive processes) is directly recommended in the occupational health monitoring discussion (8760542, p.118).} \hllightgreen{These principles directly inform the design of social applications that monitor or nudge user behaviour: any burnout\-related analytics or adaptive support features must be bounded by consent, minimization, and transparency so that AI remains a supportive wellness mechanism rather than an intrusive control tool.}\pdfcomment[icon=Note]{8760542 (p.118) and 7808872 (p.4) together link these ethical principles to application design, arguing that consent, minimization, and transparency are necessary to keep AI supportive rather than intrusive, this is a reasonable synthesis of those sources.} \hldarkgreen{Algorithmic bias and fairness concerns further constrain responsible design.}\pdfcomment[icon=Note]{Multiple sources identify algorithmic bias and fairness as central constraints on responsible design (8760542, p.118; 9300898, p.7), supporting this concise claim.} \hldarkgreen{In healthcare monitoring, biased training data and institution\-specific models can reproduce inequities across gender, race, age, and professional hierarchy, misclassifying stress levels or overlooking vulnerability in underrepresented groups \parencites[18]{8760542}.}\pdfcomment[icon=Note]{8760542 (p.118) explicitly documents how biased training data and institution‑specific models can reproduce inequities across gender, race, age, and professional hierarchy and misclassify stress or overlook underrepresented groups.} \hldarkgreen{Bias mitigation is framed as a moral obligation, requiring diverse data representation, fairness\-aware model design, continuous auditing, and attention to accessibility so that digital mental health resources remain usable for staff irrespective of literacy, role, or cultural background.}\pdfcomment[icon=Note]{The requirement that bias mitigation be treated as a moral obligation, via diverse data, fairness‑aware design, continuous auditing, and accessibility considerations, is directly advocated in 8760542 (p.118).} \hldarkgreen{Parallel analyses of AI systems across recruitment, healthcare, education, and content moderation similarly document unequal outcomes arising from biased data or flawed design assumptions, with discriminatory results reinforcing social inequalities and eroding trust among marginalized communities \parencites[7]{9300898}.}\pdfcomment[icon=Note]{Cross‑sector documentation of unequal outcomes from biased data/design in recruitment, healthcare, education, and moderation is reported in the literature (9300898, p.7; 0967765 discussion of algorithmic risks), supporting this statement.} \hldarkgreen{For a pro\-social social app, such findings impose a requirement to treat fairness and inclusion as core well\-being objectives, not peripheral technical issues, and to embed systematic bias monitoring into the product development lifecycle \parencites[7]{9300898}[18]{8760542}.}\pdfcomment[icon=Note]{The recommendation that pro‑social apps must treat fairness and inclusion as core objectives and embed bias monitoring in development lifecycles is consonant with the ethical governance recommendations in 8760542 (p.118) and is explicitly linked to organizational responsibility (8760542, p.118).} \hldarkgreen{Empirical work on AI in digital marketing suggests that users' perceptions of fairness are often mediated by whether organizations implement clear accountability mechanisms; surveys of platform users indicate that transparency and complaint channels can materially affect trust in AI interventions on social media \parencites[2]{7161497}.}\pdfcomment[icon=Note]{Empirical work from digital marketing research shows perceptions of fairness are mediated by accountability mechanisms and that transparency and complaint channels affect trust (7161497, pp.10--11; 9300898, p.7 supports transparency affecting trust).} \hllightgreen{This evidence may suggest that governance instruments for digital well\-being should explicitly require accountability processes (such as audits and redress pathways) in addition to technical fairness checks \parencites[2]{7161497}.}\pdfcomment[icon=Note]{Given the empirical findings on accountability and trust (7161497, pp.10--11) and policy literature advocating audits and redress (9300898, p.7), it is reasonable to infer that governance instruments should require accountability processes in addition to technical checks.} \hldarkgreen{Broader policy frameworks articulate risk\-based and rights\-oriented controls that are directly relevant to social media governance.}\pdfcomment[icon=Note]{Broader frameworks articulating risk‑based and rights‑oriented controls relevant to social media governance are documented (4969503, p.4; 7808872, p.4), directly supporting this claim.} \hldarkgreen{The EU Artificial Intelligence Act introduces a risk\-tiered structure and explicitly prohibits manipulative systems that deploy subliminal techniques or exploit vulnerabilities to distort behaviour, although current exceptions leave gaps in the protection of individual identity and raise difficult enforcement questions regarding the boundaries of prohibited manipulation.}\pdfcomment[icon=Note]{The EU AI Act's risk‑tiered structure and prohibition of manipulative systems that use subliminal techniques or exploit vulnerabilities, alongside noted gaps/exceptions and enforcement challenges, are explicitly discussed (4969503, p.4).} \hldarkgreen{Complementary platform regulations such as the Digital Services Act and Digital Markets Act require large platforms to conduct risk assessments that include threats to mental health and democracy, mandate annual evaluations of impacts on fundamental rights, and impose algorithmic transparency obligations related to content moderation, competition, and user rights \parencites[12]{4969503}.}\pdfcomment[icon=Note]{The Digital Services Act and Digital Markets Act requirements, risk assessments including mental health and democracy, annual impact evaluations, and algorithmic transparency obligations, are described in the EU regulatory analyses (4969503, p.4).} \hllightgreen{These instruments collectively signal that dark patterns and addictive design mechanisms are not merely a matter of design ethics but fall under formal regulatory scrutiny when they threaten mental autonomy, cognitive liberty, and mental integrity.}\pdfcomment[icon=Note]{The cited instruments and scholarship link regulatory attention to dark patterns and addictive designs when they threaten mental autonomy and integrity (4969503, p.4; 4807303 on dark patterns), so this is a reasonable synthesis.} \hldarkgreen{Governance perspectives on digital well\-being emphasise that power, agency, and accountability must be re\-conceptualized for intelligent systems.}\pdfcomment[icon=Note]{The claim that governance perspectives call for reconceptualizing power, agency, and accountability for intelligent systems is a central argument of the governance literature cited (7808872, p.2--4).} \hldarkgreen{Current national and international policies, including the EU AI Act, UNESCO's Recommendation on the Ethics of Artificial Intelligence, India's National Digital Health Mission, and National Education Policy 2020, advance AI adoption while offering limited safeguards for the emotional and cognitive effects of continuous digital engagement.}\pdfcomment[icon=Note]{The assessment that the EU AI Act, UNESCO Recommendation, India's NDHM, and NEP 2020 promote AI adoption while offering limited safeguards for emotional/cognitive effects is explicitly stated in the policy review (7808872, p.2; 2404919, p.17).} \hldarkgreen{In response, the Governance of Digital Well\-being Framework proposes five pillars for ethical oversight across education and mental health platforms: ethical accountability, algorithmic transparency, human oversight, psychological safety, and inclusive policymaking \parencites[2]{7808872}.}\pdfcomment[icon=Note]{The Governance of Digital Well‑being (GDW) Framework and its five pillars (ethical accountability, algorithmic transparency, human oversight, psychological safety, inclusive policymaking) are presented in the study (7808872, p.4).} \hldarkgreen{Translating these pillars into concrete policy imperatives yields requirements such as national\-level commissions dedicated to digital well\-being, mandatory AI impact assessments focusing on psychological autonomy and cognitive diversity, and national standards for algorithmic well\-being metrics that move beyond conventional data protection indicators toward measures of attention sustainability, authentic emotional engagement, and autonomy preservation.}\pdfcomment[icon=Note]{Concrete policy imperatives, national commissions, mandatory AI impact assessments focused on psychological autonomy and cognitive diversity, and national algorithmic well‑being metrics, are proposed in the governmental policy recommendations (7808872, p.12).} \hldarkgreen{For institutional actors, governance is operationalized through digital well\-being charters, internal ethics committees for AI oversight, and professional training in algorithmic literacy for educators, clinicians, and administrators \parencites[12]{7808872}.}\pdfcomment[icon=Note]{Institutional operationalization via digital well‑being charters, internal AI ethics committees, and professional training in algorithmic literacy is explicitly recommended (7808872, p.12; p.8).} \hldarkgreen{These mechanisms reorient governance from narrow compliance to shared responsibility for experiential outcomes, aligning with research that calls for human\-centered, well\-being\-oriented AI design where transparency, explainability, user agency, and ethical principles guide system development and evaluation \parencites[7]{9300898}[17]{2404919}.}\pdfcomment[icon=Note]{The shift from narrow compliance to shared responsibility and the alignment with human‑centred, well‑being‑oriented AI design (transparency, explainability, user agency) is reflected in the literature and tied to the cited framework (2404919, p.17; 7808872, p.12).} \hldarkgreen{Evidence from analyses of FinTech and broader AI adoption similarly points to a governance shortfall: emerging technologies create pressure for new regulatory geographies and supervisory models that can respond to privacy, fairness, and accountability risks at scale \parencites[11]{5353824}.}\pdfcomment[icon=Note]{Analyses of FinTech and AI adoption pointing to governance shortfalls and pressure for new regulatory geographies and supervisory models are discussed in the FinTech governance literature (5353824, pp.2, 11--13).} \hllightgreen{That literature appears to caution that without institutional adaptation, new oversight mechanisms, updated legal frameworks, and clearer standards for ethical use, technical fixes alone are unlikely to secure cognitive and emotional protections in digital platforms \parencites[11]{5353824}.}\pdfcomment[icon=Note]{Multiple sources caution that without institutional adaptation, new oversight, updated legal frameworks, clearer ethical standards, technical fixes alone are insufficient to secure cognitive and emotional protections (7808872, p.4; 5353824, p.13; 4590567, p.4), making this a supported synthesis.}

\section{Theoretical Perspectives Informing MindfulFeed}

\subsection{Behavioral Economics and Choice Architecture}
\label{sec:beh_econ_choice_arch} \hldarkgreen{Behavioral economics provides a structured account of why interface details such as button placement, color schemes, wording, and navigation structures can substantially shape user decisions in digital environments.}\pdfcomment[icon=Note]{Empirical reviews provided (e.g., Source 2668071, p.2; Source 4807303, p.2) explicitly state that behavioural economics explains how interface details (button placement, colors, wording, navigation) influence user decision-making.} \hldarkgreen{Minor visual and interactive elements alter the salience of options, influencing which actions users perceive as preferable or effortless, and thereby establishing core conceptual foundations for analysing deceptive or manipulative practices in online systems \parencites[2]{2668071}[12]{5678011}.}\pdfcomment[icon=Note]{The role of visual/interactive elements in altering salience and steering choices is discussed directly in Gray et al. (Source 5678011, p.12) and corroborated by broader literature summarized in Source 4807303, p.2.} \hldarkgreen{Within consent banners and privacy dialogs, designers may employ larger or brightly colored acceptance buttons, preselected boxes, and hidden configuration links to prioritise agreement over refusal, turning neutral choice architecture into interface interference oriented dark patterns that systematically steer users toward tracking or data disclosure choices \parencites[12]{5678011}.}\pdfcomment[icon=Note]{Source 5678011 (p.12 and p. ? sections quoted) explicitly documents larger/bright accept buttons, preselected boxes, hidden links and shows these are used in consent banners as interface interference/dark patterns.} \hldarkgreen{These salience manipulations exploit limited attentional resources and cognitive shortcuts, especially when users are multi tasking or processing information quickly, and can produce routinised, unreflective default behaviours even when explicit nudging elements are reduced \parencites[12]{5678011}[14]{5678011}.}\pdfcomment[icon=Note]{Source 5678011 (p.12 and p.14) links salience manipulations to limited attention, multitasking, cognitive shortcuts and reports 'effect survival' or routinised default behaviours following nudging.} \hllightgreen{For MindfulFeed, these findings delineate specific design moves that must be excluded from consent, notification, and personalization flows if the product is to avoid replicating attention diversion and traffic light metaphor biases that make psychologically costly options less visible or harder to enact \parencites[12]{5678011}.}\pdfcomment[icon=Note]{The need to exclude attention-diverting designs and traffic-light color biases is supported by Gray et al. (Source 5678011, p.12) and design recommendations in Source 4807303 (p.78/2078) ,  this is a reasonable synthesis for MindfulFeed.} \hldarkgreen{Dual process models and nudging theory further clarify the mechanisms that dark patterns employ.}\pdfcomment[icon=Note]{Dual-process models and nudging as explanatory frameworks for dark patterns are explicitly invoked in Source 4807303 (p.2) and Source 2668071 (p.2), supporting this claim.} \hldarkgreen{Kahneman's distinction between fast, intuitive $ \text{System 1} $ and slower, deliberative $ \text{System 2} $ processing shows that users under time pressure or uncertainty rely heavily on heuristics, which increases susceptibility to interface structures that exploit anchoring, loss aversion, default effects, and least effort principles \parencites[2]{2668071}[2]{4807303}.}\pdfcomment[icon=Note]{Source 4807303 (p.2) directly cites Kahneman's System 1/System 2 distinction and links time pressure/heuristics to susceptibility to anchoring, loss aversion, defaults and least-effort effects.} \hllightgreen{Thaler and Sunstein's notion of nudging conceptualises subtle changes in choice architecture that encourage particular behaviours without formally restricting options; these can be deployed for socially beneficial purposes or for organisational gains that sideline user welfare \parencites[2]{2668071}.}\pdfcomment[icon=Note]{Thaler \& Sunstein's nudging concept and its dual use for welfare or organisational gain is covered in Source 2668071 (p.2) and Source 4807303 (p.2), so this is a valid synthesis of those sources.} \hldarkgreen{Empirical work on dark patterns in e commerce websites demonstrates how countdown timers, hidden costs, and misleading consent mechanisms can systematically increase conversion and data sharing, confirming that these architectures actively harness cognitive biases rather than merely presenting information \parencites[2]{4807303}[2]{2668071}.}\pdfcomment[icon=Note]{Source 2668071 (p.2) and Source 4807303 (p.2) cite empirical e-commerce studies (e.g., countdown timers, hidden costs) showing increases in conversions and the harnessing of cognitive biases.} \hldarkgreen{In the specific context of cookie banners and consent management, preselected options and obstructive flows have been shown to reduce privacy preserving choices and nudge users toward acceptance, thereby transforming a nominally free decision into one architecturally determined by interface configuration \parencites[2]{4807303}[3]{5678011}.}\pdfcomment[icon=Note]{Automated and empirical analyses reported in Source 4807303 (p.2 and p.8) and Gray et al. (Source 5678011, p.3 and p.12) document preselected options and obstructive flows reducing privacy-preserving choices in cookie banners.} \hllightgreen{Any behavioral blueprint for MindfulFeed must therefore treat defaults, framing, and friction levels in opt in and opt out pathways as primary levers for ethical design, ensuring symmetric effort for privacy friendly actions and avoiding guilt or loss framing in calls to action \parencites[8]{4807303}[12]{5678011}.}\pdfcomment[icon=Note]{Recommendation to treat defaults, framing, and friction as primary levers, and to ensure symmetric effort, is supported by design guidance in Source 4807303 (p.78/2078) and legal/design discussion in Source 5678011 (p.12); this is a reasonable synthesis for product guidance.} \hldarkgreen{Ethical analyses distinguish between dark patterns and ethical nudges by reference to organisational intent and autonomy outcomes.}\pdfcomment[icon=Note]{The ethical distinction between dark patterns and ethical nudges based on organisational intent and autonomy outcomes is explicitly discussed in Source 2668071 (p.3) and reiterated in Gray et al. (Source 5678011, p.15).} \hldarkgreen{Dark patterns primarily manipulate users to maximise organisational metrics such as engagement, subscriptions, or sales, frequently undermining informed consent, transparency, and perceived control \parencites[2]{2668071}[2]{4807303}.}\pdfcomment[icon=Note]{Source 2668071 (p.2--3) and Source 4807303 (p.2) state that dark patterns are used to maximise organisational metrics (engagement, subscriptions, sales) and can undermine informed consent and autonomy.} \hldarkgreen{In consent interfaces, obstruction, forced action, and interface interference strategies can render rejection cumbersome or degrade service quality for those who decline tracking, thereby placing disproportionate cognitive and experiential burdens on privacy oriented users \parencites[3]{5678011}[12]{5678011}.}\pdfcomment[icon=Note]{Gray et al. (Source 5678011, p.14 and p.12) document obstruction, forced action, and interface interference in consent flows that make rejection cumbersome or impose reduced service, burdening privacy-oriented users.} \hldarkgreen{Research on "bright patterns" and educative nudges suggests limited practical feasibility in commercial settings, since many companies remain structurally incentivised to track users and to deploy architectures that favour data collection; even pro privacy nudges risk reproducing unreflective defaults and perceived lack of control when they rely solely on automatic steering rather than meaningful deliberation \parencites[15]{5678011}.}\pdfcomment[icon=Note]{Source 5678011 (p.15) explicitly argues that 'bright patterns' and educative nudges face limited practical feasibility because companies are structurally incentivised to track, and pro-privacy nudges can reproduce unreflective defaults.} \hllightgreen{For a pro social app such as MindfulFeed, these constraints imply that ethical choice architecture must foreground user comprehension and reversible, low friction controls over behaviourally optimized funnels, integrating educative elements about data consequences without subordinating them to engagement maximisation \parencites[15]{5678011}[3]{5678011}.}\pdfcomment[icon=Note]{The recommendation that ethical choice architecture should prioritise comprehension, reversible low-friction controls and educative elements for a pro-social app is supported by Gray et al. (Source 5678011, p.15) and design guidance in Source 4807303 (p.78), a reasonable synthesis.} \hldarkgreen{Aligning cancellation, consent withdrawal, and configuration pathways with principles of balanced salience, equal optical prominence for accept and reject options, and unambiguous consent satisfies both behavioural evidence on autonomy supportive design and legal expectations that withdrawing consent be as easy as giving it \parencites[12]{5678011}[3]{5678011}[2]{4807303}.}\pdfcomment[icon=Note]{Source 5678011 (p.12 and p.14) cites legal expectation (GDPR Article 7(4)) that withdrawing consent be as easy as giving it, and Source 4807303 (p.78) recommends symmetric design, directly supporting this claim.} \hllightgreen{Recent work on AIGC mediated cultural production suggests an additional layer of complexity: when algorithmic systems act as active co-creators of content and interaction, users' expectations and attentional habits are reshaped in ways that may interact with interface nudges, making it more likely that subtle salience shifts produce routinised acceptance in contexts where AI-generated companions or personalized feeds appear socially normative \parencites[4]{6567535}.}\pdfcomment[icon=Note]{Work on AIGC and cultural production (Source 6567535, pp.2--6) and literature on AI reshaping user expectations (Source 9300898, p.7) support the inference that AI co-creation can reshape norms and attention, interacting with interface nudges.} \hllightgreen{This implies MindfulFeed should consider not only superficial parity in button design, but also how AI-driven content dynamics influence perceived normality and thereby the effectiveness of consent architectures \parencites{6567535}.}\pdfcomment[icon=Note]{Combining Gray et al.'s concerns about visual parity and normative effects (Source 5678011, p.12) with AIGC literature on how content dynamics shape perceived normality (Source 6567535, pp.4--6) makes this a supported synthesis for design implications.}

\subsection{Psychological Needs and Motivation}
\label{sec:psych_needs_motivation} \hldarkgreen{Human motivational dynamics in AI\-mediated environments are shaped by the interaction between fundamental psychological needs and platform design features that systematically address, and sometimes exploit, these needs.}\pdfcomment[icon=Note]{The NGD and related literature explicitly frame motivational dynamics as an interaction between psychological needs and platform design that can amplify or exploit those needs (see 2634982 p.221; 2404919 p.8).} \hldarkgreen{Self\-Determination Theory (SDT) identifies autonomy, competence, and relatedness as core requirements for sustained motivation and well\-being.}\pdfcomment[icon=Note]{Self-Determination Theory and its three basic needs (autonomy, competence, relatedness) are explicitly stated in the sources (2404919 p.8; 2634982 p.14).} \hldarkgreen{Digital platforms explicitly target these needs by configuring personalized experiences that allow individuals to regulate content choices, by embedding gamified feedback such as achievement systems and progress indicators to reinforce competence, and by enabling community formation and social networking that foster relatedness through persistent online interaction.}\pdfcomment[icon=Note]{2404919 directly documents how digital platforms address autonomy via personalization, competence via gamification/feedback, and relatedness via online communities and social networking (2404919 p.8).} \hldarkgreen{Artificial intelligence intensifies this motivational alignment.}\pdfcomment[icon=Note]{2404919 and related syntheses state that AI enhances motivational satisfaction and thus intensifies alignment between platform features and psychological needs (2404919 p.8).} \hldarkgreen{Recommendation systems and conversational agents tailor content and interactions to individual preferences, thereby increasing perceived relevance and satisfaction.}\pdfcomment[icon=Note]{2404919 describes recommendation systems and conversational agents tailoring content and interactions to individual preferences, increasing relevance and satisfaction (2404919 p.8).} \hldarkgreen{Personalized interfaces, predictive text, and adaptive notifications collectively reduce behavioural friction, making it easier for users to continue satisfying needs within the same environment.}\pdfcomment[icon=Note]{2404919 and supporting sources describe personalized interfaces, predictive text, and adaptive notifications as reducing cognitive effort/behavioural friction and promoting continuity (2404919 p.8).} \hldarkgreen{Over time, repeated satisfaction through low\-effort interactions transforms intentional engagement into habitual routines, such as automatically checking smartphones after waking, opening social media during breaks, or browsing recommended videos before sleeping, even when users consciously recognise that their digital engagement has become excessive and difficult to regulate.}\pdfcomment[icon=Note]{2404919 presents the habit-formation pathway with the same examples and notes that such habits operate below conscious awareness despite recognition of excessive engagement (2404919 p.8).} \hllightgreen{For a pro\-social app such as MindfulFeed, these dynamics have direct design implications.}\pdfcomment[icon=Note]{This is a reasonable design-oriented inference supported by sources that link motivational dynamics to practical design implications (2634982 p.217; 2404919 p.8).} \hldarkgreen{SDT explains both the attractiveness and the potential psychological risks of AI\-mediated experiences.}\pdfcomment[icon=Note]{2404919 explicitly argues that SDT accounts for both attractiveness and psychological risks of AI-mediated engagement (2404919 p.8).} \hldarkgreen{Excessive algorithmic personalization may gradually erode genuine autonomy by subtly shaping behavioural options through predictive recommendations, shifting users from active selection to guided patterns of engagement.}\pdfcomment[icon=Note]{2404919 directly warns that excessive algorithmic personalization can erode genuine autonomy by shaping behavioural choices via predictive recommendations (2404919 p.8).} \hldarkgreen{In parallel, persuasive technology principles describe how reminders, social comparison, gamification, reward systems, and behavioural feedback are intentionally configured to influence attitudes and behaviours without overt coercion.}\pdfcomment[icon=Note]{Persuasive Technology Theory and its mechanisms (reminders, social comparison, gamification, rewards, feedback) are listed and characterised in 2404919 (3.6) and related sources (4807303).} \hldarkgreen{AI systems learn which persuasive strategies are most effective for specific individuals, thereby increasing the potency of nudges that anchor users in ongoing interaction cycles.}\pdfcomment[icon=Note]{2404919 states that AI increases persuasive effectiveness by learning which strategies generate maximal behavioural responses for individuals (2404919 3.6).} \hldarkgreen{Uses and Gratifications Theory adds a complementary lens for product decisions.}\pdfcomment[icon=Note]{2404919 and 2634982 discuss Uses and Gratifications Theory as a complementary lens for understanding user motivations and product decisions (2404919 3.5; 2634982 p.217).} \hldarkgreen{Individuals adopt and continue to use digital technologies to satisfy diverse psychological and social needs, including information seeking, entertainment, identity development, emotional regulation, social interaction, professional networking, and education.}\pdfcomment[icon=Note]{2404919 enumerates the diverse needs users seek to satisfy (information, entertainment, identity, emotional regulation, social interaction, professional networking, education) (2404919 3.5).} \hldarkgreen{AI expands the breadth and intensity of gratifications by supplying highly personalized streams of entertainment, adaptive learning experiences, predicted consumer offerings, emotionally responsive conversational agents, and individualized wellness guidance.}\pdfcomment[icon=Note]{2404919 provides examples of AI expanding gratifications through personalized entertainment, adaptive learning, predicted offerings, conversational agents, and wellness guidance (2404919 3.5).} \hldarkgreen{This diversification of gratifications strengthens habitual engagement and demonstrates that problematic use does not arise solely from dark\-patterned interface design, but also from repeated reliance on digital environments to meet emotional and psychological needs that may otherwise remain unmet.}\pdfcomment[icon=Note]{2404919 and 2634982 link diversified/personalized gratifications to strengthened habitual engagement and note that problematic use arises also from reliance on digital environments to meet unmet needs, not only dark patterns (2404919 3.5; 2634982 p.217).} \hllightgreen{Within MindfulFeed's development record, motivational theories therefore function as operational constraints on feature ideation and prototyping.}\pdfcomment[icon=Note]{Sources recommend that motivational theories inform design and intervention priorities; applying them as operational constraints in a product development record is a reasonable synthesis (2634982 p.217; 5375614 p.6).} \hllightgreen{SDT requires that autonomy support be authentic rather than superficial, so personalization controls, recommendation explanations, and adjustable notification regimes must be implemented to preserve self\-directed choice rather than to lock users into pre\-defined paths.}\pdfcomment[icon=Note]{2404919 and 2634982 emphasize authentic autonomy support via transparency, controls, and explanations; recommending personalization controls and adjustable notifications is a supported design inference (2404919 p.8; 2634982 p.16).} \hllightgreen{Competence reinforcement through feedback and progress indicators should be oriented toward meaningful skill building or reflective self\-knowledge instead of raw engagement metrics, while relatedness features need to promote constructive social connection rather than purely comparative or competitive dynamics.}\pdfcomment[icon=Note]{Sources note competence is often provided superficially via gamification and caution against undermining authentic fulfillment; steering feedback toward meaningful skill-building and constructive relatedness is a supported recommendation (2404919 p.8; 2634982 p.14--16).} \hllightgreen{Persuasive technology insights justify embedding goal tracking and behavioural feedback only when they are transparently presented as well\-being tools rather than covert engagement maximizers.}\pdfcomment[icon=Note]{Persuasive technology literature and ethics-focused sources argue for transparency and framing of behaviour-change tools as well-being supports rather than covert engagement maximizers (2404919 3.6; 4807303).} \hllightgreen{Evidence on habit formation and gratifications advises that design sprints within the Empathize and Define phases focus on mapping which psychological needs different personas attempt to satisfy through social media, and how AI\-driven personalization currently shapes those efforts.}\pdfcomment[icon=Note]{Intervention and NGD literature recommend mapping needs across personas and examining how AI personalization shapes gratifications, so focusing design sprints on these mappings is a reasoned synthesis (2634982 p.216--221; 2404919 3.4--3.5).} \hllightgreen{Low\- and mid\-fidelity prototypes should deliberately test alternative ways of meeting autonomy, competence, and relatedness needs with higher transparency and user control, while empirical usability evaluations must capture whether features intended as supportive begin to function as de facto dependency loops.}\pdfcomment[icon=Note]{The recommendation to test prototypes for autonomy/competence/relatedness support and to evaluate whether features create dependency loops follows directly from calls for empirical testing and evaluation in the sources (2634982 p.18; 2634982 p.16).} \hllightgreen{Anchoring product decisions in SDT, Uses and Gratifications, and Persuasive Technology Theory ensures that motivational mechanisms in MindfulFeed are harnessed for sustainable digital well\-being rather than for intensified engagement \parencites[8]{2404919}.}\pdfcomment[icon=Note]{Anchoring decisions in SDT, U\&G, and Persuasive Tech to prioritize wellbeing over raw engagement is a defensible synthesis supported by the cited frameworks and policy recommendations (2634982 p.217; 2404919 p.8; 4590567).} \hldarkgreen{Evidence from domains where AI surveillance and automated monitoring have been deployed suggests a parallel risk: systems that reduce human oversight and accelerate automated detection can increase public mistrust and raise fraught questions about data retention, secondary uses, and accountability \parencites[2]{7025237}.}\pdfcomment[icon=Note]{Surveillance literature documents that automated monitoring reducing human oversight can increase mistrust and raises concerns about retention, secondary uses, and accountability (7025237 p.3--6; 9300898 p.7).} \hllightgreen{Such findings imply that MindfulFeed should not only provide user controls and transparent explanations, but also consider independent audit mechanisms and clear limits on data reuse to avoid replicating the kinds of trust deficits observed in other AI\-heavy contexts \parencites[2]{7025237}.}\pdfcomment[icon=Note]{Sources recommend user controls, transparent explanations, independent audits, and limits on data reuse to address trust deficits in AI-heavy contexts (7025237 p.6--11; 0235650 p.8).}

\subsection{Algorithmic Personalization and Dark Patterns}
\label{sec:alg_personalization_dark_patterns} \hllightgreen{Algorithmic personalization operates as a central mechanism through which dark patterns scale and adapt to individual users in contemporary digital environments.}\pdfcomment[icon=Note]{Sources discuss how algorithmic personalization amplifies manipulative practices and attention capture in digital platforms (9300898 p.5; 2404919 p.8), supporting a synthesis that personalization enables dark patterns to scale and adapt to individuals.} \hldarkgreen{Personalized recommendation systems on social media, streaming services, e commerce platforms, and news aggregators use machine learning to tailor content to inferred preferences, thereby reducing information overload and increasing perceived relevance for many users.}\pdfcomment[icon=Note]{Source 9300898 explicitly states that social media, streaming, e‑commerce and news aggregators use machine learning to tailor content, reducing information overload and increasing perceived relevance (9300898 p.5).} \hllightgreen{At the same time, these systems tend to surface emotionally charged and sensational content because such material drives higher engagement, which has been empirically associated with increased polarization, exposure to misinformation, and emotional exhaustion \parencites[5]{9300898}.}\pdfcomment[icon=Note]{Sources report that recommendation systems can surface emotionally charged/sensational content because it drives engagement and link this dynamic to polarization, misinformation exposure, and emotional exhaustion (9300898 p.5; 0235650 p.4).} \hllightgreen{This engagement maximization dynamic interacts directly with dark pattern design: when attention capture is the dominant objective, interface configurations that privilege certain actions, conceal alternatives, or obstruct exit paths become instrumentally valuable, reinforcing compulsive use and undermining cognitive well being.}\pdfcomment[icon=Note]{The literature links engagement‑maximizing objectives to interface manipulations (interface interference, obstruction, privileging actions) that promote compulsive use and harm well‑being (5678011 p.4; 4807303 p.8), making this a reasonable synthesis.} \hllightgreen{For MindfulFeed, these findings delineate a critical boundary condition for algorithmic features: recommendation logic must be constrained by well being criteria rather than pure engagement metrics, and interface structures around personalized content must avoid exploitative manipulation.}\pdfcomment[icon=Note]{Well‑being oriented AI literature and design recommendations argue aligning recommendation logic to well‑being criteria and avoiding exploitative interfaces (9300898 p.5; 4807303 p.8), supporting this normative design boundary for MindfulFeed.} \hllightgreen{Research on dark patterns identifies a range of mechanisms that are particularly potent when combined with personalization.}\pdfcomment[icon=Note]{Multiple studies identify mechanisms of dark patterns and note their potency when combined with personalization (4807303 p.2; 4692551 p.2), so this general claim is supported as a synthesis.} \hldarkgreen{Foundational taxonomies describe patterns such as trick questions, privacy zuckering, misdirection, disguised ads, roach motels, and sneaking, each defined by specific manipulative strategies \parencites[2]{4807303}[4]{5678011}.}\pdfcomment[icon=Note]{Brignull's original catalogue and subsequent taxonomies list patterns such as trick questions, privacy zuckering, misdirection, disguised ads, roach motels, and sneaking (4807303 p.2).} \hldarkgreen{Subsequent work systematizes these into higher level categories including nagging, obstruction, sneaking, interface interference, and forced action, with explicit attention to how interfaces privilege certain choices or require users to perform unwanted actions to continue accessing functionality \parencites[4]{5678011}.}\pdfcomment[icon=Note]{Gray et al. and related work systematize patterns into higher‑level categories including nagging, obstruction, sneaking, interface interference, and forced action (5678011 p.4).} \hldarkgreen{Empirical experiments demonstrate that dark pattern consent mechanisms can reduce privacy preserving decisions by measurable margins, and that the design of cookie banners strongly determines consent withdrawal rates, indicating that nominally free choices are often architecturally dictated \parencites[2]{4807303}[4]{5678011}.}\pdfcomment[icon=Note]{Empirical findings reported include that dark‑pattern consent mechanisms reduce privacy‑preserving choices (up to \textasciitilde{}23 percentage points) and cookie banner design strongly affects consent withdrawal (4807303 p.2).} \hllightgreen{Within an AI personalized environment, these mechanisms can be targeted at specific user segments whose behavioural data indicate higher susceptibility, making manipulative designs harder to detect and resist.}\pdfcomment[icon=Note]{Sources note that personalization enables targeting of susceptible user segments and complicates detection/resistance to manipulative designs (2668071 p.3; 9300898 p.7), supporting this inferential claim.} \hllightgreen{MindfulFeed's product development record must therefore treat any personalization of consent, settings, or content controls as ethically sensitive, ensuring that tailored interfaces do not exploit user specific vulnerabilities.}\pdfcomment[icon=Note]{Regulatory and ethical literature emphasizes sensitivity around tailored consent/settings and the risk of exploiting vulnerabilities, recommending cautious design of personalized interfaces (5678011 p.15; 9300898 p.7).} \hldarkgreen{Behavioural economics provides further insight into how algorithmic personalization amplifies dark pattern effects.}\pdfcomment[icon=Note]{Behavioral economics is cited as providing key explanatory mechanisms for why dark patterns work, as discussed in the literature review (4807303 p.2).} \hldarkgreen{Dual process theory situates dark patterns as leveraging fast, automatic $ \text{System 1} $ processing while bypassing slower, deliberative $ \text{System 2} $ oversight.}\pdfcomment[icon=Note]{The sources explicitly invoke Kahneman's dual‑process theory to explain how dark patterns exploit fast System 1 processing and bypass System 2 oversight (4807303 p.2).} \hldarkgreen{Anchoring, loss aversion, default effects, and least effort principles explain why preselected options, countdown timers, and asymmetric opt in versus opt out flows reliably steer decisions \parencites[2]{4807303}.}\pdfcomment[icon=Note]{Anchoring, loss aversion, default effects and least‑effort principles are discussed as mechanisms underlying preselected options, countdown timers and asymmetric opt flows (4807303 p.2).} \hllightgreen{When personalization learns which specific anchors or frames work best for an individual, these biases are engaged with increasing precision.}\pdfcomment[icon=Note]{Literature notes that personalization increases persuasive precision by learning effective anchors/frames for individuals (2404919 p.8; 2668071 p.3), supporting this reasonable inference.} \hldarkgreen{Studies note that sophisticated personalization technologies can tailor deceptive interfaces to consumer characteristics, thereby intensifying manipulative practices and making dark patterns more difficult to detect and resist.}\pdfcomment[icon=Note]{Studies and reviews explicitly state that sophisticated personalization can tailor deceptive interfaces to consumer characteristics, intensifying manipulative practices (2668071 p.3; 4807303 p.2).} \hllightgreen{For MindfulFeed, this evidence imposes a design constraint: adaptive algorithms may be used to reduce friction for autonomy supportive actions, but not to calibrate salience, defaults, or emotional framing in ways that systematically favour data disclosure or prolonged engagement.}\pdfcomment[icon=Note]{Well‑being and governance sources recommend using adaptive algorithms to support autonomy while prohibiting calibration of salience/defaults to promote disclosure or prolonged engagement (9300898 p.5; 4807303 p.8), supporting this design constraint as a synthesis.} \hldarkgreen{Ethical and regulatory analyses connect these technical dynamics to privacy and autonomy concerns.}\pdfcomment[icon=Note]{Ethical and regulatory analyses in the provided literature explicitly connect technical personalization dynamics to privacy and autonomy concerns (9300898 p.7).} \hldarkgreen{Contemporary platforms frequently deploy complex privacy policies, confusing consent flows, and defaults that encourage unintended data disclosure, limiting user understanding of data collection, processing, and sharing \parencites[3]{2668071}.}\pdfcomment[icon=Note]{Multiple sources document complex privacy policies, confusing consent flows and default settings that encourage unintended disclosure and limit user understanding (2668071 p.3; 5678011 p.15).} \hldarkgreen{As personalization, targeted advertising, and recommendation systems become more sophisticated, researchers argue that such technologies can amplify manipulative practices by tailoring interfaces against user interests, thereby raising the need for stronger ethical guidelines and regulatory frameworks that prioritise welfare alongside business objectives \parencites[3]{2668071}[7]{9300898}.}\pdfcomment[icon=Note]{Sources argue that as personalization and targeted advertising grow more sophisticated they can amplify manipulative practices and call for stronger ethical and regulatory frameworks (2668071 p.3; 9300898 p.7).} \hldarkgreen{Legal scholarship on consent banners under the GDPR and ePrivacy Directive underscores that interface design is central to lawful processing: consent must be freely given, specific, informed, and unambiguous, and dark patterns that rely on obstruction, forced action, or interface interference are incompatible with these requirements \parencites[4]{5678011}.}\pdfcomment[icon=Note]{Legal scholarship cited explains that GDPR/ePrivacy require freely given, specific, informed and unambiguous consent and that dark patterns relying on obstruction or interface interference conflict with these requirements (5678011 p.4; 0967765 p.4).} \hllightgreen{MindfulFeed's algorithmic personalization strategy must therefore be aligned with data protection by design, avoiding both overt and emergent dark patterns, and subjecting personalized flows to systematic ethical and legal review.}\pdfcomment[icon=Note]{Guidance in policy and ethics literature recommends aligning personalization with data protection by design and subjecting personalized flows to ethical and legal review (9300898 p.7; 4807303 p.8), supporting this normative recommendation.} \hldarkgreen{Evidence on digital well being oriented AI design reinforces this orientation.}\pdfcomment[icon=Note]{The literature on digital well‑being oriented AI design emphasizes and reinforces human‑centered and well‑being objectives (9300898 p.5).} \hldarkgreen{Human centred and well being oriented frameworks argue for systems that prioritize transparency, explainability, user agency, and ethical principles, with success measured in terms of human flourishing rather than engagement alone \parencites[5]{9300898}[7]{9300898}.}\pdfcomment[icon=Note]{Source 9300898 explicitly advocates for frameworks prioritizing transparency, explainability, user agency and measuring success by human flourishing rather than engagement (9300898 p.5--7).} \hllightgreen{In this context, personalization should be deployed to support psychological stability and autonomy, not to deepen dependency.}\pdfcomment[icon=Note]{Well‑being and persuasive technology literature argue personalization should support psychological stability and autonomy rather than deepen dependency (9300898 p.5; 2404919 p.8), making this a supported normative claim.} \hldarkgreen{For the Empathize and Define phases in MindfulFeed's development, this implies collecting user narratives about how recommendation systems and consent interfaces have previously contributed to fatigue or loss of control, and using these insights to specify requirements such as symmetric effort in opt out paths, plain language in personalized prompts, and explicit explanations of why particular content or choices are being surfaced \parencites[2]{4807303}[3]{2668071}[4]{5678011}.}\pdfcomment[icon=Note]{Sources recommend collecting user experiences to inform requirements like symmetric opt‑out effort, plain language, and explanations for personalized prompts (2668071 p.3; 5678011 p.15), directly supporting these development practices.}

\section{Methodological Rationale and Design Thinking Process}

\subsection{Selection of Stanford d.School Framework}
\label{sec:selection_ds_framework} \hlorange{Selection of the Stanford d.School design thinking framework is grounded in methodological choices that prioritise governance, ethical accountability, and psychological outcomes in AI\-mediated environments.}\pdfcomment[icon=Note]{None of the provided sources mention the Stanford d.School framework or claim its selection is grounded in methodological choices prioritising governance, ethical accountability, and psychological outcomes; the sources discuss governance frameworks but not d.School (see 7808872 p.3--8).} \hldarkgreen{Research on digital well\-being governance adopts qualitative, multi\-layered designs that integrate policy analysis, theoretical synthesis, and case illustration to address complex socio\-technical questions, explicitly favouring depth, contextual understanding, and interdisciplinary integration over narrow technical evaluation \parencites[3]{7808872}[7]{7808872}.}\pdfcomment[icon=Note]{Directly supported by Source 7808872 which states the study adopts a qualitative, multi-layered design integrating policy analysis, conceptual synthesis, and comparative case examination and privileges depth and interdisciplinarity (7808872 p.7).} \hlorange{This orientation towards interpretivist, governance\-focused inquiry aligns closely with design thinking's emphasis on understanding lived experience and experiential outcomes before proposing solutions.}\pdfcomment[icon=Note]{While 7808872 emphasizes interpretivist, governance-focused inquiry and understanding human experience, none of the provided sources explicitly link that orientation to design thinking's emphasis on lived experience and experiential outcomes (7808872 p.3--4).} \hldarkgreen{The Governance of Digital Well\-being Framework operationalises this interpretive stance through four analytical dimensions: accountability, transparency, oversight, and well\-being outcomes, all examined via thematic content analysis of policies and case evidence \parencites[8]{7808872}.}\pdfcomment[icon=Note]{Explicitly supported: 7808872 lists the four analytical dimensions (accountability, transparency, oversight, well-being outcomes) and describes thematic content analysis as the analytic method (7808872 p.8).} \hlorange{These dimensions map directly onto the kinds of questions that each phase of the Stanford d.School process can structure in a product development record for MindfulFeed.}\pdfcomment[icon=Note]{No source provided documents a direct mapping between these governance dimensions and the phases of the Stanford d.School process for a product development record for MindfulFeed; d.School is not cited in the sources (see 7808872 for dimensions but no d.School mapping).} \hlorange{For instance, accountability and human oversight can be systematically explored in early empathic research and problem definition, while transparency and psychological safety can be operationalised during ideation and prototyping of features that avoid dark patterns and respect cognitive autonomy.}\pdfcomment[icon=Note]{The particular mapping of accountability/human oversight to empathic research and transparency/psychological safety to ideation/prototyping (linked specifically to d.School phases and dark-pattern avoidance) is not supported by the provided sources, which do not describe d.School phase mappings.} \hldarkgreen{Evidence indicates that contemporary AI deployments in education and mental health are governed primarily through data protection, cybersecurity, and algorithmic fairness, leaving emotional stability, cognitive autonomy, and social connection comparatively under\-specified in policy instruments \parencites[3]{7808872}[4]{7808872}.}\pdfcomment[icon=Note]{Supported by 7808872 which notes policymakers focus on data protection, cybersecurity, and algorithmic fairness while psychological and emotional impacts are often overlooked in AI governance in education and mental health (7808872 p.3--4, p.8).} \hldarkgreen{In parallel, studies of dark patterns document extensive use of manipulative interface strategies, strong behavioural impacts on consent and decision quality, and limited consumer awareness of such practices \parencites[2]{4807303}[2]{5127387}[3]{4692551}.}\pdfcomment[icon=Note]{Supported by multiple dark-pattern sources: 5127387 and 4807303 document prevalence and users' limited awareness, and empirical work (e.g., 4807303, 5127387) reports substantial behavioural impacts of manipulative interface strategies (5127387 p.2; 4807303 p.2072).} \hllightgreen{This combination of governance blind spots and manipulative interface prevalence creates a design problem that is as much ethical and experiential as it is technical.}\pdfcomment[icon=Note]{Reasonable synthesis: 7808872 documents governance blind spots around psychological impacts and dark-pattern literature documents manipulative interfaces, together supporting the claim that the problem is ethical/experiential as well as technical (7808872 p.3--4; 5127387 p.2).} \hllightgreen{A phased design thinking framework provides a structured way to translate these abstract governance deficits and interaction harms into concrete product requirements, prototypes, and test protocols that can be traced through an evidence\-based development record.}\pdfcomment[icon=Note]{Plausible academic synthesis: 7808872 describes governance deficits and the methodological sections emphasize structured analysis, supporting the idea that a phased design process can translate deficits into requirements and traceable artefacts, though no source explicitly names a 'phased design thinking' pipeline for this purpose (7808872 p.7--8; 8760542 p.19 on traceable review procedures).} \hldarkgreen{The methodological approach used in digital well\-being governance research relies on iterative coding, triangulation of policy, academic, and case data, and explicit procedures to ensure credibility, transferability, and dependability \parencites[8]{7808872}.}\pdfcomment[icon=Note]{Directly supported by 7808872 which describes iterative coding, triangulation of policy/academic/case data, and prioritising credibility, transferability, and dependability in its qualitative approach (7808872 p.7--8, 3.4--3.5).} \hldarkgreen{Similarly, systematic reviews of AI\-based burnout interventions employ staged processes of identification, screening, eligibility assessment, and inclusion, supported by structured coding frameworks, thematic clustering, and mixed methods appraisal tools to sustain analytical integrity \parencites[19]{8760542}.}\pdfcomment[icon=Note]{Supported by 8760542 which details PRISMA-based staged processes (identification, screening, eligibility, inclusion), structured coding, thematic clustering, and use of appraisal tools to maintain analytical integrity in systematic reviews (8760542 p.119/Methods section).} \hllightgreen{Selecting the Stanford d.School framework allows MindfulFeed's development documentation to emulate this procedural clarity: each phase can be associated with defined activities, data sources, and analytic steps, such as empathy interviews, cognitive walkthroughs, prototype evaluations, and usability diaries, all logged as traceable artefacts analogous to coded documents in policy analysis or screened studies in systematic reviews.}\pdfcomment[icon=Note]{Partial synthesis: sources describe activities like empathy interviews, cognitive walkthroughs, prototype evaluations and logging qualitative artefacts (7173856; 7808872), so mapping these to phases of a design process as traceable artefacts is a reasonable inference though d.School is not explicitly cited (7173856 p.4; 7808872 p.7--8).} \hldarkgreen{Research on dark patterns highlights mechanistic clarity and non\-overlapping taxonomies as central design principles for classifying manipulative strategies, reflecting a need for exhaustive yet structured categorisation in interface assessment \parencites[2]{4807303}[2]{5127387}.}\pdfcomment[icon=Note]{Directly supported by 4807303 which states taxonomic design principles including mechanistic clarity and non-overlap as central to classifying dark patterns (4807303 p.2072).} \hldarkgreen{Cognitive walkthrough studies on social platforms and mindfulness apps further show that expert evaluators can systematically identify combinations of dark patterns across account creation, content engagement, and termination flows when guided by clear methodological protocols \parencites[3]{4692551}[4]{7173856}.}\pdfcomment[icon=Note]{Supported by 4692551 and 7173856 which describe cognitive walkthroughs and expert evaluators identifying dark-pattern combinations across platform flows, using clear protocols and tasks (4692551 p.3; 7173856 p.4).} \hllightgreen{Embedding these taxonomic insights within the Define and Ideate phases of the d.School process ensures that MindfulFeed's problem statements and concept generation explicitly reference documented dark pattern categories, rather than relying on informal intuitions, and that subsequent prototypes can be evaluated against recognised harms such as obstruction, sneaking, confirm shaming, and roach motel behaviours.}\pdfcomment[icon=Note]{Reasonable synthesis: the dark-pattern taxonomies and enumerated harms (obstruction, sneaking, confirm shaming, roach motel) are documented in 5127387 and 4807303, and integrating these into early design phases to evaluate prototypes is a logical inference (5127387 p.2; 4807303 p.2072).} \hldarkgreen{Finally, governance research on digital well\-being insists that algorithmic systems be accountable to the moral conditions of human life, positioning attention, emotion, and cognition as contested frontiers that require ethical reorientation in power, agency, and accountability \parencites[3]{7808872}[4]{7808872}.}\pdfcomment[icon=Note]{Directly supported by 7808872 which argues governance must make algorithmic systems accountable to moral conditions and identifies attention, emotion, and cognition as contested frontiers requiring ethical reorientation (7808872 p.3--4, p.8).} \hllightgreen{A design thinking framework that begins with empathic investigation and iteratively cycles through prototyping and testing offers a practical methodology to realise this reorientation within a single pro\-social app: each phase can explicitly incorporate governance pillars such as ethical accountability, transparency, human oversight, psychological safety, and inclusive stakeholder input, thereby aligning product\-level decisions with broader policy ambitions for human\-centred AI systems \parencites[4]{7808872}[6]{5375614}.}\pdfcomment[icon=Note]{Supported as a synthesis: 7808872 and 5375614 articulate governance pillars (ethical accountability, transparency, human oversight, psychological safety) and emphasize human-centred AI in education, making it reasonable to infer that an empathic, iterative design process could incorporate these pillars into product-level decisions (7808872 p.4; 5375614 p.6).}

\subsection{Comparison with Alternative Design Frameworks}
\label{sec:comparison_alt_frameworks} \hllightgreen{Evaluating the suitability of the selected design thinking framework requires comparison with alternative methodological approaches that have been applied to digital well\-being, dark patterns, and AI governance.}\pdfcomment[icon=Note]{The sources present multiple methodological options (qualitative governance frameworks in 7808872 and HCI/field methods in 7173856 and 4692551), so comparing a design thinking framework with alternatives applied to digital well‑being, dark patterns, and AI governance is a reasonable synthesis (see 7808872 p.7 and 7173856 p.4).} \hldarkgreen{Qualitative, multi\-layered policy analysis designs integrate policy review, conceptual synthesis, and comparative case examination to construct governance frameworks such as the Governance of Digital Well\-being Framework.}\pdfcomment[icon=Note]{Source 7808872 explicitly describes a qualitative, multi‑layered design that integrates policy analysis, conceptual synthesis, and comparative case examination to construct a governance framework (7808872 p.7).} \hldarkgreen{These approaches privilege contextual depth and interpretive analysis of accountability, transparency, oversight, and well\-being outcomes rather than structured product development activities \parencites[7]{7808872}[8]{7808872}.}\pdfcomment[icon=Note]{Source 7808872 states the qualitative approach privileges depth and interpretive analysis around accountability, transparency, oversight, and well‑being outcomes (7808872 p.8), matching the claim.} \hlorange{In contrast, the Stanford d.School process used in MindfulFeed explicitly organises activities around empathic research, problem definition, ideation, prototyping, and testing, which more directly supports the creation and iterative evaluation of concrete interface artefacts, consent flows, and dark\-pattern countermeasures within a product development record.}\pdfcomment[icon=Note]{None of the provided sources mention the Stanford d.School process being used in MindfulFeed or describe that specific organisation of activities for MindfulFeed; no direct support in the listed sources.} \hllightgreen{Systematic review methodologies provide another point of comparison.}\pdfcomment[icon=Note]{Systematic review methodologies are discussed in the sources (e.g., PRISMA usage in 8760542), so stating they provide another point of comparison is a supported synthesis (8760542 p.119).} \hldarkgreen{PRISMA\-guided reviews of AI\-based burnout interventions employ staged identification, screening, eligibility, and inclusion procedures, coupled with structured coding frameworks, thematic clustering, and formal quality appraisal tools such as the Mixed Methods Appraisal Tool.}\pdfcomment[icon=Note]{Source 8760542 documents PRISMA‑guided reviews using staged identification, screening, eligibility, inclusion procedures and quality appraisal with the Mixed Methods Appraisal Tool (8760542 p.119).} \hldarkgreen{These designs achieve methodological transparency, replicability, and quantitative documentation of trends in intervention effectiveness and predictive accuracy, but they focus primarily on synthesising existing empirical studies rather than generating and testing new interface solutions.}\pdfcomment[icon=Note]{Source 8760542 argues PRISMA and related review designs deliver transparency, replicability and quantitative synthesis while focusing on synthesising existing empirical studies rather than producing new interface solutions (8760542 p.119).} \hlorange{For MindfulFeed, such systematic review procedures are informative for evidence gathering and risk mapping, yet they lack the phase\-based structure for producing and iterating low\-, mid\-, and high\-fidelity prototypes that is central to the product brief \parencites[19]{8760542}.}\pdfcomment[icon=Note]{The claim about MindfulFeed specifically and its product brief structure is not supported by the provided sources; none mention MindfulFeed or its product brief.} \hldarkgreen{Within human--computer interaction and user experience research, expert\-led cognitive walkthroughs represent a distinct methodological paradigm.}\pdfcomment[icon=Note]{Source 4692551 documents expert‑led cognitive walkthroughs in HCI/UX as a methodological paradigm applied to dark pattern research (4692551 p.3).} \hldarkgreen{Cognitive walkthrough studies on large social platforms conduct task\-based evaluations across content engagement, account creation, and termination flows to identify combinations of dark patterns and analyse how these combinations affect user experience and engagement.}\pdfcomment[icon=Note]{Source 4692551 describes task‑based cognitive walkthroughs across content engagement, account creation, and termination flows on platforms to analyse dark pattern combinations and their effects (4692551 p.3).} \hldarkgreen{Sampling of multiple platforms and pilot testing of procedures ensure coverage of diverse interaction modalities and refinement of feature categorisation and evaluator guidelines.}\pdfcomment[icon=Note]{Source 4692551 details sampling multiple platforms and conducting pilot testing to refine procedures and evaluator guidelines (4692551 p.3).} \hllightgreen{This approach yields detailed diagnostic insight into existing interface harms and dark pattern configurations, yet it is not inherently organised around participatory ideation or prototype generation for alternative, autonomy\-supportive designs \parencites[2]{4692551}[3]{4692551}.}\pdfcomment[icon=Note]{Source 4692551 indicates cognitive walkthroughs yield diagnostic insights into interface harms and dark pattern configurations, and it does not frame them as participatory ideation/prototyping methods (4692551 p.3), supporting a synthesized claim.} \hlorange{Integrating cognitive walkthrough techniques into the Test phase of the d.School framework therefore enhances evaluative rigour while preserving the overall innovation\-oriented trajectory of the product development record.}\pdfcomment[icon=Note]{No source documents integrating cognitive walkthroughs specifically into the d.School Test phase for MindfulFeed; the combination with d.School is not supported by the provided materials.} \hldarkgreen{Multi\-method field studies of mindfulness applications illustrate another alternative framework.}\pdfcomment[icon=Note]{Source 7173856 presents multi‑method field studies of mindfulness applications as an alternative framework (7173856 p.4, methodology section).} \hldarkgreen{Over several weeks, combined heuristic evaluation, surveys, diaries, and user interviews are used to examine dark patterns and user perceptions in popular apps.}\pdfcomment[icon=Note]{Source 7173856 describes over several weeks using heuristic evaluation, surveys, diaries, and user interviews to examine dark patterns and user perceptions in mindfulness apps (7173856 p.4--5).} \hldarkgreen{Heuristic evaluation applies an established dark pattern taxonomy, including work by Brignull and subsequent projects, while surveys and diaries provide longitudinal data on engagement, discontinuation, and user attitudes.}\pdfcomment[icon=Note]{Source 7173856 reports using heuristic evaluation with an established dark pattern taxonomy (integrating Brignull and others) and collecting survey/diary longitudinal data on engagement and discontinuation (7173856 p.4--5).} \hldarkgreen{This multi\-stage design mirrors certain aspects of empathic and evaluative research, but it is structured primarily for understanding existing apps rather than for guiding the stepwise construction of a new pro\-social platform with explicit regulatory and theoretical anchors \parencites[2]{7173856}[4]{7173856}.}\pdfcomment[icon=Note]{Source 7173856 notes the multi‑stage design mirrors empathic/evaluative research and is structured to understand existing apps rather than to stepwise construct a new pro‑social platform (7173856 p.4; methodology).} \hlorange{For MindfulFeed, diary and interview components can be embedded within Empathize and Test phases, yet the overarching design thinking framework remains more suitable for mapping desirability, feasibility, and ethical constraints into iterative prototypes.}\pdfcomment[icon=Note]{The specific claim about embedding diary/interview components into MindfulFeed's Empathize and Test phases and that design thinking is more suitable is not supported by the provided sources (no mention of MindfulFeed).} \hldarkgreen{Dark pattern research itself introduces taxonomic and automated detection frameworks that could, in principle, be treated as design methodologies.}\pdfcomment[icon=Note]{Multiple sources describe taxonomic and automated detection frameworks for dark patterns (e.g., 4807303/IRJAEM describing taxonomy and DarkScan; 2832310 and 4807303 discuss detection approaches) supporting the claim (4807303 p.2072).} \hldarkgreen{Comprehensive taxonomies emphasise mechanistic clarity, non\-overlap, and exhaustiveness, organising dozens of patterns into a small number of top\-level categories.}\pdfcomment[icon=Note]{Source 4807303/IRJAEM outlines taxonomic design principles emphasizing mechanistic clarity, non‑overlap, and exhaustiveness in organising patterns into top‑level categories (4807303 p.2072; taxonomic principles section).} \hldarkgreen{Automated detection pipelines such as DarkScan combine static DOM analysis, text classification, visual detection, and dynamic behavioural tracing to identify manipulative strategies and to quantify behavioural and psychological impacts on decision quality, trust, and well\-being.}\pdfcomment[icon=Note]{Source 4807303 describes automated detection pipelines like DarkScan that combine static DOM analysis, text classification, visual detection, and dynamic behavioural tracing to detect manipulative strategies (4807303 p.2072).} \hllightgreen{These frameworks are optimised for classification, measurement, and policy recommendation rather than for participatory co\-design of alternative interaction models.}\pdfcomment[icon=Note]{The sources (4807303 and related detection literature) show these frameworks are optimised for classification, measurement, and policy recommendation; concluding they are not primarily for participatory co‑design is a reasonable inference (4807303 p.2072).} \hllightgreen{When compared with the d.School structure, taxonomic and detection approaches function as analytical backbones that inform problem framing, constraint identification, and success metrics, but they do not replace the need for phases dedicated to ideation, prototyping, and iterative user testing \parencites[2]{4807303}.}\pdfcomment[icon=Note]{The sources support that taxonomies/detection approaches inform problem framing, constraints, and metrics (4807303 p.2072), so concluding they complement but do not replace ideation/prototyping phases is a defensible synthesis.} \hldarkgreen{Finally, interaction criticism and legal\-ethical analyses of consent banners embody interpretive design\-research frameworks focused on uncovering emergent dark patterns and evaluating compliance with consent requirements.}\pdfcomment[icon=Note]{Source 5678011 discusses interaction criticism and legal‑ethical analysis of consent banners as interpretive design‑research frameworks focused on emergent dark patterns and compliance evaluation (5678011 p.15).} \hldarkgreen{Such work reveals previously uncategorised strategies like consent walls and reduced service, discusses bright patterns and educative nudges, and interrogates the burdens placed on end users by current consent regimes.}\pdfcomment[icon=Note]{Source 5678011 explicitly identifies consent walls and reduced service as previously uncategorised strategies, discusses bright patterns/educative nudges, and interrogates burdens on end users (5678011 p.15).} \hllightgreen{This perspective highlights power asymmetries and regulatory constraints, yet it stops short of prescribing a procedural pathway for generating and validating alternative interfaces at scale.}\pdfcomment[icon=Note]{Source 5678011 highlights power asymmetries and regulatory constraints in consent systems and notes the work does not offer a full procedural pathway for generating/validating alternative interfaces at scale, supporting this synthesis (5678011 p.15).} \hlorange{Within MindfulFeed's development record, interaction criticism is incorporated as an evaluative lens in later phases, while the d.School framework provides the operational sequence through which evidence, ethics, and user narratives are translated into concrete design decisions, prototypes, and documented test outcomes \parencites[15]{5678011}.}\pdfcomment[icon=Note]{There is no evidence in the provided sources that a development record for 'MindfulFeed' exists or that interaction criticism is incorporated into its later phases alongside a d.School operational sequence; this specific project claim is unsupported.}

\section{Empathize Phase: Understanding Digital Burnout and User Needs}

\subsection{Qualitative Exploration of Digital Burnout}
\label{sec:qual_explore_burnout} \hldarkgreen{Empathic investigation of digital burnout within MindfulFeed's product development record requires qualitative methods capable of capturing how manipulative design strategies and algorithmic environments shape emotional stability, cognitive autonomy, and social connectedness.}\pdfcomment[icon=Note]{Source 7808872 (p.8) defines well‑being outcomes including emotional stability, cognitive autonomy, and social connectedness and describes a qualitative, interpretive approach suitable for capturing how algorithmic and design features affect these experiential dimensions.} \hldarkgreen{Governance research on digital well\-being adopts thematic content analysis guided by four dimensions, namely accountability, transparency, oversight, and well\-being outcomes, using interpretive coding procedures to link administrative and technological design with experiential consequences for users' cognition and emotion \parencites[7]{7808872}[8]{7808872}.}\pdfcomment[icon=Note]{Source 7808872 (p.8) explicitly states a thematic content analysis guided by four dimensions (accountability, transparency, oversight, well‑being outcomes) and iterative coding linking administrative/technological design to experiential outcomes.} \hllightgreen{This analytical structure directly informs the qualitative exploration strategy for MindfulFeed: semi\-structured data collection must explicitly probe how users attribute responsibility for harmful engagement loops, how visible and explainable recommendation and notification systems appear, what mechanisms they experience as human oversight or redress, and which emotional or social outcomes they associate with prolonged platform use.}\pdfcomment[icon=Note]{Source 7808872 (p.7--8) outlines the analytical framework linking governance design to experiential outcomes; recommending semi‑structured probes about responsibility, transparency, and oversight is a reasonable methodological adaptation supported by that framework.} \hldarkgreen{Dark pattern research in social media and educational technologies underscores that manipulative interface combinations can exacerbate distraction during learning and exposure to harmful content, lowering student achievement and undermining the educational potential of platforms that might otherwise support wellbeing \parencites[3]{4692551}.}\pdfcomment[icon=Note]{Source 4692551 (p.3) cites Himawan and Rosy describing how dark patterns in social media and educational tech can distract learners, increase exposure to harmful content, and lower student achievement, undermining educational potential.} \hldarkgreen{Cognitive walkthrough studies across major social media services such as YouTube, LinkedIn, WhatsApp, and Telegram show that UI/UX experts, when guided by structured tasks, can systematically identify combinations of dark patterns in content creation, engagement, account creation, and account termination flows, revealing how these patterns interact to shape user behaviour and overall experience \parencites[2]{4692551}[3]{4692551}.}\pdfcomment[icon=Note]{Source 4692551 (p.2--3) reports cognitive walkthroughs by UI/UX experts across YouTube, LinkedIn, WhatsApp, and Telegram using structured tasks to identify dark pattern combinations and their interactions affecting user behaviour.} \hllightgreen{These findings justify including expert walkthroughs alongside user interviews in the Empathize phase, with each walkthrough documented as a qualitative artefact that maps specific pattern clusters to hypothesised burnout drivers such as heightened cognitive load, difficulty disengaging, or obstructed account exit.}\pdfcomment[icon=Note]{Source 4692551 (p.3) documents expert cognitive walkthroughs, and Source 7173856 (p.4--5) integrates expert reviews with user interviews; combining expert walkthroughs with interviews and mapping pattern clusters to hypothesised burnout drivers is a supported synthesis of these methods.} \hldarkgreen{Multi\-method fieldwork on mindfulness applications provides an additional template for exploring burnout\-related experiences in app ecosystems that explicitly address psychological health.}\pdfcomment[icon=Note]{Source 7173856 (p.4) describes a multi‑method fieldwork approach on mindfulness apps, providing a clear template for exploring burnout‑related experiences in apps that address psychological health.} \hldarkgreen{Over nine weeks, heuristic evaluation, surveys, diaries, and user interviews were combined to examine how dark patterns in Calm, Headspace, and Insight Timer influence user perceptions and engagement, with ethical clearance obtained in advance and standardized actions such as sign\-up and deletion used to identify manipulative designs.}\pdfcomment[icon=Note]{Source 7173856 (p.4) reports a nine‑week study combining heuristic evaluation, surveys, diaries, and interviews on Calm, Headspace, and Insight Timer, notes ethical clearance, and describes standardized actions (sign‑up, deletion) to identify manipulative designs.} \hldarkgreen{Diary studies over five days allowed participants to use assigned apps flexibly while logging usage, activities, reasons for non\-usage, and perceptions of trust and engagement, thereby supporting familiarity with the app and sensitivity to subtle manipulative cues \parencites[4]{7173856}.}\pdfcomment[icon=Note]{Source 7173856 (p.5) specifies a five‑day diary stage where participants used assigned apps flexibly while logging usage, activities, reasons for non‑usage, and perceptions, intended to foster familiarity and sensitivity to subtle manipulative cues.} \hldarkgreen{Subsequent interviews used a three\-part structure that first elicited general insights on mindfulness and design features, then presented expert review videos of dark patterns to prompt reflection, and finally employed guided application walkthroughs where users shared screens and discussed account deletion or unsubscribing processes in real time.}\pdfcomment[icon=Note]{Source 7173856 (p.5) describes a three‑part interview structure: general discussion on mindfulness and design, viewing expert review videos to prompt reflection, and guided application walkthroughs including account deletion/unsubscribe tasks.} \hldarkgreen{Data analysis followed an inductive content analysis procedure, with repeated reading of transcripts, coding of similar information, and theme construction around perception, motivation, behaviour, and dark patterns, revealing user types (skeptics, interested, experienced) and documenting correlations between dark pattern exposure and negative behaviours such as fleeing the application \parencites[5]{7173856}.}\pdfcomment[icon=Note]{Source 7173856 (p.5) reports inductive content analysis with repeated transcript reading, coding, theme construction (perception, motivation, behaviour, dark patterns), identification of user types (skeptics, interested, experienced), and links between dark pattern exposure and negative behaviours like fleeing the app.} \hllightgreen{Adapting this structure to MindfulFeed, the qualitative exploration of digital burnout will similarly combine diaries of social media use, video\-based expert walkthroughs of problematic flows, and semi\-structured interviews focused on attentional strain, emotional exhaustion, and perceived loss of control.}\pdfcomment[icon=Note]{Source 7173856 (p.4--5) supports combining diaries, expert walkthroughs, and semi‑structured interviews; extending those methods to probe attentional strain, emotional exhaustion, and perceived loss of control is a reasonable methodological adaptation.} \hldarkgreen{Normative analysis of algorithmic addictive designs directed at older adults identifies aADs as systems that use predictive analytics and reinforcement loops to subtly shape behaviour, eroding self\-directed decision making and impairing autonomy in ways comparable to coercive contracts that are invalid in law.}\pdfcomment[icon=Note]{Source 4590567 (p.4) characterises algorithmic addictive designs (aADs) as using predictive analytics and reinforcement loops that erode self‑directed decision‑making and impair autonomy, and draws an analogy to coercive contracts in legal terms.} \hldarkgreen{Excessive technology use is linked to mental health issues, reduced productivity, and increased risks of substance abuse, yet regulatory models often treat aADs as isolated design flaws rather than structural market failures, suggesting that qualitative inquiries must explicitly address how users understand their own responsibility versus systemic manipulation in burnout trajectories.}\pdfcomment[icon=Note]{Source 4590567 (p.4) links excessive technology use to mental health harms, productivity losses, and substance‑use risks and critiques regulatory models that treat aADs as isolated design flaws rather than systemic market failures, motivating qualitative inquiry into user responsibility vs systemic manipulation.} \hldarkgreen{Behaviourally informed risk disclosure concepts such as concise labels indicating addiction risk scores, typical time commitments, and available disengagement tools provide a lens for interview questions about whether existing platforms communicate engagement risks adequately and how such labels might alter subjective experiences of burnout \parencites[4]{4590567}.}\pdfcomment[icon=Note]{Source 4590567 (4.2.1) proposes behaviourally informed 'aADs Labels' featuring addiction risk scores, typical time‑commitment indicators, and summaries of disengagement tools, directly supporting the described disclosure concepts.} \hldarkgreen{Culturally responsive regulation perspectives further emphasize that autonomy, care, and responsibility are interpreted differently across contexts, which argues for sampling diverse user groups and exploring how intergenerational expectations, social isolation, or rural marginalization intersect with digital burnout narratives \parencites[5]{4590567}.}\pdfcomment[icon=Note]{Source 4590567 (p.5) argues for culturally responsive regulation, noting that autonomy, care, and responsibility vary across contexts and recommending sampling diverse groups to explore intergenerational expectations, social isolation, and rural marginalization.} \hllightgreen{Integrating these governance and dark pattern insights, MindfulFeed's qualitative exploration will generate empathy maps, coded themes, and illustrative quotes that document how manipulative architectures, opaque recommendations, and inadequate disclosure contribute to burnout, thereby grounding subsequent Define phase problem statements in systematically analysed experiential evidence \parencites[2]{4807303}[2]{2668071}[7]{7808872}[5]{7173856}.}\pdfcomment[icon=Note]{Sources 2668071, 7808872, and 7173856 provide methodological and governance foundations (thematic coding, qualitative frameworks, and multi‑method dark pattern studies) that reasonably support producing empathy maps, coded themes, and quotes to ground Define‑phase problem statements.} \hllightgreen{Recent policy and media analyses further suggest that transparency obligations and explicit content labelling may alter how users make sense of algorithmic curation; in particular, mandatory disclosures and clearer markers for algorithmic personalization are likely to affect perceptions of accountability and editorial responsibility \parencites[9]{0235650}.}\pdfcomment[icon=Note]{Source 0235650 (p.8--9) and 7808872 (p.7--8) discuss transparency obligations, labeling, and the role of disclosures; inferring that such measures may change how users perceive algorithmic curation and accountability is a supported policy inference.} \hllightgreen{Such provisions, which aim to make algorithmic influence more visible and to require human oversight in editorial workflows, will be considered as prompts in interviews and as contextual variables when interpreting diaries and walkthrough artefacts \parencites[9]{0235650}.}\pdfcomment[icon=Note]{Source 0235650 (p.9) describes legal transparency requirements (e.g., AI Act labeling) and human editorial responsibility, and treating these provisions as interview prompts and contextual variables is a reasonable methodological choice supported by those sources.}

\subsection{User Interview Profiles and Triggers}
\label{sec:user_interview_profiles_triggers} \hllightgreen{Interview sampling for MindfulFeed's empathic work can operationalise profiles that mirror established qualitative studies on youth social media use and mindfulness apps.}\pdfcomment[icon=Note]{This statement reasonably synthesises material from youth social media qualitative work (Bambili study, Source 1161579) and mindfulness-app mixed-method findings (Source 7173856) that produced user profiles; the claim is an inference combining both literatures.} \hldarkgreen{A purposive sampling strategy is appropriate for selecting information rich participants capable of articulating detailed experiences of digital fatigue, dark patterns, and emotional distress, as purposive sampling is explicitly recommended for qualitative designs seeking depth over breadth in youth social media contexts.}\pdfcomment[icon=Note]{Source 1161579 explicitly describes purposive sampling as appropriate for selecting information-rich participants in qualitative youth social media studies and justifies depth-over-breadth sampling (p.5).} \hldarkgreen{Target groups can include students and young adults in dense educational settings similar to Bambili, where social media is an unavoidable component of daily life and interaction, ensuring diversity of experiences in terms of academic use, entertainment, communication, and psychological impacts such as cyberbullying and misinformation.}\pdfcomment[icon=Note]{Source 1161579 identifies students and young adults in Bambili as the target population and reports social media as an unavoidable component of youth life with themes including academic use, entertainment, communication and issues like cyberbullying and misinformation (p.4--5).} \hldarkgreen{Semi structured interviews provide a flexible yet focused format for exploring burnout triggers.}\pdfcomment[icon=Note]{Source 1161579 states semi-structured interviews offer flexibility while ensuring focus for exploring core questions and sensitive topics (p.5), supporting their use to explore burnout triggers.} \hldarkgreen{Interview guides should use open ended questions that invite narratives about motivations, routines, and challenges in social media use, while enabling probing into sensitive topics like emotional strain, distraction, and manipulative experiences, paralleling procedures shown to be effective for discussing cyberbullying, online scams, and distress in conversational settings.}\pdfcomment[icon=Note]{Source 1161579 describes interview guides using open-ended questions that invite narratives and enable probing into sensitive issues such as cyberbullying and emotional distress (p.5), matching the procedures described.} \hldarkgreen{Each interview can be planned for approximately $15$ to $20$ minutes to maintain concentration and reduce respondent fatigue, with audio recording used for accurate transcription and subsequent thematic analysis, following manual coding practices that keep researchers closely connected to emerging patterns in subjective experience \parencites[5]{1161579}.}\pdfcomment[icon=Note]{Source 1161579 reports interviews lasting 15--20 minutes, audio-recorded for transcription, with manual coding to stay close to emerging patterns (p.5); Source 7173856 also uses recording and transcription in thematic analysis (p.5).} \hllightgreen{Profile construction within the MindfulFeed record can draw directly on mixed methods findings from mindfulness application research.}\pdfcomment[icon=Note]{Source 7173856 is a mixed-methods study of mindfulness apps whose findings (diary, interviews, heuristic evaluation) provide profiles and insights that can reasonably inform profile construction for MindfulFeed (p.3--5).} \hldarkgreen{Users of mental health and mindfulness apps have been categorised into skeptics, interested, and experienced types, differentiated by their intrinsic and extrinsic motivations and associated behaviours.}\pdfcomment[icon=Note]{Source 7173856 explicitly identifies three user types ,  skeptics, interested, and experienced , differentiated by motivations and behaviours (p.5, Findings section).} \hldarkgreen{Skeptics predominantly exhibit extrinsic motivations and may be more sensitive to manipulative strategies, while experienced users show primarily intrinsic motivations and persist despite encountering dark patterns.}\pdfcomment[icon=Note]{Source 7173856 reports skeptics as having primarily extrinsic motivations and experienced users as primarily intrinsically motivated who persist despite manipulative designs (p.5, Findings).} \hldarkgreen{Interested users, with mixed motivations, often report confusion and a need for guidance, making them particularly vulnerable to complex, rigid, or pressured interaction flows.}\pdfcomment[icon=Note]{Source 7173856 describes 'interested' users as having mixed motivations, often getting lost in complex apps and needing guidance, which makes them vulnerable to confusing or pressured flows (p.8).} \hllightgreen{These three motivational profiles can be repurposed for MindfulFeed's interviews, with targeted prompts and scenarios designed for each type to capture distinct burnout triggers and reactions to persuasive or manipulative design.}\pdfcomment[icon=Note]{Repurposing the three motivational profiles for targeted interview prompts is a reasonable methodological adaptation inferred from the user-typing in Source 7173856 and aligns with qualitative targeting described in Source 1161579 (p.5).} \hldarkgreen{Diary informed interview structures further refine trigger identification.}\pdfcomment[icon=Note]{Source 7173856 used a diary stage prior to interviews to refine understanding and sensitivity to design nuances, supporting the claim that diary-informed interviews refine trigger identification (p.5).} \hldarkgreen{In the referenced mindfulness study, participants first completed a five day diary capturing app usage, activities, experiences, and reasons for non usage, fostering familiarity and sensitivity to design nuances.}\pdfcomment[icon=Note]{Source 7173856 describes a five-day diary capturing app usage, activities, experiences, and reasons for non-usage, and notes this fostered familiarity and sensitivity (p.5).} \hldarkgreen{Subsequent interviews used a three part format: initial discussion of general app usage and dark patterns, viewing of expert review videos to elicit reflections on negative impacts, and guided walkthroughs where participants shared screens and demonstrated account deletion or unsubscribing processes.}\pdfcomment[icon=Note]{Source 7173856 details a three-part interview: discussion of general app usage and dark patterns, viewing expert review videos to prompt reflections, and guided app walkthroughs including account deletion/unsubscribing demonstrations (p.5).} \hllightgreen{Adopting this sequence for MindfulFeed allows interviews to surface triggers at multiple levels: narrative reflections on burnout, reactions to observed dark patterns, and real time commentary on navigation difficulties, obstructions, or confusing flows in termination and configuration tasks.}\pdfcomment[icon=Note]{This is a reasonable synthesis/inference from the three-part sequence in Source 7173856 showing interviews elicit narrative reflections, reactions to observed dark patterns, and real-time walkthrough commentary (p.5).} \hldarkgreen{Analytical procedures for these interviews should follow inductive content analysis.}\pdfcomment[icon=Note]{Source 7173856 states the researchers chose an inductive content analysis approach for interview data (p.5, Data analysis).} \hldarkgreen{Semi structured interviews are transcribed, read repeatedly to identify similar and relevant information, and coded into categories that then form themes such as perception, motivation, behaviour, and dark patterns, as implemented in mindfulness app research \parencites[5]{7173856}.}\pdfcomment[icon=Note]{Source 7173856 describes transcribing semi-structured interviews, repeatedly reading transcripts to identify similar information, coding into categories and forming themes such as perception, motivation, behaviour, and dark patterns (p.5).} \hllightgreen{Themes can then be mapped to specific burnout triggers, including cognitive overload from notifications, emotional strain linked to rigid progress streaks, or frustration arising from short trial periods that prevent habit formation and understanding of benefits; in prior work, users reported that $7$ to $14$ day trials were insufficient relative to habit formation estimates, leading to confusion and disengagement.}\pdfcomment[icon=Note]{Source 7173856 and its findings discuss dual triggers like notifications and streaks and explicitly report that 7--14 day trials were insufficient for habit formation leading to confusion/disengagement; cognitive overload from notifications is also discussed in related sources (7173856 p.8; 7811174 on emotional triggers).} \hldarkgreen{Recognising dual constraints, where features like streaks or notifications are beneficial for some users but induce stress or guilt for others, is essential for interpreting heterogeneous trigger responses and for identifying design elements that must be configurable rather than universal \parencites[8]{7173856}.}\pdfcomment[icon=Note]{Source 7173856 explicitly recognises dual constraints where features (streaks, notifications) help some users but cause stress or guilt for others and recommends configurability (p.8, Conclusions).} \hldarkgreen{To maintain methodological trustworthiness, interview data for MindfulFeed should be triangulated with other artefacts such as diaries and expert walkthrough observations, echoing multi method approaches that combine heuristic evaluation, survey, diary, and interviews to understand dark patterns' impact on perception and engagement \parencites[4]{7173856}.}\pdfcomment[icon=Note]{Source 7173856 employed and advocates triangulation across heuristic evaluation, survey, diary, and interviews to understand dark patterns' impact (p.3--5), supporting triangulating interviews with diaries and walkthroughs.} \hllightgreen{Credibility, transferability, and dependability can be enhanced through detailed contextual descriptions of participant demographics, replicable interview guides, and transparent coding schemes, in line with governance oriented thematic content analysis that emphasises accountability, transparency, oversight, and well being outcomes in AI mediated environments \parencites[8]{7808872}.}\pdfcomment[icon=Note]{Source 7808872 and Source 716.... (7173856) discuss enhancing credibility, transferability, dependability via contextual descriptions, replicable guides and transparent coding and link thematic analysis to accountability, transparency, oversight and well-being outcomes in governance-oriented analyses (7808872 p.8; 7173856 p.5).}

\section{Define Phase: Framing the MindfulFeed Design Challenge}

\subsection{Synthesis of Empathize Insights into Design Criteria}
\label{sec:synth_empathize_design_criteria} \hllightgreen{Empathic work on digital well\-being and dark patterns provides concrete dimensions that can be translated into design criteria for MindfulFeed.}\pdfcomment[icon=Note]{The claim is a reasonable synthesis: governance and dark-pattern literature provide dimensions and design recommendations that can be translated into product criteria (see 7808872 p.3, p.7; 4807303 p.8).} \hldarkgreen{The Governance of Digital Well\-being research operationalises experiential concerns through four analytical dimensions: accountability, transparency, oversight, and well\-being outcomes, each explicitly focused on emotional stability, cognitive autonomy, and social connectedness in AI\-mediated environments \parencites[3]{7808872}[8]{7808872}.}\pdfcomment[icon=Note]{Source 7808872 explicitly defines the four analytical dimensions (accountability, transparency, oversight, well‑being outcomes) and links them to emotional stability, cognitive autonomy, and social connectedness (7808872 p.8).} \hllightgreen{These dimensions become top\-level requirement categories in the Define phase.}\pdfcomment[icon=Note]{This is a reasonable inference from the framework-to-practice logic in 7808872 and the recommendations to operationalise governance categories into implementation phases (7808872 p.7--8; 4807303 p.8).} \hllightgreen{Interface and feature concepts must be specified in ways that allow later traceability: for every planned flow or control, the product record should document which accountability mechanism is visible to users, how transparency about AI operations or data use is provided, which human oversight or feedback channel is available, and which well\-being outcomes the feature is intended to protect.}\pdfcomment[icon=Note]{Sources recommend documenting user journeys and specifying how accountability, transparency, oversight and well‑being are manifested in interfaces (4807303 p.8; 7808872 p.3, p.8), supporting this traceability recommendation as a synthesis.} \hllightgreen{Methodological procedures used in governance studies and dark pattern detection further shape how empathize insights are converted into actionable criteria.}\pdfcomment[icon=Note]{Both governance-method literature and dark‑pattern research discuss methodological procedures that shape translation of insights into criteria (7808872 p.3--4; 4692551 p.103; 4807303 p.2), so this is a valid synthesis.} \hldarkgreen{Qualitative policy research emphasises multi\-source triangulation, detailed contextual descriptions, and transparent coding to ensure credibility, transferability, and dependability \parencites[8]{7808872}.}\pdfcomment[icon=Note]{Source 7808872 explicitly emphasises multi-source triangulation, detailed contextual descriptions, and transparent coding to ensure credibility, transferability, and dependability in qualitative policy research (7808872 p.8, 3.5).} \hldarkgreen{Dark pattern work likewise stresses mechanistic clarity, non\-overlapping taxonomies, and exhaustive categorisation to support reliable identification and impact assessment \parencites[2]{4807303}[3]{4692551}.}\pdfcomment[icon=Note]{4807303 sets out taxonomic design principles including mechanistic clarity, non-overlap, and exhaustiveness for dark-pattern taxonomies (4807303 p.3), directly supporting this claim; 4692551 also discusses taxonomy and methodology (4692551 p.103).} \hllightgreen{In the product development record, these principles imply that user narratives, diary entries, and expert walkthrough observations must be coded against a clearly defined dark pattern taxonomy and the four GDW dimensions, yielding a matrix of "pain points by mechanism" that directly informs design criteria.}\pdfcomment[icon=Note]{Combining methodological guidance from governance (7808872 p.8) and dark-pattern taxonomy recommendations (4807303 p.3) reasonably supports coding narratives against taxonomies and GDW dimensions to form a pain‑points-by-mechanism matrix.} \hldarkgreen{For instance, if cognitive walkthroughs reveal obstruction and interface interference during account termination flows on major platforms, this is recorded as a specific problem mechanism to be reversed through symmetric, low\-friction exit designs \parencites[3]{4692551}[2]{4807303}.}\pdfcomment[icon=Note]{4807303 provides an empirical example of slow, obstructive cancellation flows (long cancellation times, high abandonment) and recommends symmetric, low‑friction exit designs (4807303 p.8; p.78--78 excerpt showing cancellation metrics).} \hllightgreen{Empathize findings from mindfulness app research add granular detail on how manipulative designs translate into negative user experience patterns.}\pdfcomment[icon=Note]{Mindfulness/HCI literature documents how commercialised mindfulness apps can employ manipulative designs that harm user experience (3667034 p.8; 7173856 p.2), supporting this general statement as a synthesis.} \hlorange{In a nine\-week multi\-method study, heuristic evaluation, diaries, surveys, and semi\-structured interviews documented how Brignull\-type dark patterns such as bait and switch, confirm shaming, disguised ads, forced continuity, friend spam, and hidden cost appear within Calm, Headspace, and Insight Timer, and how users respond behaviourally and emotionally to these techniques \parencites[2]{7173856}.}\pdfcomment[icon=Note]{No provided source documents a specific nine‑week multi‑method study listing those Brignull patterns within Calm, Headspace, and Insight Timer with the described methods; the exact study details are not present in the sources.} \hlorange{Inductive content analysis organised interview data into themes of perception, motivation, behaviour, and dark patterns, and identified user types (skeptics, interested, experienced) with distinct sensitivities and motivations \parencites[5]{7173856}.}\pdfcomment[icon=Note]{The specific claim about inductive content analysis yielding those exact themes and user types is not documented in the provided sources; no source presents that exact coding outcome or typology.} \hllightgreen{Translating these insights into MindfulFeed's Define phase requires explicit pattern\-to\-criterion mapping.}\pdfcomment[icon=Note]{This is a reasonable methodological inference: translating empirical pattern findings into explicit pattern‑to‑criterion mappings follows from the frameworks and recommendations in 7808872 and 4807303 (7808872 p.7; 4807303 p.8).} \hllightgreen{Each observed pattern category becomes a banned mechanism, and each theme becomes a design objective: avoid guilt\-inducing copy for "confirm shaming", avoid hidden paywalls or short trial periods that undermine understanding, and ensure that any continuity or reminders are presented as configurable support tools rather than enforced obligations.}\pdfcomment[icon=Note]{4807303 and related dark‑pattern literature recommend prohibiting certain manipulative mechanisms and converting harms into concrete design objectives (4807303 p.8; 5678011 p.15), so framing observed patterns as banned mechanisms and themes as objectives is a supported synthesis.} \hlorange{Diary and guided walkthrough structures used in the mindfulness study also suggest practical execution details for MindfulFeed.}\pdfcomment[icon=Note]{The assertion that diary and guided walkthrough structures from 'the mindfulness study' provide practical execution details for MindfulFeed is not directly supported by the provided sources; the specific study and its procedures are not cited in the materials given.} \hlorange{Participants kept five\-day diaries capturing usage, activities, reasons for non\-usage, and trust and engagement perceptions, then participated in interviews that combined general reflection, viewing of expert dark pattern reviews, and live screen\-sharing for account deletion tasks \parencites[4]{7173856}.}\pdfcomment[icon=Note]{No provided source documents participants keeping five‑day diaries and performing the exact interview and live screen‑sharing procedure described; these detailed study procedures are not present in the cited materials.} \hllightgreen{For the Define phase, the product record should synthesise these types of artefacts into empathy maps and trigger lists, indicating which interface moments most frequently correlate with cognitive overload, emotional strain, or perceived loss of control.}\pdfcomment[icon=Note]{Synthesising artefacts into empathy maps and trigger lists is a reasonable inference from qualitative best practices and the recommendation to document user journeys and contextual descriptions (7808872 p.8; 4807303 p.8).} \hllightgreen{Each trigger is then formalised as a design criterion, such as limiting notification volume, eliminating countdown timers tied to mental health content, or ensuring that account deletion and data export can be completed in a small, documented number of steps.}\pdfcomment[icon=Note]{Examples like limiting notifications, removing countdown timers, and ensuring manageable account deletion steps align with concrete recommendations in 4807303 (symmetric journeys, plain language, eliminate manipulative timers) and governance concerns in 7808872 (7808872 p.3; 4807303 p.8).} \hldarkgreen{Governance analysis of digital well\-being underscores that AI systems can erode agency by quantifying learning patterns and mental states into optimisable data streams, exposing individuals to surveillance, behavioural nudging, and psychological conditioning \parencites[3]{7808872}.}\pdfcomment[icon=Note]{7808872 explicitly warns that AI systems can erode agency by quantifying learning patterns and mental states into optimisable data streams, exposing users to surveillance, nudging, and conditioning (7808872 p.3; p.8).} \hllightgreen{Combined with dark pattern research showing measurable impacts on decision quality, trust, and psychological well\-being, and demonstrating how consent interfaces can reduce privacy\-preserving choices by significant margins \parencites[2]{4807303}, empathize insights must be synthesized into criteria that explicitly guard against agency erosion.}\pdfcomment[icon=Note]{4807303 reports measurable impacts of dark‑pattern exposure on decision quality and psychological outcomes (4807303 p.2), and 5678011/4807303 discuss how consent interfaces reduce privacy‑preserving choices, supporting this combined synthesis (4807303 p.2; 5678011 p.4).} \hldarkgreen{This involves requiring plain language in calls to action, enforcing symmetric complexity for opt\-in and opt\-out journeys, banning consent walls and reduced\-service strategies that tie access to data disclosure, and mandating that AI\-driven personalisation be explainable and adjustable by users \parencites[15]{5678011}[8]{4807303}.}\pdfcomment[icon=Note]{4807303 explicitly recommends plain language, symmetric complexity for opt‑in/opt‑out, banning consent walls/reduced‑service tactics, and making personalization explainable/adjustable (4807303 p.8).} \hllightgreen{Design criteria framed in this way create a direct bridge between users' burnout narratives and concrete interface and feature constraints, providing a structured problem definition that subsequent Ideate, Prototype, and Test activities can operationalise and evaluate.}\pdfcomment[icon=Note]{This is a reasonable synthesis: framing design criteria to connect user burnout narratives to interface constraints follows from the governance-to-design arguments and recommendations in 7808872 and 4807303 (7808872 p.3; 4807303 p.8).} \hldarkgreen{Evidence from related empirical work on short\-form video consumption suggests that very high exposure can be associated with reduced attention spans and frequent, repeated app access, an outcome that may parallel the cognitive overload seen with manipulative interface patterns; such findings imply we should treat high\-frequency content loops and reward scheduling as risk vectors when specifying well\-being constraints \parencites[10]{2464839}.}\pdfcomment[icon=Note]{Empirical work on short‑form video consumption (2464839) documents high exposure, frequent app access, and a significant association with reduced attention spans and frequent app openings, supporting treating high‑frequency loops as risk vectors (2464839 p.6--9, p.10).}

\subsection{Primary Persona Development}
\label{sec:primary_persona_development} \hldarkgreen{Primary persona construction for MindfulFeed must be anchored in empirically derived motivational and behavioural profiles rather than hypothetical user archetypes.}\pdfcomment[icon=Note]{The methodological discussion and findings in the diary/interview study recommend grounding persona construction in empirical motivational and behavioural profiles (see 7173856, p.5).} \hldarkgreen{Mixed methods research on mindfulness applications identifies three distinct user types, differentiated by the nature of their motivations and their interaction with dark patterns: skeptics, interested users, and experienced users.}\pdfcomment[icon=Note]{The mixed-methods study explicitly reports three user types (skeptics, interested, experienced) differentiated by motivations and interactions with dark patterns (7173856, p.5).} \hldarkgreen{Skeptics exhibit predominantly extrinsic motivations, display higher sensitivity to manipulative strategies, and are more likely to disengage or "flee" when confronted with opaque flows or pressured calls to action.}\pdfcomment[icon=Note]{7173856 describes skeptics as primarily extrinsically motivated, sensitive to manipulative strategies, and prone to fleeing or disengaging when confronted with opaque or pressured flows (7173856, p.5).} \hldarkgreen{Interested users show a blend of intrinsic and extrinsic motivations, frequently reporting confusion and a need for guidance in navigating complex app features.}\pdfcomment[icon=Note]{The study characterises interested users as showing a mix of intrinsic and extrinsic motivations and reports confusion and need for guidance in navigating app features (7173856, p.5).} \hldarkgreen{Experienced users are primarily intrinsically motivated, maintain more regular usage, and are more inclined to purchase applications they perceive as beneficial and use frequently.}\pdfcomment[icon=Note]{7173856 reports experienced users as primarily intrinsically motivated, with more regular usage and a higher likelihood of purchase linked to perceived benefit and frequent use (7173856, p.5).} \hllightgreen{Defining a primary persona for MindfulFeed requires selecting the profile whose vulnerabilities and needs align most closely with the project brief's emphasis on counteracting dark patterns and digital burnout.}\pdfcomment[icon=Note]{This is a reasonable prescriptive inference from the study's emphasis on matching persona choice to project aims (empirical profiling in 7173856, p.5) and governance goals (7808872, p.8).} \hldarkgreen{Findings indicate that the presence of manipulative designs correlates with negative behaviours, particularly abandonment, and that motivational composition shapes how different users interpret and tolerate these experiences \parencites[5]{7173856}.}\pdfcomment[icon=Note]{The research finds manipulative designs correlate with negative behaviours such as abandonment, and that motivational composition shapes interpretation and tolerance of these experiences (7173856, p.5).} \hllightgreen{For an evidence based product record focused on ethical, autonomy supporting interaction, the interested user type provides the most critical baseline.}\pdfcomment[icon=Note]{Selecting the interested user as a baseline is an interpretive recommendation combining the empirical typology (7173856, p.5) with product design priorities; sources support this as a defensible choice.} \hllightgreen{This profile is simultaneously attracted to perceived benefits and susceptible to confusion, making it an appropriate focal point for design criteria that must reduce cognitive load, clarify flows, and avoid exploitative reinforcement.}\pdfcomment[icon=Note]{The characterization (attracted to benefits but susceptible to confusion) synthesises findings on motivations and reported confusion (7173856, p.5) with design recommendations to reduce cognitive load (4807303, p.2078).} \hldarkgreen{Semi structured interviews in mindfulness studies demonstrate that such users can articulate nuanced perceptions of design features and dark patterns when prompted with examples and guided walkthroughs, offering rich qualitative material for persona detail \parencites[2]{7173856}[5]{7173856}.}\pdfcomment[icon=Note]{7173856 describes semi-structured interviews that used examples and guided walkthroughs enabling participants to articulate nuanced perceptions of design features and dark patterns (7173856, p.5).} \hldarkgreen{Operationalising this primary persona requires systematic integration of diary entries, interview codes, and guided walkthrough observations.}\pdfcomment[icon=Note]{The study operationalises persona building by integrating diary entries, interview coding, and guided walkthrough observations (7173856, p.5, data analysis section).} \hldarkgreen{In the referenced research, a five day diary phase captured usage, activities, experiences, and reasons for non usage, fostering familiarity and sensitivity to subtle manipulative cues in Calm, Headspace, and Insight Timer \parencites[7173856][4]{7173856}.}\pdfcomment[icon=Note]{7173856 documents a five-day diary phase capturing usage, activities, experiences and reasons for non-usage to foster familiarity and sensitivity to manipulative cues (7173856, p.5).} \hldarkgreen{Subsequent interviews followed a three part structure: open discussion of mindfulness and app usage, exposure to expert dark pattern review videos, and live screen sharing to examine routine tasks and account deletion or unsubscribing processes \parencites[2]{7173856}[4]{7173856}.}\pdfcomment[icon=Note]{The three-part interview structure (general discussion, expert dark-pattern review videos, guided walkthroughs including account deletion) is described in 7173856 (7173856, p.5).} \hldarkgreen{Content analysis grouped transcribed data into codes and themes under perception, motivation, behaviour, and dark patterns, revealing how interested users interpret design elements, experience stress, and respond to pressured flows \parencites[7173856][5]{7173856}.}\pdfcomment[icon=Note]{7173856 reports content analysis grouping transcribed data into codes/themes under perception, motivation, behaviour, and dark patterns and links these to how interested users interpret design elements (7173856, p.5).} \hllightgreen{For MindfulFeed, persona attributes should be derived through equivalent coding, with fields such as typical usage context, motivational drivers, trust thresholds, burnout triggers, and reactions to obstructive or guilt framed interfaces populated directly from these thematic clusters.}\pdfcomment[icon=Note]{This prescriptive recommendation follows directly from the paper's methodological approach (use of coding to derive persona fields in 7173856, p.5) and is a reasonable synthesis for MindfulFeed.} \hldarkgreen{Trustworthiness requirements in governance oriented qualitative research further constrain persona development.}\pdfcomment[icon=Note]{7808872 emphasises trustworthiness requirements in governance-oriented qualitative research and notes these constrain development of outputs such as personas (7808872, p.8).} \hldarkgreen{Thematic content analysis in AI governance work emphasises credibility through data triangulation, transferability through detailed contextual descriptions, and dependability via transparent coding procedures \parencites[8]{7808872}.}\pdfcomment[icon=Note]{7808872 explicitly highlights credibility via triangulation, transferability through contextual descriptions, and dependability via transparent coding in thematic analyses (7808872, p.4).} \hllightgreen{Applying these criteria, the primary persona must be documented with explicit references to source artefacts: diary excerpts demonstrating feelings of anxiety or irritation when access is restricted, interview segments describing difficulty identifying manipulative designs due to "dark pattern blindness," and walkthrough observations revealing confusion or frustration in account management tasks \parencites[7173856][5]{7173856}.}\pdfcomment[icon=Note]{Combining the governance trustworthiness criteria (7808872, p.4) with the diary/interview artefacts described in 7173856 supports documenting personas with explicit source artefacts (7173856, p.5).} \hldarkgreen{Research notes that repeated exposure can normalise dark patterns, leading users to justify or accept them despite residual discomfort, which is a critical experiential dimension that should appear in the persona's attitudinal profile and burnout risk assessment \parencites[7173856][5]{7173856}.}\pdfcomment[icon=Note]{7173856 notes normalization of dark patterns with repeated exposure and users' tendency to justify or accept them despite discomfort, which the authors identify as important for attitudinal profiles (7173856, p.5).} \hllightgreen{Dark pattern detection and governance studies also inform the behavioural envelope within which the primary persona operates.}\pdfcomment[icon=Note]{This synthesises dark-pattern detection literature (e.g., taxonomies and tools in 4807303) and governance studies (7808872) to argue they inform behavioural envelopes for personas.} \hldarkgreen{Cognitive walkthroughs on YouTube, LinkedIn, WhatsApp, and Telegram show that experts, when guided by structured tasks, can identify combinations of dark patterns across content engagement and account creation or termination, and that these combinations shape user experience and interaction difficulty \parencites[3]{4692551}.}\pdfcomment[icon=Note]{4692551 documents cognitive walkthroughs on YouTube, LinkedIn, WhatsApp, and Telegram by UI/UX experts identifying dark-pattern combinations across engagement and account flows and their effect on experience (4692551, p.103).} \hldarkgreen{Comprehensive taxonomies and automated detection pipelines such as DarkScan document measurable impacts of dark pattern exposure on decision quality, platform trust, and psychological well being, and recommend symmetric design principles ensuring that opting out, cancelling, or exercising privacy preferences is no more complex than the corresponding opt in flows \parencites[2]{4807303}[8]{4807303}.}\pdfcomment[icon=Note]{4807303 presents a comprehensive taxonomy, the DarkScan automated detection pipeline, reports measurable impacts on decision quality, trust and well-being, and recommends symmetric design principles (4807303, pp.2072, 2078).} \hllightgreen{MindfulFeed's primary persona therefore must explicitly encode expectations for symmetric pathways, plain language, and low friction autonomy safeguards, treating any deviation from these expectations as a probable burnout or disengagement trigger.}\pdfcomment[icon=Note]{This recommendation logically follows from DarkScan/policy recommendations to adopt symmetric pathways and plain language to protect autonomy (4807303, p.2078) and is a reasonable design inference for MindfulFeed.} \hldarkgreen{Finally, governance analysis frames attention, emotion, and cognition as contested frontiers in AI mediated environments, with digital well being defined as a problem of accountability, transparency, oversight, and well being outcomes \parencites[7808872][8]{7808872}.}\pdfcomment[icon=Note]{7808872 frames attention, emotion and cognition as contested frontiers in AI-mediated environments and defines digital well-being via accountability, transparency, oversight, and well-being outcomes (7808872, pp.3,8).} \hllightgreen{The primary persona for MindfulFeed should thus be structured around these four dimensions, specifying how the interested user type perceives responsibility for harmful engagement loops, understands or fails to understand algorithmic operations, experiences available oversight or complaint mechanisms, and links interface encounters to emotional stability and social connectedness.}\pdfcomment[icon=Note]{Structuring a persona around the four governance dimensions (7808872, p.8) and applying them to the interested user type is a reasonable synthesis of the governance framework and the user-typology (7173856, p.5).} \hllightgreen{Anchoring persona development in this multi dimensional, empirically coded profile ensures that subsequent ideation and prototyping activities remain aligned with documented user needs and governance principles rather than speculative assumptions \parencites[7808872][8]{7808872}[2]{4807303}[3]{4692551}.}\pdfcomment[icon=Note]{This is a defensible synthesis: anchoring persona development in empirically coded profiles aligns with governance principles and empirical methods described in 7808872 (pp.3--4), 4807303 (pp.2072--2078) and 7173856 (p.5).} \hldarkgreen{This body of work also signals that digitalisation requires coordinated technological, legal and social measures and that algorithmic audit and transparency are often necessary to prevent new forms of discrimination or privacy harm \parencites[2]{0967765}.}\pdfcomment[icon=Note]{Several sources emphasise that digitalisation requires coordinated technical, legal and social measures and call for algorithmic audit/transparency to prevent discrimination and privacy harms (0967765, pp.2--4; 4807303, pp.2072--2078).} \hldarkgreen{Examples from national systems further suggest that decentralised data exchange layers and immutable logging can increase traceability and accountability, which complements the requirement for explicit traceable artefacts in persona records \parencites[2]{0967765}.}\pdfcomment[icon=Note]{0967765 describes national examples (e.g., Estonia's X-Road with immutable logging features) showing how decentralised exchange and logging increase traceability and accountability (0967765, pp.3--4).}

\subsection{Strategic How\-Might\-We Questions}
\label{sec:strategic_how_might_we} \hllightgreen{Strategic How\-Might\-We questions in the product development record must translate documented harms and governance concerns into open, solution\-oriented prompts that remain tightly anchored in empirical findings on dark patterns, cognitive bias, and digital well\-being governance.}\pdfcomment[icon=Note]{This is a reasonable synthesis of the empirical dark-pattern literature and governance framing found in the provided sources (e.g., the empirical taxonomy and impacts in 4807303, p.2072 and governance framing in 7808872, p.3).} \hldarkgreen{Dark pattern taxonomies grounded in mechanistic clarity, non\-overlap, and exhaustiveness show that manipulative strategies such as sneaking, obstruction, misdirection, forced action, confirm shaming, roach motels, and drip pricing deliberately exploit cognitive tendencies including default bias, scarcity bias, framing effects, and loss aversion to steer users into decisions they did not intend to make \parencites[2]{5127387}[2]{4807303}.}\pdfcomment[icon=Note]{Source 4807303 presents a taxonomy built on mechanistic clarity, non-overlap, exhaustiveness and lists patterns (sneaking, roach motel, misdirection, forced action, confirm shaming, drip pricing) and links these to cognitive biases (default effect, scarcity, framing, loss aversion) (4807303, p.2072).} \hllightgreen{Within the Define phase, strategic questions must therefore ask how interface flows and consent dialogs can be architected so that these mechanisms are structurally impossible rather than merely discouraged, while still satisfying legal requirements that consent be freely given, specific, informed, unambiguous, and revocable under the GDPR and ePrivacy regimes \parencites[4]{5678011}[2]{4807303}.}\pdfcomment[icon=Note]{The GDPR consent requirements and the connection to consent-dialog design are documented in 5678011 (legal requirements for freely given/specific/informed/unambiguous/revocable consent, see 5678011, p.4) and the taxonomy/design guidance in 4807303; the sentence is a reasonable synthesis of those sources.} \hldarkgreen{Evidence indicates that dark\-pattern consent mechanisms reduce privacy\-preserving choices by up to $23$ percentage points, and that cookie banner design significantly determines consent withdrawal rates, demonstrating that apparent freedom of choice can be overridden by interface configuration \parencites[2]{4807303}.}\pdfcomment[icon=Note]{4807303 explicitly reports that dark-pattern consent mechanisms reduced privacy-preserving choices by up to 23 percentage points and cites work showing cookie banner design affects consent withdrawal rates (4807303, p.2072).} \hllightgreen{A critical How\-Might\-We prompt for MindfulFeed is therefore: How might consent and privacy flows be structured so that opting out, withdrawing consent, or limiting tracking is at least as salient and effortless as opting in, eliminating asymmetric effort and routinised conditioning effects identified in experimental work on consent dialogs \parencites[2]{4807303}[14]{5678011}?}\pdfcomment[icon=Note]{The question is a logical design prompt grounded in experimental findings about consent dialog effects and routinised conditioning discussed in 5678011 (see discussion of effect survival and educative/bright patterns, 5678011, p.15), so this is a valid academic synthesis.} \hldarkgreen{Complementary legal scholarship frames tracking walls and reduced service designs as unlawful or legally actionable, yet notes that such patterns continue to appear because current consent systems place substantial decision burdens on users in complex, cognitively demanding contexts \parencites[4]{5678011}[15]{5678011}.}\pdfcomment[icon=Note]{5678011 discusses that tracking walls and reduced-service designs can be legally actionable and notes these patterns persist because current consent systems place substantial decision burdens on users (5678011, pp.14--15).} \hldarkgreen{This motivates further questions about reducing cognitive load: How might MindfulFeed minimise the informational and attentional burden imposed at consent points, while still conveying meaningful data consequences through approaches such as "data visceralization" that connect users' feelings with their data in tangible, emotionally appreciable ways \parencites[15]{5678011}?}\pdfcomment[icon=Note]{5678011 explicitly mentions the notion of reducing informational/attentional burden and cites 'data visceralization' as a technique to make data consequences emotionally appreciable (5678011, p.15).} \hldarkgreen{Governance analysis conceptualises digital well\-being as a problem of accountability, transparency, oversight, and well\-being outcomes, highlighting how AI systems in education and mental health can erode human agency through surveillance, behavioural nudging, and psychological conditioning when learning patterns and mental states are optimised as data streams.}\pdfcomment[icon=Note]{Source 7808872 conceptualizes digital well-being as a governance problem framed by accountability, transparency, oversight and well-being outcomes and describes how AI systems in education/mental health can erode agency via surveillance, nudging and conditioning (7808872, p.3).} \hldarkgreen{Strategic How\-Might\-We questions must therefore address the governance frontier directly: How might interface and AI\-driven features in MindfulFeed make accountability for algorithmic decisions visible to end users, clarify which actors bear responsibility for psychological outcomes, and provide oversight channels that permit feedback and correction when recommendation or notification behaviours contribute to cognitive strain or emotional distress \parencites[3]{7808872}[8]{7808872}?}\pdfcomment[icon=Note]{7808872 lays out accountability, transparency and oversight as core governance dimensions and discusses making algorithmic operations and responsibilities visible to users (7808872, pp.8, 3), supporting the proposed questions.} \hldarkgreen{Because attention, emotion, and cognition are described as contested frontiers analogous to "new oil", and because governance aims to function both as shield and compass for human dignity, a further question arises: How might MindfulFeed embed transparent explanations of algorithmic operations affecting feed ranking, prompts, and behavioural feedback so that users can understand and contest potential nudging or conditioning \parencites[3]{7808872}?}\pdfcomment[icon=Note]{7808872 uses the explicit formulation that 'If data is the new oil, then attention, emotion, and cognition are its most contested frontiers' and argues governance should act as shield and compass (7808872, p.3).} \hldarkgreen{Awareness research shows that while dark patterns are widely visible on digital platforms, most users fail to recognise them as intentional manipulation, often misclassifying them as mere bad design and demonstrating particularly low comprehension of data\-driven dark patterns such as cookie consent manipulation \parencites[2]{5127387}.}\pdfcomment[icon=Note]{Source 5127387 reports that dark patterns are widely visible yet most users fail to recognise them as intentional manipulation and specifically notes low comprehension of data-driven patterns like cookie consent manipulation (5127387, p.2).} \hllightgreen{This calls for How\-Might\-We questions that combine design and education: How might MindfulFeed surface manipulative design risks in a way that increases user literacy about dark patterns without reproducing unreflective defaults or perceived lack of control associated with some "bright patterns" and pro\-privacy nudges \parencites[15]{5678011}[2]{5127387}?}\pdfcomment[icon=Note]{The call to combine design and education aligns with discussion in 5678011 about combining 'bright'/'educative' nudges and warnings and cautions about unreflective defaults (5678011, p.15), and 5127387 documents the need to increase user literacy (5127387, p.2); this is a reasonable synthesis.} \hllightgreen{Since end\-user centric consent regimes place substantial and possibly unwarranted pressure on individuals to make fully informed choices under uncertainty, another strategic question is: How might responsibility for ethical data practices be redistributed within MindfulFeed's ecosystem, for example by embedding default safeguards and constraint mechanisms at the system level, so that user choices operate within a protective envelope instead of being the sole barrier to exploitation \parencites[4]{5678011}[3]{7808872}?}\pdfcomment[icon=Note]{5678011 argues the GDPR places substantial pressure on end users to provide consent and questions whether burden should rest solely on users (5678011, p.15), and governance sources (7808872) discuss redistributing responsibility, so the recommendation is a supported inference.} \hldarkgreen{Finally, comprehensive taxonomies and automated detection pipelines such as DarkScan demonstrate that dark\-pattern exposure measurably degrades decision quality, platform trust, and psychological wellbeing, and they provide evidence\-based design and policy recommendations rooted in regulatory frameworks \parencites[2]{4807303}.}\pdfcomment[icon=Note]{4807303 describes both a comprehensive taxonomy and an automated detection pipeline (DarkScan) and reports measurable impacts of dark-pattern exposure on decision quality, trust and psychological wellbeing and offers design/policy recommendations (4807303, p.2072).} \hllightgreen{Strategic How\-Might\-We questions must integrate these diagnostics into the development record, asking: How might MindfulFeed use taxonomic clarity and automated checks as internal design constraints to ensure that no pattern category from established lists is instantiated in any prototype flow, and how might behavioural and trust metrics from usability testing be monitored as primary success indicators rather than engagement or conversion \parencites[2]{4807303}[2]{5127387}?}\pdfcomment[icon=Note]{4807303 presents taxonomy and automated detection work and derives recommendations for integrating these tools into design processes (4807303, p.2072); synthesising this into product-development constraints and monitoring behavioural/trust metrics is a reasonable academic inference.} \hllightgreen{Framing How\-Might\-We prompts in this empirically grounded way keeps subsequent ideation tethered to demonstrated harms, governance dimensions, and legal thresholds, providing a structured bridge from burnout narratives and dark\-pattern evidence to actionable, autonomy\-supportive design requirements.}\pdfcomment[icon=Note]{This is a synthesis claim about process and benefits of empirically grounded prompts; the components (linking harms to governance and legal thresholds) are supported across sources (e.g., 4807303 on harms/taxonomy and 7808872 on governance), so LIGHTGREEN is appropriate.} \hldarkgreen{Empirical work on platform amplification and early detection suggests useful operational details for monitoring and intervention: algorithms can create distinct promotion dynamics across platforms and early-warning models may detect escalation windows of a few hours, which appears to matter for when automated checks should fire and when human review is required \parencites[11]{8597215}.}\pdfcomment[icon=Note]{The platform-amplification and early-detection findings (distinct promotion dynamics and early-warning windows of hours) are reported in 8597215 with quantified platform differences and early-detection windows (e.g., 4--14 hours) (8597215, pp.4, 11).} \hldarkgreen{That same research reports that culturally informed interventions markedly reduce escalation compared with generic fact-checking, implying that automated detection should be paired with response protocols that are sensitive to cultural and linguistic context rather than relying solely on blunt suppression \parencites[11]{8597215}.}\pdfcomment[icon=Note]{8597215 explicitly reports culturally informed interventions reduced escalation markedly (67\% success) versus generic fact-checking (23\%), supporting the claim that detection should be paired with culturally sensitive response protocols (8597215, pp.9--10).}

\section{Ideate Phase: Concept Generation for Pro\-Social Social Media}

\subsection{Ideation Methods and Creative Exploration}
\label{sec:ideation_methods_creative_exploration} \hllightgreen{Ideation for MindfulFeed must operationalise ethical, governance, and dark pattern evidence into concrete concept generation procedures that can later be prototyped and evaluated.}\pdfcomment[icon=Note]{The recommendation to operationalise ethical, governance, and dark pattern evidence into concept generation is a reasonable synthesis of the governance-to-design framing in 7808872 (p.8) and dark-pattern design literature (4807303, p.2), though it is a prescriptive inference rather than a direct quote.} \hldarkgreen{Qualitative governance work on digital well\-being demonstrates the usefulness of multi\-layered analysis and thematic content coding across accountability, transparency, oversight, and well\-being outcomes.}\pdfcomment[icon=Note]{Source 7808872 (p.8) explicitly outlines a multi-layered qualitative approach and thematic content analysis using the four dimensions accountability, transparency, oversight, and well-being outcomes.} \hllightgreen{These dimensions can be treated as ideation lenses: for each candidate feature, structured workshops should deliberately explore how accountability is made visible, which transparency mechanisms are embedded, what oversight channels are feasible, and how emotional stability, cognitive autonomy, and social connectedness are supported in practice.}\pdfcomment[icon=Note]{Treating the four governance dimensions as 'ideation lenses' is a reasonable academic synthesis drawing on the four thematic dimensions in 7808872 (p.8) applied to design practice in the dark-pattern and governance literature (4807303, p.2).} \hldarkgreen{Ideation sessions can be organised around these four lenses, using prompts derived from coded policy and case evidence to systematically challenge feature concepts that might inadvertently reproduce manipulative or opaque interaction patterns \parencites[7]{7808872}[8]{7808872}.}\pdfcomment[icon=Note]{7808872 (p.8) describes using coded policy and case evidence with thematic prompts; the specific claim about organising sessions around those lenses and challenging manipulative patterns is directly supported by that methodological description.} \hllightgreen{Dark pattern taxonomies and detection research provide further scaffolding for creative exploration.}\pdfcomment[icon=Note]{Multiple sources on dark pattern taxonomies and detection (4807303, p.2; 4807303, p.1; 4692551, p.2) support the claim that these research strands can scaffold creative exploration, a reasonable synthesis.} \hldarkgreen{A five\-tier taxonomy built on mechanistic clarity, non\-overlap, and exhaustiveness identifies 27 specific patterns organised into categories that exploit cognitive biases such as anchoring, loss aversion, default effects, and least effort principles.}\pdfcomment[icon=Note]{Source 4807303 (p.2 and p.3) explicitly presents a five-tier taxonomy organised into top-level categories and states it organises 27 specific patterns, and 4807303/4807303 link these patterns to cognitive-bias mechanisms (anchoring, loss aversion, defaults, least effort).} \hllightgreen{Ideation activities should therefore incorporate pattern checklists and cognitive bias maps as constraint tools: brainstorming boards can list archetypal dark patterns (for example roach motel, privacy zuckering, misdirection, disguised ads) and associated bias mechanisms, with design teams required to demonstrate explicitly how emerging concepts avoid each mechanism or invert it toward autonomy support.}\pdfcomment[icon=Note]{Listing archetypal dark patterns (roach motel, privacy zuckering, misdirection, disguised ads) and recommending checklists is supported by foundational taxonomies and catalogs in 4807303 (p.1--2) and Brignull references summarized in 4807303 (p.1).} \hldarkgreen{This turns taxonomic knowledge into practical creative boundaries that prevent generation of concepts reliant on asymmetric opt\-in versus opt\-out flows, countdown timers for engagement, or preselected consent options that undermine privacy\-preserving choices \parencites[2]{4807303}.}\pdfcomment[icon=Note]{The examples of asymmetric opt-in/opt-out flows, countdown timers, and preselected consent options as problematic mechanisms are directly discussed as manipulative mechanisms in 4807303 (p.2) and 4807303 (p.2) which link these to reduced privacy-preserving choices.} \hldarkgreen{Expert\-led cognitive walkthrough methodology on social platforms such as YouTube, LinkedIn, WhatsApp, and Telegram offers a template for ideation grounded in interaction scenarios.}\pdfcomment[icon=Note]{Source 4692551 (p.3) documents expert-led cognitive walkthroughs across YouTube, LinkedIn, WhatsApp, and Telegram and frames this methodology as a template for analysing interaction scenarios.} \hldarkgreen{UI/UX specialists in that study executed twelve comprehensive tasks across content creation, engagement, account creation, and termination to identify combinations of dark patterns and analyse their impact on user experience.}\pdfcomment[icon=Note]{4692551 (p.3) states that UI/UX specialists completed twelve comprehensive tasks across content creation, engagement, account creation, and termination to identify dark-pattern combinations, matching this sentence directly.} \hllightgreen{MindfulFeed's ideation workshops can repurpose these task structures as generative prompts: for each key journey (onboarding, feed browsing, notification management, account exit), teams sketch low\-fidelity flow diagrams and then annotate them with potential dark pattern risks identified in the walkthrough research.}\pdfcomment[icon=Note]{Repurposing walkthrough task structures as generative prompts for ideation (onboarding, browsing, notifications, exit) is a reasonable application inference from the task-based walkthrough methodology described in 4692551 (p.3).} \hllightgreen{Subsequent creative refinements focus on replacing risky segments with alternative interaction strategies that foreground clarity, symmetric effort, and cognitive relief, using the combinations documented in the walkthroughs as explicit anti\-patterns \parencites[3]{4692551}.}\pdfcomment[icon=Note]{Replacing risky segments with alternative interaction strategies and using documented combinations as anti-patterns is a plausible design synthesis grounded in the walkthrough analyses reported in 4692551 (p.3) and dark-pattern literature (4807303, p.2).} \hldarkgreen{Systematic review methodology in AI\-based burnout interventions adds a further layer of rigor to concept generation.}\pdfcomment[icon=Note]{Source 8760542 (p.19) describes a systematic review methodology applied to AI-based burnout/stress interventions, supporting the claim that such methodology adds rigor to concept generation.} \hldarkgreen{The PRISMA\-guided review employed staged processes of identification, screening, eligibility, and inclusion, supported by structured coding frameworks, thematic clustering, and quality appraisal via the Mixed Methods Appraisal Tool.}\pdfcomment[icon=Note]{8760542 (p.19) explicitly states the PRISMA-guided review stages of identification, screening, eligibility, and inclusion, plus structured coding and quality appraisal via MMAT.} \hllightgreen{Translating these procedures into ideation practice, candidate features for MindfulFeed can pass through analogous stages: initial idea capture, screening against governance and dark pattern criteria, eligibility assessment based on anticipated impact on stress or burnout, and inclusion into prototype backlogs only when preliminary evidence and ethical transparency requirements are satisfied.}\pdfcomment[icon=Note]{Translating PRISMA-like staged processes into ideation (capture, screening, eligibility, inclusion) is a reasonable methodological synthesis grounded in 8760542 (p.19) which describes those stages and coding frameworks.} \hllightgreen{Feature matrices can record variables such as anticipated computational method, data sources, and targeted outcomes, mirroring the coding of intervention type and measured effects in the review and thereby maintaining analytical integrity in concept selection \parencites[19]{8760542}.}\pdfcomment[icon=Note]{8760542 (p.19) documents coding of intervention type, methods, data sources, and outcomes; suggesting feature matrices mirror this coding in ideation is a logical extension and synthesis of that review practice.} \hldarkgreen{Research on digital aggression and algorithmic amplification points to multifaceted strategies, including algorithmic tweaks that prioritise constructive discourse over divisive content and digital literacy education to help users recognise manipulative tactics designed to provoke outrage.}\pdfcomment[icon=Note]{Source 2937051 (p.5) explicitly recommends multifaceted strategies to address digital aggression, including algorithmic tweaks prioritising constructive discourse and digital literacy education to recognise manipulative tactics.} \hllightgreen{These insights can be incorporated into ideation by defining concept clusters around constructive feed logic and literacy\-enhancing interface elements.}\pdfcomment[icon=Note]{Incorporating insights about constructive feed logic and literacy elements into ideation is a reasonable design-oriented inference from 2937051 (p.5), which suggests these strategies.} \hllightgreen{For example, ideation sprints can focus on alternative ranking algorithms that deprioritise outrage cues, combined with explanatory microcopy that trains users to identify emotionally manipulative patterns in their broader digital environment.}\pdfcomment[icon=Note]{The concrete example of sprints on alternative ranking algorithms and explanatory microcopy synthesises recommendations in 2937051 (p.5) about deprioritising outrage cues and promoting digital literacy, a reasonable inference.} \hllightgreen{Each concept in these clusters is then evaluated against hypotheses about anonymity, algorithmic influence, and moderation efficacy, ensuring that creative exploration remains aligned with documented relationships between social media algorithms and online aggression \parencites[5]{2937051}.}\pdfcomment[icon=Note]{Evaluating concepts against hypotheses about anonymity, algorithmic influence, and moderation efficacy aligns with the documented relationships and research gaps discussed in 2937051 (p.5), making this a supported synthesis.} \hldarkgreen{Finally, methodological principles of credibility, transferability, and dependability in qualitative governance research demand that ideation artefacts be systematically documented.}\pdfcomment[icon=Note]{7808872 (p.8) explicitly prioritizes credibility, transferability, and dependability for qualitative governance research and thus supports the claim that these principles demand systematic documentation.} \hldarkgreen{Session outputs such as annotated user journeys, pattern\-risk maps, and governance lens canvases should be archived and coded, enabling subsequent phases to trace how specific MindfulFeed features emerged from evidence\-constrained creative exploration rather than ad hoc brainstorming.}\pdfcomment[icon=Note]{7808872 (p.8) describes using NVivo and systematic coding to enable traceability, directly supporting the recommendation to archive and code session outputs for traceability in later phases.} \hldarkgreen{This documentation practice directly mirrors the use of NVivo for thematic organisation and the emphasis on detailed contextual description and transparent coding in governance studies, thereby embedding methodological integrity into the ideation phase itself \parencites[7]{7808872}[8]{7808872}[19]{8760542}.}\pdfcomment[icon=Note]{7808872 (p.8) documents NVivo use for thematic organisation and transparent coding, and 8760542 (p.19) details PRISMA/MMAT practices, together directly supporting the claim that this documentation mirrors those governance study methods.}

\subsection{Conceptual Design Space for MindfulFeed}
\label{sec:conceptual_design_space_mindfulfeed} \hllightgreen{Conceptual boundaries for MindfulFeed's design space are strongly shaped by empirical dark pattern research and governance analyses of digital well\-being.}\pdfcomment[icon=Note]{This is a reasonable synthesis: Source 4807303 (p.2--3) provides empirical dark-pattern research (taxonomy, DarkScan) and Source 7808872 (pp.2--4) frames digital well‑being as a governance problem, supporting the claim that both empirics and governance shape conceptual boundaries.} \hldarkgreen{A five\-tier taxonomy built on mechanistic clarity, non\-overlap, and exhaustiveness organises $27$ specific patterns into five top\-level categories, providing a structured negative design space that MindfulFeed must systematically avoid.}\pdfcomment[icon=Note]{Directly supported by Source 4807303 which states a five‑tier taxonomy built on mechanistic clarity, non‑overlap, exhaustiveness and organises 27 specific patterns into five top‑level categories (pp.2--3).} \hldarkgreen{Archetypal patterns such as trick questions, sneak\-into\-basket, roach motel, privacy zuckering, misdirection, and disguised ads exemplify manipulative mechanisms that exploit cognitive biases and undermine autonomous decision making.}\pdfcomment[icon=Note]{Source 4807303 cites Brignull's original catalogue and explicitly lists archetypal dark patterns including trick questions, sneak‑into‑basket, roach motel, privacy zuckering, misdirection, and disguised ads (p.2).} \hllightgreen{Any conceptual feature canvas for MindfulFeed must therefore treat these named patterns, and the broader categories they instantiate, as explicit exclusion constraints when ideating consent flows, onboarding, feed interaction, and account management.}\pdfcomment[icon=Note]{Supported by recommendations in Source 4807303 (pp.7--8) which advise treating dark patterns as prohibited practices and designing symmetric, non‑manipulative flows, so treating them as exclusion constraints is a valid synthesis.} \hllightgreen{Evidence on cognitive mechanisms further narrows the permissible design space.}\pdfcomment[icon=Note]{Source 4807303 (b, p.2--3) documents cognitive mechanisms exploited by dark patterns, which reasonably implies that such evidence narrows permissible design choices.} \hldarkgreen{Dark patterns leverage $ \text{System 1} $ processing, anchoring bias, loss aversion, default effects, and the principle of least effort through devices such as decoy pricing, countdown timers, preselected options, and asymmetric opt\-in versus opt\-out flows.}\pdfcomment[icon=Note]{Source 4807303 (b, p.2--3) explicitly states that dark patterns leverage System 1 processing and specific mechanisms (anchoring, loss aversion, default effects, principle of least effort) and gives examples like decoy pricing, countdown timers, preselected options, and asymmetric opt‑in/opt‑out.} \hllightgreen{Conceptually, MindfulFeed's interaction models must invert these mechanisms: any use of defaults, scarcity messages, or effort asymmetries is treated as high risk and either prohibited or reconfigured to support autonomy, for instance by equalising the number of steps and time required for engagement versus disengagement actions.}\pdfcomment[icon=Note]{Source 4807303 (pp.7--8) recommends symmetric design and equalising opt‑in/opt‑out complexity; treating defaults, scarcity messages, and effort asymmetries as high risk is a reasonable policy synthesis from those recommendations.} \hldarkgreen{This inversion is not speculative; experimental work has shown that dark\-pattern consent mechanisms reduce privacy\-preserving choices by up to $23$ percentage points, and cookie banner designs significantly determine consent withdrawal rates, confirming that interface architecture can override nominal freedom of choice \parencites[2]{4807303}.}\pdfcomment[icon=Note]{Source 4807303 (b, p.2--3) reports experimental findings that dark‑pattern consent mechanisms reduced privacy‑preserving choices by up to 23 percentage points and that cookie‑banner designs significantly affect consent withdrawal, directly supporting the claim.} \hldarkgreen{Within MindfulFeed's design space, consent and privacy features are therefore constrained to plain language, symmetric effort, and the absence of countdowns or guilt framing.}\pdfcomment[icon=Note]{Source 4807303 recommendations (pp.7--8) call for plain language, symmetric effort in user journeys, and avoidance of guilt/countdown framing, directly supporting these constraints on consent and privacy features.} \hllightgreen{Governance perspectives extend these constraints beyond microinteractions to macro design objectives.}\pdfcomment[icon=Note]{Source 7808872 (pp.2--4) frames digital well‑being as a governance problem and argues governance perspectives should inform higher‑level design objectives, supporting this extension from micro to macro.} \hldarkgreen{Digital well\-being is framed as a governance problem centred on accountability, transparency, oversight, and well\-being outcomes, with AI systems in education and mental health currently operating in regulatory grey zones where surveillance, behavioural nudging, and psychological conditioning erode human agency \parencites[2]{7808872}[3]{7808872}.}\pdfcomment[icon=Note]{Source 7808872 (pp.2--3, and analytical framework p.8) explicitly frames digital well‑being as a governance problem centred on accountability, transparency, oversight, and well‑being outcomes and notes AI in education/mental health operating in regulatory grey zones with surveillance, nudging, and conditioning (pp.2--3, 8).} \hllightgreen{Translating this into a conceptual design space, MindfulFeed's feature set must only include AI\-mediated components that can be made visibly accountable, explainable, and subject to human feedback or correction, and that are explicitly oriented toward emotional stability, cognitive autonomy, and social connection rather than engagement maximisation.}\pdfcomment[icon=Note]{Source 7808872 (pp.3--4) and Source 4807303 (pp.7--8) emphasize accountability, transparency, human oversight and orientation toward emotional/cognitive well‑being, supporting this normative design constraint as a synthesis.} \hllightgreen{Algorithms that simulate empathy or predict affective states without clear oversight channels fall outside the acceptable design space, given concerns about subtle erosion of the capacity to think, choose, and feel autonomously in digital environments \parencites[3]{7808872}.}\pdfcomment[icon=Note]{Source 7808872 (pp.2--4) discusses simulated empathy and the erosion of human agency; combined with governance emphasis on oversight, this supports excluding algorithms that simulate affect without clear oversight.} \hldarkgreen{The GDW Framework's five pillars of ethical accountability, algorithmic transparency, human oversight, psychological safety, and inclusive policymaking act as additional structuring dimensions for concept generation.}\pdfcomment[icon=Note]{Source 7808872 (pp.3--4) explicitly presents the GDW Framework and its five pillars: ethical accountability, algorithmic transparency, human oversight, psychological safety, and inclusive policymaking.} \hllightgreen{Features that cannot be mapped to at least accountability, transparency, and psychological safety should be excluded at the ideation stage, since governance analysis argues that systems must be accountable to the moral conditions of human life, not only to regulators or efficiency metrics \parencites[4]{7808872}.}\pdfcomment[icon=Note]{Source 7808872 (pp.3--4) argues systems should be accountable to moral conditions of human life beyond mere regulatory compliance; coupled with Source 4807303 recommendations (pp.7--8), excluding features that cannot map to key pillars is a supported synthesis.} \hllightgreen{This means, for example, that any persuasive feature envisaged for MindfulFeed must allow users to understand why a particular prompt appears, who defined the behavioural objective, and how negative psychological impacts will be monitored and mitigated.}\pdfcomment[icon=Note]{Transparency and accountability recommendations in Sources 7808872 (pp.3--4) and 4807303 (pp.7--8) support the concrete example that persuasive features must disclose rationale, objectives, and monitoring/mitigation plans.} \hldarkgreen{Automated detection work with DarkScan also informs the design space by defining measurable thresholds of harm.}\pdfcomment[icon=Note]{Source 4807303 (pp.2, 7--8) describes DarkScan and its role in automated detection and severity scoring (DPSS), indicating that automated detection informs measurable harm thresholds.} \hldarkgreen{A comprehensive empirical investigation combining a five\-category taxonomy, automated DOM and visual analysis, a severity scoring model (DPSS), and a controlled user experiment demonstrates that dark patterns are pervasive and produce measurable behavioural and psychological harm, degrading decision quality, platform trust, and psychological wellbeing \parencites[2]{4807303}.}\pdfcomment[icon=Note]{Source 4807303 (conclusion and pp.2--8) reports a comprehensive empirical investigation combining a five‑category taxonomy, DOM/visual analysis, DPSS, and a controlled experiment, finding pervasive dark patterns with measurable behavioural and psychological harms to decision quality, trust, and wellbeing.} \hldarkgreen{Policy and design recommendations derived from these findings include symmetric design principles, plain language copy, organisational policies that prohibit dark patterns, and success metrics that incorporate user trust and long\-term retention rather than short\-term conversion \parencites[8]{4807303}.}\pdfcomment[icon=Note]{Source 4807303 (p.8, Recommendations) explicitly lists symmetric design principles, plain language, organisational policies prohibiting dark patterns, and aligning success metrics to user trust and long‑term retention.} \hllightgreen{Conceptually, MindfulFeed's design space is therefore bounded by commitments to symmetric pathways for opt\-out and cancellation, non\-manipulative language, and evaluation criteria focused on trust and wellbeing rather than conversion funnels.}\pdfcomment[icon=Note]{This is a reasonable synthesis of Source 4807303 recommendations (pp.7--8) and governance emphasis in Source 7808872 (pp.3--4) that bind the design space to symmetric opt‑out/cancellation, non‑manipulative language, and trust/wellbeing metrics.} \hllightgreen{Finally, trustworthiness criteria from governance research impose documentation expectations on the conceptual design space itself.}\pdfcomment[icon=Note]{Source 7808872 (pp.3, 8) and Source 4807303 (p.8) discuss trustworthiness and transparency expectations, which reasonably imply documentation obligations for the conceptual design space.} \hldarkgreen{Thematic content analysis emphasises credibility through data triangulation, transferability through detailed contextual description, and dependability through transparent coding and iterative peer debriefing \parencites[8]{7808872}.}\pdfcomment[icon=Note]{Source 7808872 (p.8) explicitly details thematic content analysis trustworthiness practices: credibility via data triangulation, transferability via contextual description, and dependability via transparent coding and peer debriefing.} \hllightgreen{Applied to MindfulFeed, this requires that every proposed concept be traceable to specific dark pattern evidence or governance principles and that design rationales be recorded as artefacts subject to later audit.}\pdfcomment[icon=Note]{Source 7808872 (pp.3, 8) and Source 4807303 (p.8) endorse traceability, documentation of user journeys and transparency reports, supporting the requirement that concepts be traceable and recorded for audit as a synthesis.} \hllightgreen{Concepts lacking such evidence linkage or documentation fall outside the robust, evidence\-based design space appropriate for an ethical, pro\-social social app.}\pdfcomment[icon=Note]{Both Source 4807303 (pp.7--8) and Source 7808872 (pp.3--4, 8) argue for evidence‑based design and documentation; excluding concepts lacking evidence linkage is a defensible inference from these sources.}

\section{Prototype Phase: Progressive Design of MindfulFeed}

\subsection{Prototype v1: Low\-Fidelity Batch Feed Concept}
\label{sec:prototype_v1_batch_feed} \hllightgreen{Prototype v1 is scoped as a low\-fidelity batch feed concept that operationalizes evidence on dark patterns, consent requirements, and personalization fatigue into a constrained interaction model.}\pdfcomment[icon=Note]{The claim synthesizes design choices grounded in dark-pattern scholarship and personalization-fatigue mitigation; sources discuss prototypes and translating evidence into constrained interaction models (see 0270098 p.6 and 4807303 p.2).} \hldarkgreen{Dark pattern taxonomies identify mechanisms such as sneaking, obstruction, misdirection, forced action, confirm shaming, roach motel, and drip pricing, all of which deliberately exploit cognitive biases like default bias, scarcity bias, framing, and loss aversion to steer users into unintended actions \parencites[2]{5127387}[2]{4807303}.}\pdfcomment[icon=Note]{Taxonomies and associated cognitive biases are explicitly listed in the literature (5127387 p.2 lists types including sneaking, obstruction, misdirection, confirm shaming, roach motel, drip pricing and common biases; 4807303 p.2 links patterns to biases).} \hldarkgreen{Experimental work shows that dark\-pattern consent mechanisms can reduce privacy\-preserving choices by up to $23$ percentage points, and that cookie banner design strongly determines consent withdrawal rates, demonstrating that interface architecture can override nominal freedom of choice \parencites[2]{4807303}.}\pdfcomment[icon=Note]{The empirical magnitudes cited are reported in the provided review: 4807303 p.2 reports a reduction of up to 23 percentage points and notes cookie-banner design strongly determines consent withdrawal rates.} \hllightgreen{To avoid reproducing these harms, the batch feed prototype is explicitly framed as a non\-scrolling, session\-bounded layout that presents a finite set of items without countdowns, auto\-advance, or preselected engagement defaults.}\pdfcomment[icon=Note]{Avoiding urgency cues, preselected defaults, and continuous feeds is consistent with recommendations to reduce dark-pattern effects and personalization fatigue (0270098 p.6; 5127387 p.2 describes urgency/scarcity as manipulative).} \hllightgreen{Low\-fidelity feed sketches are structured around cards arranged in discrete batches, with explicit "next batch" controls instead of infinite scroll.}\pdfcomment[icon=Note]{UI-based implementations and alternatives to continuous engagement (e.g., discrete controls instead of infinite scroll) are advocated in the UI literature (0270098 p.6 describes concrete UI examples and session-level heuristics).} \hldarkgreen{This configuration responds to analyses of dark patterns that emphasize how asymmetric opt\-in versus opt\-out flows and least\-effort principles can trap users in roach motel configurations where entry is easy but exit is difficult \parencites[2]{5127387}[2]{4807303}.}\pdfcomment[icon=Note]{4807303 p.2 explicitly discusses asymmetric opt‑in/opt‑out flows, the principle of least effort, and roach‑motel style traps as mechanisms exploited by dark patterns.} \hllightgreen{Each batch page in the prototype therefore contains a visible "end of batch" marker and a neutral "pause" affordance rather than engagement prompts framed through urgency or scarcity.}\pdfcomment[icon=Note]{Design elements such as visible end markers and neutral pause affordances map to fatigue-aware nudges and avoidance of urgency cues recommended in the framework (0270098 p.6 Table; 0270098 p.6 discusses breathing-space options).} \hllightgreen{Interface copy is deliberately stripped of emotional pressure, avoiding confirm shaming language that would induce guilt for disengaging \parencites[2]{5127387}.}\pdfcomment[icon=Note]{Avoiding confirm‑shaming and emotionally coercive copy is supported by dark-pattern taxonomies and guidance that list confirm shaming as a harmful pattern (5127387 p.2) and legal/ethical critiques (5678011 p.4).} \hllightgreen{In addition, decision points such as saving items or adjusting preferences are designed with symmetric effort: accepting and declining have equal visual prominence and identical click depth, aligning with legal expectations that consent be freely given, specific, informed, unambiguous, and easy to withdraw \parencites[4]{5678011}[2]{4807303}.}\pdfcomment[icon=Note]{Combining design symmetry (equal prominence and click depth) with legal consent principles is a reasonable synthesis: GDPR consent requirements are cited (5678011 p.4) and the design rationale is consistent with recommended best practices (0270098 p.6).} \hllightgreen{Consent and personalization controls embedded into the batch feed follow legal and governance guidance on consent banners and digital well\-being.}\pdfcomment[icon=Note]{The statement draws on legal consent guidance and governance perspectives on digital well‑being found in the sources (5678011 p.4 on consent banners; 7808872 p.3 on governance/digital well‑being).} \hldarkgreen{GDPR and the ePrivacy Directive require that consent be presented in a manner clearly distinguishable from other matters, with meaningful choice evidenced by the ability to withdraw consent, and they caution that making access conditional on consent can invalidate reliance on consent \parencites[4]{5678011}.}\pdfcomment[icon=Note]{5678011 p.4 directly summarizes GDPR/ePrivacy requirements: consent must be clearly distinguishable, withdrawable, and making access conditional on consent can invalidate reliance on consent.} \hllightgreen{In the low\-fidelity concept, batch\-level personalization toggles are surfaced at the top of the feed rather than buried in secondary menus.}\pdfcomment[icon=Note]{Making personalization controls accessible rather than buried is explicitly recommended in the tiered personalization guidance (0270098 p.6 states settings should not be buried under legalese or multi‑tab menus).} \hldarkgreen{Drawing on tiered personalization settings work, prototype controls allow selection of "Minimal," "Balanced," or "Maximized" personalization, explicitly mapped to tracking granularity, content diversity, and data retention, and presented as user sovereignty features rather than compliance formalities.}\pdfcomment[icon=Note]{0270098 p.6 explicitly prescribes 'Minimal,' 'Balanced,' and 'Maximized' tiers mapped to tracking granularity, visual density, and retention, and frames them as user‑sovereignty features.} \hldarkgreen{These settings are accessible without navigating multi\-tab menus or dense legalese, reflecting evidence that accessible, user\-selectable personalization tiers can reduce fatigue and strengthen perceived autonomy.}\pdfcomment[icon=Note]{0270098 p.6 states these tiers should be accessible (not buried) and argues they reduce fatigue and strengthen perceived autonomy, supporting the claim directly.} \hldarkgreen{The batch feed prototype also integrates mitigation strategies against personalization fatigue documented in UI frameworks.}\pdfcomment[icon=Note]{0270098 p.6 documents mitigation strategies against personalization fatigue in a UI framework and presents UI‑based implementations that prototypes can integrate.} \hldarkgreen{Adaptive exposure control, diversity\-enriched recommendations, exposure freshness windows, transparent recommendation logic, fatigue\-aware UI nudges, and tiered personalization are described as six strategies that shift personalization from extraction to empathy.}\pdfcomment[icon=Note]{0270098 p.6 Table 3 enumerates the six strategies listed (adaptive exposure control, diversity‑enriched recommendations, freshness windows, transparent logic, fatigue‑aware nudges, tiered personalization) and frames the shift from extraction to empathy.} \hldarkgreen{At low fidelity, this translates into simple placeholders: each item card includes a short "Why this item?" tooltip zone for later implementation of explanation logic; batch headers reserve space for "freshness indicators" that can mark when content has been recently updated or deliberately rotated; and a "breathing space" button is positioned alongside navigation controls to prototype fatigue\-aware nudges that offer quick exits or reflective pauses instead of additional content.}\pdfcomment[icon=Note]{0270098 p.6 describes explanation tooltips ('Why this item?'), freshness indicators, and breathing‑space options as concrete UI placeholders, matching the low‑fidelity placeholders described.} \hllightgreen{These layout decisions are annotated in the design record against the corresponding framework dimensions (algorithmic transparency, emotional salience, user sovereignty), ensuring traceability from empirical mitigation strategies to concrete interface structures \parencites[6]{0270098}.}\pdfcomment[icon=Note]{0270098 maps layout decisions to framework dimensions (algorithmic transparency, emotional salience, user sovereignty), so annotating design records for traceability is a reasonable synthesis of that guidance (0270098 p.6).} \hldarkgreen{Governance perspectives on digital well\-being define attention, emotion, and cognition as contested frontiers and argue that AI systems in education and mental health currently operate in regulatory grey zones where surveillance, behavioural nudging, and psychological conditioning can erode human agency \parencites[3]{7808872}.}\pdfcomment[icon=Note]{7808872 p.3 explicitly frames attention, emotion, and cognition as contested frontiers and states that AI systems in education and mental health operate in regulatory grey zones where surveillance and nudging can erode agency.} \hllightgreen{Although this prototype remains low\-fidelity and non\-functional, batch feed flows are documented to prevent implicit nudging through opaque ranking or simulated empathy.}\pdfcomment[icon=Note]{The claim that a documented low‑fidelity prototype seeks to prevent implicit nudging via opaque ranking or simulated empathy aligns with governance and UI guidance emphasizing explainability and avoidance of manipulative ranking (7808872 p.3; 0270098 p.6).} \hllightgreen{Wireframe annotations specify that recommendation ordering must later be explainable and adjustable by users, and that no batch navigation element may depend on hidden data processing.}\pdfcomment[icon=Note]{Requiring recommendation ordering to be explainable and preventing navigation elements from relying on hidden processing follows recommendations for explainability and user control in the literature (0270098 p.6; 5375614 p.6).} \hllightgreen{This aligns the concept with governance pillars of accountability, transparency, oversight, and psychological safety, positioning the batch feed as a structural countermeasure against quiet architectures of dependency \parencites[3]{7808872}[8]{7808872}.}\pdfcomment[icon=Note]{Connecting the concept to governance pillars (accountability, transparency, oversight, psychological safety) and positioning it as countermeasure against dependency echoes themes in the governance literature (7808872 p.3).} \hllightgreen{Finally, the prototype documentation includes an AI usage note for future iterations, informed by educational AI governance.}\pdfcomment[icon=Note]{Including an AI usage note informed by educational AI governance is consistent with calls for governance guidance in education contexts (5375614 p.6; 7808872 p.3).} \hldarkgreen{Personalization and learning analytics in education are said to require explanations understandable to users and evaluation beyond prediction accuracy, including impacts on different groups and long\-term consequences of labeling \parencites[6]{5375614}.}\pdfcomment[icon=Note]{5375614 p.6 explicitly argues that personalization and learning‑analytics need explanations understandable to users and evaluations beyond accuracy, including group impacts and long‑term labeling consequences.} \hllightgreen{In the batch feed record, any planned AI component (for example, session\-level heuristics for adaptive exposure) is listed in an appendix with intended function, data inputs, and oversight mechanisms, thereby conforming to demands for transparency, human oversight, and differentiated use that supports learning rather than replacing intellectual processes \parencites[6]{5375614}[3]{7808872}.}\pdfcomment[icon=Note]{Listing planned AI components with function, inputs, and oversight to satisfy transparency and human‑oversight demands aligns with literature recommending such documentation and oversight in educational AI (5375614 p.6; 7808872 p.3).}

\subsection{Prototype v2: Mid\-Fidelity Algorithm Control Dashboard}
\label{sec:prototype_v2_alg_control_dashboard} \hllightgreen{Mid\-fidelity prototyping of an algorithm control dashboard is positioned as a response to documented risks around opaque personalization, dark\-pattern consent design, and erosion of autonomy in AI\-mediated environments.}\pdfcomment[icon=Note]{Source 7808872 (p.2) documents risks from opaque personalization and erosion of agency in AI environments and Source 0270098 (p.6) offers UI prototyping responses (tiering, fatigue mitigation), so this is a reasonable synthesis of those sources.} \hldarkgreen{Governance analysis identifies accountability, transparency, oversight, and well\-being outcomes as core dimensions for evaluating AI systems that affect cognition and emotion, noting that current platforms often operate in regulatory grey zones where surveillance, behavioural nudging, and psychological conditioning undermine agency \parencites[2]{7808872}[8]{7808872}.}\pdfcomment[icon=Note]{Source 7808872 (p.8 and p.2) explicitly frames governance dimensions as accountability, transparency, oversight, and well-being outcomes and notes platforms operating in regulatory grey zones with surveillance, nudging, and conditioning.} \hllightgreen{Within this evidence base, the dashboard concept is specified as an interface that exposes algorithmic operations, allows users to configure personalization intensity, and surfaces consent status in a way that facilitates oversight rather than control.}\pdfcomment[icon=Note]{Source 0270098 (p.6) specifies interfaces that expose recommendation logic, allow tiered personalization intensity and emphasize accessible consent; this sentence synthesizes those design goals with governance themes from 7808872 (p.2).} \hldarkgreen{Design of mid\-fidelity controls draws directly on UI\-based personalization fatigue mitigation strategies.}\pdfcomment[icon=Note]{Source 0270098 (p.6, Table 3 and section 5) explicitly grounds mid‑fidelity control design in UI‑based personalization fatigue mitigation strategies (fatigue‑aware UI nudges, explanation tooltips, etc.).} \hldarkgreen{Tiered personalization settings are recommended as a way to respect differing comfort levels with tracking and tailoring, by offering "Minimal", "Balanced", and "Maximized" options that affect recommendation logic, visual density, tracking granularity, and data retention.}\pdfcomment[icon=Note]{Source 0270098 (p.6, '4.6 Tiered Personalization Settings') explicitly recommends 'Minimal', 'Balanced', and 'Maximized' tiers affecting recommendation logic, visual density, tracking granularity, and data retention.} \hldarkgreen{Crucially, these settings must be accessible and not buried under legalese or multi\-tab menus, and the act of choosing among tiers is framed as a gesture of respect rather than a compliance exercise.}\pdfcomment[icon=Note]{Source 0270098 (p.6) directly states these settings must be accessible (not buried under legalese) and frames choosing tiers as a gesture of respect rather than a compliance exercise.} \hldarkgreen{The same framework lists transparent recommendation logic and fatigue\-aware UI nudges as key strategies, implemented through explanation tooltips, preference dashboards, breathing space options, and engagement monitoring to reduce overwhelm and strengthen trust \parencites[6]{0270098}.}\pdfcomment[icon=Note]{Source 0270098 (p.6, Table 3) lists transparent recommendation logic and fatigue‑aware UI nudges and maps implementations to explanation tooltips, preference dashboards, breathing space options, and engagement monitoring.} \hllightgreen{In the mid\-fidelity dashboard, these strategies are translated into wireframed modules: a tier selector rendered as equal\-prominence buttons, a "Why this content?" explanation zone mapped to specific feed segments, and a fatigue panel that reserves space for future breathing\-space configurations.}\pdfcomment[icon=Note]{Source 0270098 (p.6) and 1912034 (p.7) discuss wireframe‑style UI elements such as tier selectors and 'Why this content?' explanations; the specific button prominence and reserved fatigue panel are reasonable design translations of those sources.} \hllightgreen{Dark pattern research provides constraints on how consent and personalization controls can be structured.}\pdfcomment[icon=Note]{Sources on dark patterns (4807303 p.2; 5353824 p.5) outline constraints and legal/regulatory considerations that inform how consent and personalization controls should be structured, so this is a supported general claim.} \hldarkgreen{Brignull's catalogue and subsequent taxonomies enumerate manipulative strategies such as trick questions, sneak\-into\-basket, roach motel, privacy zuckering, misdirection, and disguised ads, all grounded in cognitive biases including $ \text{System 1} $ reliance, anchoring, loss aversion, default effects, and least effort.}\pdfcomment[icon=Note]{Source 4807303 (p.2) documents Brignull's catalogue and lists patterns (trick questions, sneak‑into‑basket, roach motel, privacy zuckering, misdirection, disguised ads) and links them to cognitive biases including System 1 reliance, anchoring, loss aversion, default effects, and least effort.} \hldarkgreen{Randomised trials and consent experiments have shown that aggressive dark patterns can increase subscription uptake dramatically and that dark\-pattern consent mechanisms reduce privacy\-preserving choices by up to $23$ percentage points, while cookie banner design significantly determines consent withdrawal rates.}\pdfcomment[icon=Note]{Source 4807303 (p.2) cites randomized trials (e.g., Luguri \& Strahilevitz) showing dramatic increases in uptake and reports dark‑pattern consent mechanisms reducing privacy‑preserving choices by up to 23 percentage points and notes cookie banner design effects.} \hldarkgreen{These findings confirm that apparent freedom of choice can be architecturally determined by asymmetrical opt\-in versus opt\-out flows and preselected options \parencites[2]{4807303}.}\pdfcomment[icon=Note]{Source 4807303 (p.2) explicitly links apparent freedom of choice to architectural features like asymmetrical opt‑in/opt‑out flows and preselected defaults, supporting this claim.} \hllightgreen{Accordingly, the mid\-fidelity dashboard mockups explicitly avoid prechecked boxes, countdown timers, or buried "off" switches.}\pdfcomment[icon=Note]{Given the empirical findings on the harms of prechecked boxes, timers, and similar dark patterns (4807303 p.2), stating that mockups avoid these elements is a supported design decision inferred from the literature.} \hlorange{Each toggle for tracking, data retention, or personalization intensity is designed with symmetric visual weight and click depth for enabling and disabling, and batch diagrams of navigation flows document equal step counts for entering and exiting personalized modes.}\pdfcomment[icon=Note]{None of the provided sources describe symmetric visual weight, equal click depth for toggles, or batch diagrams documenting equal step counts for entering/exiting modes; this specific implementation detail has no direct support in the sources.} \hllightgreen{Ethical governance in AI\-enabled monitoring environments further informs dashboard content.}\pdfcomment[icon=Note]{Source 8760542 (p.18) and 7808872 (p.2) discuss how ethical governance for AI‑enabled monitoring informs design and content of interfaces used in occupational or wellbeing contexts, supporting this linkage.} \hldarkgreen{Occupational health analyses stress that privacy and informed consent are central when data streams include psychological states and behavioural patterns, and that transparent consent protocols must clearly inform users about what data are collected, how they are processed, and for what purposes.}\pdfcomment[icon=Note]{Source 8760542 (p.18) explicitly states that privacy and informed consent are central when data include psychological states and behavioural patterns and calls for transparent consent protocols describing collected data, processing, and purposes.} \hldarkgreen{Consent should be iterative and revocable, allowing opt\-out without penalty; data minimization should restrict retention to information necessary for legitimate objectives; and algorithmic transparency must make predictive functioning explainable and auditable.}\pdfcomment[icon=Note]{Source 8760542 (p.18) explicitly recommends iterative, revocable consent allowing opt‑out without penalty, data minimization for retention, and algorithmic transparency for explainability and auditability.} \hldarkgreen{Absent such transparency, AI risks being perceived as an intrusive management tool rather than a supportive wellness mechanism \parencites[18]{8760542}.}\pdfcomment[icon=Note]{Source 8760542 (p.18) states that absence of transparency risks eroding trust and leads AI to be perceived as intrusive management rather than supportive wellness, directly supporting this sentence.} \hllightgreen{Reflecting these requirements, the mid\-fidelity dashboard allocates explicit sections for "Data in use", "Purpose of analysis", and "Current consent state", with flow diagrams indicating how a user can revoke consent at any time without service degradation.}\pdfcomment[icon=Note]{Sources advocating transparent 'data in use' and consent state displays (0270098 p.6; 8760542 p.18) support allocating dashboard sections for data, purposes, and consent and suggest flow diagrams for revocation, making this a reasonable synthesis.} \hldarkgreen{Educational AI research elaborates how personalization and learning analytics should be governed, emphasizing that algorithms must not determine paths in a closed\-off way and that students need to understand recommendations, have choices, and be able to reject or adjust profiles.}\pdfcomment[icon=Note]{Source 5375614 (p.6) explicitly argues that personalization and learning analytics must not deterministically lock student paths and that students should understand, choose, reject, or adjust recommendations/profiles.} \hldarkgreen{Recent literature calls for human oversight, transparency, fairness, accountability, and explainability, with algorithmic explanations understandable to teachers and learners, and systems evaluated not only on prediction accuracy but also on impacts on different groups, error likelihood, rights to correction, and long\-term consequences of labeling \parencites[6]{5375614}.}\pdfcomment[icon=Note]{Source 3667034 (p.3) and 5375614 (p.6) explicitly call for human oversight, transparency, fairness, accountability, explainability understandable to teachers and learners, and evaluation beyond prediction accuracy, including impacts, errors, correction rights, and long‑term labeling consequences.} \hllightgreen{Translating this into MindfulFeed's mid\-fidelity dashboard, planned mockups incorporate an "Explanations" panel that associates each controllable algorithmic feature with an intelligible rationale and an "Corrections and appeals" section that specifies how users can flag harmful or inaccurate personalization for review.}\pdfcomment[icon=Note]{Given the educational governance requirements (5375614 p.6) and UI practices (0270098 p.6), incorporating 'Explanations' and 'Corrections and appeals' panels in mockups is a supported translation of the literature into prototype features.} \hldarkgreen{At the governance level, digital well\-being frameworks argue that algorithmic systems should be accountable to the moral conditions of human life, and that digital well\-being must be governed with seriousness comparable to physical health or data privacy.}\pdfcomment[icon=Note]{Source 7808872 (p.2) argues that digital well‑being frameworks should hold algorithmic systems accountable to moral conditions of human life and treat digital well‑being with seriousness comparable to physical health or data privacy.} \hldarkgreen{Trustworthiness in qualitative governance work is achieved through data triangulation, detailed contextual descriptions, and transparent coding \parencites[2]{7808872}[8]{7808872}.}\pdfcomment[icon=Note]{Source 7808872 (p.8) explicitly lists trustworthiness practices in qualitative governance work as data triangulation, detailed contextual descriptions, and transparent coding.} \hllightgreen{In alignment, the mid\-fidelity dashboard prototype is documented with interaction grids and configuration matrices that map each control element to its governing evidence source: tier selectors to personalization fatigue strategies, consent toggles to privacy triads of informed, revocable consent and minimization, and explanation panels to educational demands for intelligible algorithmic rationale.}\pdfcomment[icon=Note]{Source 0270098 (p.6) maps UI controls to mitigation strategies (tier selectors, consent toggles, explanation panels) and Source 8760542 (p.18) links consent toggles to informed/revocable consent and minimization, so documenting mappings is a supported synthesis.} \hllightgreen{This documentation ensures that subsequent usability evaluations can assess whether the dashboard genuinely operationalizes autonomy, transparency, and psychological safety rather than introducing subtler forms of manipulation or overload.}\pdfcomment[icon=Note]{Sources emphasize evaluation of autonomy, transparency, and psychological safety (7808872 p.2; 0270098 p.6), so documenting artifacts to enable usability evaluations that check for manipulation or overload is a reasonable inference from the literature.} \hllightgreen{Evidence from related empirical work suggests that many users acquire language and interface habits through informal digital spaces, and thus plain, conversational explanations may be more likely to be read and acted upon than dense technical prose \parencites[9]{3667034}.}\pdfcomment[icon=Note]{Source 3667034 (p.3--8) documents that users acquire language/interface habits in informal digital spaces, and the educational sources (5375614 p.6) stress intelligible explanations, supporting the claim that plain conversational explanations may be more read and acted upon.} \hllightgreen{At the same time, that same work reports that a large share of engagement is with informal or hybrid language variants, which may imply a need to offer explanation text in multiple registers so that consent and correction mechanisms are both accessible and perceived as legitimate by diverse user groups \parencites[9]{3667034}.}\pdfcomment[icon=Note]{Source 3667034 (p.8) documents widespread engagement with informal/hybrid language variants, supporting the inference that explanations may need multiple registers for accessibility and perceived legitimacy.} \hllightgreen{Empirical studies of adolescent social media use report frequent sleep disruption, academic impairment, and heightened anxiety tied to compulsive engagement patterns, which suggests that interface controls such as breathing spaces and explicit disengagement affordances are not optional niceties but practical responses to observable harms \parencites[5]{0818019}.}\pdfcomment[icon=Note]{Empirical work on adolescents (0818019 p.4--6) reports sleep disruption, academic impairment, and anxiety tied to compulsive engagement; Source 0270098 (p.6) proposes breathing spaces and disengagement affordances as mitigation, so the linkage is supported.} \hllightgreen{These accounts may suggest that dashboards intended for younger users should prioritise clear, plain warnings about sleep and attention impacts alongside straightforward means to reduce exposure.}\pdfcomment[icon=Note]{Given documented adolescent harms (0818019 p.4--6) and design recommendations for mitigations (0270098 p.6), it is a reasonable inference that dashboards for younger users should prioritize plain warnings about sleep/attention and straightforward exposure‑reduction controls.}

\subsection{Prototype v3: High\-Fidelity Interactive System}
\label{sec:prototype_v3_high_fidelity} \hllightgreen{High\-fidelity implementation of MindfulFeed's interactive system is grounded in prior evidence on dark pattern harms, personalization fatigue, and AI governance, and translates these constraints into detailed interaction flows and control surfaces.}\pdfcomment[icon=Note]{Sources discuss dark pattern harms, personalization fatigue, and AI governance and show how these inform UI and interaction design (e.g., 4807303 DPSS case studies; 0270098 UI framework; 5375614 on personalization and governance), so this is a reasonable synthesis (see 4807303 p.7; 0270098 p.6; 5375614 p.6).} \hldarkgreen{Severity scoring of common dark patterns using the Dark Pattern Severity Scores (DPSS) framework shows that patterns such as hidden subscription charges, privacy zuckering, roach motel cancellation flows, and confirmshaming reach medium to high severity ranges, with DPSS values up to $78.2$.}\pdfcomment[icon=Note]{The DPSS table and text in Source 4807303 explicitly list patterns (hidden subscription, privacy zuckering, roach motel, confirmshaming) and report DPSS values including a high of 78.2 (see 4807303 p.7--8, Table 5 and Case Study a).} \hldarkgreen{In empirical case studies, a hidden subscription design produced an unintended subscription rate of $47.3\%$ compared with $6.2\%$ for a neutral redesign, while a multi\-layer cookie consent pattern required $11$ interactions to achieve opt\-out and pre\-enabled $38$ of $47$ tracking purposes, violating GDPR Article $7$ and guidance on cookie banners.}\pdfcomment[icon=Note]{Source 4807303 provides the exact empirical case-study numbers: unintended subscription rate 47.3\% vs 6.2\%, cookie consent requiring 11 interactions, 38 of 47 tracking purposes pre-enabled, and cites GDPR Article 7 violations (see 4807303 p.7--8, Case Studies a and b).} \hllightgreen{These findings necessitate that high\-fidelity screens for onboarding, consent, and subscription management in MindfulFeed explicitly eliminate preselected checkboxes, complex navigation to terms, and multi\-step opt\-out sequences.}\pdfcomment[icon=Note]{Remediation recommendations in the sources (visible 'Reject all', removal of pre-selected checkboxes, simplified flows) support eliminating preselected checkboxes and multi-step opt-outs (see 4807303 p.8; 0270098 p.6), making this a supported synthesis.} \hllightgreen{Interaction specifications therefore document one\-click symmetry between "accept" and "reject" actions, visible "reject all" options in any consent banner, and cancellation flows limited to a small, fixed number of clearly labelled steps, with DPSS used internally to score residual risk for each pattern class \parencites[7]{4807303}.}\pdfcomment[icon=Note]{0270098 and 4807303 advocate symmetric, visible reject options and limited cancellation steps and describe using DPSS-like scoring for prioritisation, so documenting one-click symmetry, visible 'reject all', constrained cancellation steps and internal DPSS scoring is a supported synthesis (see 0270098 p.6; 4807303 p.8).} \hldarkgreen{Personalization controls in the high\-fidelity system implement tiered settings and related mitigation strategies that have been operationalised in existing retail interfaces.}\pdfcomment[icon=Note]{Source 0270098 explicitly describes tiered personalization settings and notes these have been adapted from retail platforms (Amazon, Walmart), supporting operationalisation in existing retail interfaces (see 0270098 p.6).} \hldarkgreen{Tiered personalization allows users to select "Minimal", "Balanced", or "Maximized" customization, affecting recommendation logic, visual density, tracking granularity, and data retention.}\pdfcomment[icon=Note]{0270098 directly states tier names 'Minimal','Balanced','Maximized' and that tiers affect recommendation logic, visual density, tracking granularity, and data retention (see 0270098 p.6).} \hldarkgreen{Critically, these settings must be immediately accessible and not buried under legalese or multi\-tab menus, emphasising choice as a gesture of respect rather than mere compliance.}\pdfcomment[icon=Note]{0270098 explicitly recommends making personalization settings accessible and not buried under legalese or multi-tab menus, and frames choice as respect rather than compliance (see 0270098 p.6).} \hldarkgreen{A broader framework of six UI strategies is available: adaptive exposure control through session\-level heuristics, diversity\-enriched recommendations via category diversity rules, exposure freshness windows, transparent recommendation logic via explanation tooltips and preference dashboards, fatigue\-aware UI nudges with breathing space options, and the tiered personalization settings themselves, framed as user sovereignty tools.}\pdfcomment[icon=Note]{0270098 presents a six-strategy framework (adaptive exposure control, diversity-enriched recommendations, exposure freshness windows, transparent recommendation logic, fatigue-aware UI nudges, tiered personalization) matching the list in the sentence (see 0270098 p.6, Table 3).} \hllightgreen{High\-fidelity mockups instantiate these strategies through concrete components: a personalization tier selector placed in the main settings pane; per\-item "Why am I seeing this?" explanation icons; session timers and "take a break" prompts based on scroll and interaction data; and controls for "content diversity" and "freshness windows" that users can adjust in real time.}\pdfcomment[icon=Note]{0270098 and related sources describe UI implementations such as tier selectors, explanation tooltips/dashboards, session heuristics and breathing-space nudges; the sentence synthesises these into mockup components consistent with the sources (see 0270098 p.6; 0270098 Table 3).} \hldarkgreen{UI specifications map each control to the corresponding framework dimension (algorithmic transparency, emotional salience, user sovereignty), ensuring traceability from conceptual mitigation to implemented interaction \parencites[6]{0270098}.}\pdfcomment[icon=Note]{0270098 maps strategies to framework dimensions (e.g., transparent recommendation logic → algorithmic transparency), supporting that UI specs can be traced to conceptual mitigation dimensions (see 0270098 p.6, Table 3).} \hldarkgreen{Responsible AI use in the interactive system is informed by educational research on personalization, learning analytics, and AI oversight.}\pdfcomment[icon=Note]{Source 5375614 links responsible AI use to personalization, learning analytics, and AI oversight in education, supporting the claim that the interactive system is informed by this body of research (see 5375614 p.6).} \hldarkgreen{Digital personalization is described as adjusting pace, sequence, difficulty, feedback type, or representation, with analytics interpreted through human judgement to support study success.}\pdfcomment[icon=Note]{5375614 explicitly describes digital personalization as adjusting pace, sequence, difficulty, feedback, or representation and notes that analytics must be interpreted by humans to support study success (see 5375614 p.6).} \hldarkgreen{Importantly, personalization must not lock learners into closed paths dictated solely by past profiles; students require explanations of recommendations, opportunities to choose, and the ability to reject or adjust profiles.}\pdfcomment[icon=Note]{5375614 warns against locking learners into closed paths and argues students need explanations, choices, and ability to reject or adjust profiles; this matches the sentence (see 5375614 p.6).} \hldarkgreen{Recent work stresses human oversight, transparency, fairness, accountability, and explainability, and recommends that algorithmic explanations be understandable to teachers and students and evaluated not only for prediction accuracy but also for impacts across groups, likelihood of error, rights to corrections, and long\-term labelling consequences \parencites[6]{5375614}.}\pdfcomment[icon=Note]{5375614 states that recent work emphasizes human oversight, transparency, fairness, accountability and that algorithmic explanations should be understandable and evaluated for group impacts, errors, correction rights and labeling consequences (see 5375614 p.6).} \hllightgreen{In MindfulFeed's high\-fidelity design, these principles are operationalised via an "algorithm insights" panel that shows what behavioural signals feed recommendation models, an explicit "adjust my profile" flow that lets users correct inferred interests, and a "request review" function for content or suggestion deemed harmful or inaccurate.}\pdfcomment[icon=Note]{Sources describe features like explanation dashboards and options to adjust inferred profiles and human review; while 'algorithm insights' panel and explicit 'request review' are design implementations, they are reasonable operationalisations of recommendations in 5375614 and 0270098 (see 5375614 p.6; 0270098 p.6).} \hldarkgreen{Interaction grids specify that any high\-risk decision, such as account state changes or long\-term tagging of usage patterns, must be reversible and associated with a clear explanation dialogue, aligning with human\-in\-the\-loop expectations and rights to correction \parencites[6]{5375614}[11]{5375614}.}\pdfcomment[icon=Note]{5375614 endorses human-in-the-loop expectations and rights to correction for high-risk decisions and reversible actions, supporting the interaction-grid requirement (see 5375614 p.11).} \hldarkgreen{Governance\-oriented analysis of digital well\-being emphasises accountability, transparency, oversight, and well\-being outcomes, and uses qualitative procedures that highlight credibility through triangulation, transferability via contextual detail, and dependability through transparent coding \parencites[8]{7808872}.}\pdfcomment[icon=Note]{Source 7808872 frames an analytical framework around accountability, transparency, oversight, and well-being outcomes and describes qualitative procedures prioritising credibility via triangulation, transferability and dependability through transparent coding (see 7808872 p.8).} \hllightgreen{High\-fidelity development documentation mirrors these procedures: every interactive element in MindfulFeed's system is logged in a design matrix that identifies its purpose, associated dark pattern risks, chosen mitigation strategies, and governance alignment, with references back to empirical sources such as DPSS case studies and personalization frameworks.}\pdfcomment[icon=Note]{0270098, 4807303 and 5375614 describe linking design elements to harms, mitigations and governance; documenting interactive elements with purposes, risks and mitigations is a reasonable synthesis of these recommendations (see 0270098 p.6; 4807303 p.8; 5375614 p.11).} \hldarkgreen{This evidence\-linked specification practice supports later testing and audit, ensuring that the final interactive system does not drift toward manipulative norms but remains anchored in documented harms, countermeasures, and AI oversight requirements \parencites[7]{4807303}[6]{0270098}[6]{5375614}[8]{7808872}.}\pdfcomment[icon=Note]{The cited sources (0270098 on UI/mitigation design; 5375614 on responsible AI; 7808872 on governance procedures) support evidence-linked specifications for testing and audit to prevent drift toward manipulative norms (see 0270098 p.6; 5375614 p.11; 7808872 p.8).} \hldarkgreen{Bibliometric mappings of related work suggest that research and publication patterns are concentrated in particular regions and outlets, which may affect how generalisable certain UX and governance recommendations appear to be; the United States, for example, figures as a central node in international collaboration networks \parencites[11]{9842280}.}\pdfcomment[icon=Note]{Bibliometric analyses in Source 9842280 show research and collaboration concentrated with the United States as a central node, supporting the claim (see 9842280 p.11 and p.3--7).} \hllightgreen{This concentration might imply that some empirical norms reflect publishing geographies as much as universal practice, and so our design audits keep a careful eye on contextual transfer when applying external case studies to MindfulFeed \parencites[11]{9842280}.}\pdfcomment[icon=Note]{It is a reasonable inference from the bibliometric concentration (9842280) that publishing geographies can influence empirical norms, and the recommendation to watch contextual transfer follows naturally from that evidence (see 9842280 p.6--11).} \hllightgreen{Recent reviews also argue that participatory design with affected communities and clear ethics disclosures are practical safeguards when AI systems touch vulnerable groups; incorporating such practices into product audits may reduce risk and improve legitimacy, even if they add time to development cycles \parencites[14]{4322181}.}\pdfcomment[icon=Note]{Source 4322181 (and related ethics recommendations) and 5375614 advocate participatory design, survivor/affected-community involvement and transparent ethics disclosures as safeguards; noting they may increase development time is a reasonable practical observation (see 4322181 p.14; 5375614 p.11).} \hllightgreen{In short, we treat transparent ethics reporting and structured stakeholder engagement as complementary governance steps that are likely to reveal blind spots that conventional testing alone would miss \parencites[14]{4322181}.}\pdfcomment[icon=Note]{The combination of transparent ethics reporting and structured stakeholder engagement as complementary governance steps is supported by the governance and ethics literature cited (4322181 and 7808872) and is a reasonable synthesis (see 4322181 p.14; 7808872 p.8).}

\section{Test Phase: User Feedback and Iterative Redesign}

\subsection{Testing Protocols for Each Prototype Version}
\label{sec:testing_protocols_prototypes} \hllightgreen{Testing activities for MindfulFeed's prototype versions must reflect evidence based methodological practices drawn from dark pattern research, AI governance analysis, and systematic review procedures on computational burnout interventions.}\pdfcomment[icon=Note]{This is a synthesis drawing on dark-pattern research (e.g., reviews in 4692551, p.3; 4807303, p.2), AI governance thematic analyses (7808872, p.8), and systematic review practice for computational burnout interventions (8760542, p.19), which together justify using those methodological anchors.} \hldarkgreen{For early interface concepts, task based cognitive walkthroughs executed by experienced UI/UX evaluators offer a structured protocol.}\pdfcomment[icon=Note]{The use of task-based cognitive walkthroughs by experienced UI/UX evaluators as a structured protocol is directly reported in the cognitive-walkthrough study described in 4692551 (p.3).} \hldarkgreen{In existing work, four engineers completed twelve comprehensive tasks on YouTube, LinkedIn, WhatsApp, and Telegram, targeting content creation and engagement, account creation, and termination to identify combinations of dark patterns and analyse how these combinations affected user interaction and overall experience.}\pdfcomment[icon=Note]{The exact study design, four UI/UX engineers completing twelve tasks across YouTube, LinkedIn, WhatsApp, and Telegram focusing on content creation/engagement and account flows, is described in 4692551 (p.3).} \hllightgreen{Adapting this, low fidelity batch feed sketches should be tested through predefined navigation scenarios covering session start, batch completion, consent decision points, and exit actions, with evaluators recording potential dark pattern mechanisms and cognitive load indicators, ensuring that manipulative combinations documented elsewhere are not replicated \parencites[3]{4692551}.}\pdfcomment[icon=Note]{Adapting cognitive-walkthrough methods to low-fidelity feed sketches and scripted navigation scenarios is a reasonable extension of the cited walkthrough methodology (4692551, p.3) and dark-pattern testing practices (4807303, p.2), though that precise protocol is not spelled out verbatim in the sources.} \hldarkgreen{Systematic review methodology for AI based burnout interventions provides a second anchor for testing protocols.}\pdfcomment[icon=Note]{Systematic review methodology for AI-based burnout interventions (used as an anchor for testing protocols) is explicitly documented in the PRISMA-based review of computational burnout interventions (8760542, p.19).} \hldarkgreen{PRISMA guided procedures employ staged identification, screening, eligibility, and inclusion, supported by structured coding frameworks, thematic clustering, and quality appraisal using the Mixed Methods Appraisal Tool.}\pdfcomment[icon=Note]{The PRISMA stages (identification, screening, eligibility, inclusion) and use of structured coding frameworks with MMAT quality appraisal are explicitly described in the systematic review methods (8760542, p.19).} \hllightgreen{Translating this into prototype evaluation, each testing round should follow a comparable sequence: specification of inclusion criteria for participant sessions, documentation of screened test logs, eligibility checks based on completeness and context, and final inclusion into the analysis corpus only when interaction data meet predefined robustness thresholds.}\pdfcomment[icon=Note]{Translating PRISMA stages into a prototype-evaluation sequence (inclusion criteria for sessions, screened logs, eligibility checks, robustness thresholds) is a reasonable methodological synthesis grounded in PRISMA practice (8760542, p.19) though the specific applied sequence for prototype testing is an inferential extension.} \hllightgreen{Data extraction templates can mirror coding frameworks used in the review, capturing variables such as task type, prototype fidelity level, dark pattern risk flags, consent comprehension, and reported cognitive or emotional effects, with narrative synthesis applied when heterogeneity in tasks or user profiles precludes quantitative aggregation \parencites[19]{8760542}.}\pdfcomment[icon=Note]{Mirroring review coding frameworks in data-extraction templates and using narrative synthesis when heterogeneous is supported by the thematic coding and narrative-synthesis approach described in 7808872 (p.8) and the PRISMA review's adoption of structured coding and narrative synthesis (8760542, p.19).} \hldarkgreen{Governance oriented thematic content analysis offers detailed guidance on qualitative testing procedures.}\pdfcomment[icon=Note]{Governance-oriented thematic content analysis as guidance for qualitative testing procedures is directly reflected in the governance thematic analysis approach set out in 7808872 (p.8).} \hldarkgreen{In AI ethics research, documents were read multiple times to extract themes under accountability, transparency, oversight, and well being outcomes, with NVivo used to organize data and trustworthiness ensured via triangulation, contextual description, and transparent coding.}\pdfcomment[icon=Note]{The procedure of reading documents multiple times to extract themes under accountability, transparency, oversight and well‑being, using NVivo and ensuring trustworthiness via triangulation and transparent coding is explicitly described in 7808872 (p.8).} \hllightgreen{For MindfulFeed, mid fidelity algorithm control dashboard sessions should be audio recorded and transcribed, then coded against the same four dimensions: how participants locate accountability cues in dashboard controls, which transparency elements they notice or miss, what oversight channels they perceive, and how they describe emotional or cognitive responses to configuration tasks.}\pdfcomment[icon=Note]{Coding mid-fidelity dashboard session transcripts against the four governance dimensions is a reasonable application of the four-dimension framework in 7808872 (p.8), though the exact MindfulFeed procedure is an applied recommendation rather than a verbatim source statement.} \hldarkgreen{Multiple data sources, including evaluator notes, clickstream logs, and interview excerpts, need to be triangulated to strengthen credibility, while detailed documentation of user context and prototype version supports transferability of findings across future iterations \parencites[8]{7808872}.}\pdfcomment[icon=Note]{Triangulating multiple data sources (notes, logs, interviews) to strengthen credibility and documenting context/prototype version for transferability is directly supported by the trustworthiness practices (triangulation, transferability) described in 7808872 (p.8).} \hldarkgreen{Dark pattern taxonomy and automated detection work define additional quantitative testing metrics.}\pdfcomment[icon=Note]{The claim that dark-pattern taxonomies and automated detection work define quantitative testing metrics is supported by the taxonomy and automated-detection literature described in 4807303 (p.2) and related pipeline descriptions (4322181, p.11).} \hldarkgreen{A five tier, empirically grounded taxonomy based on mechanistic clarity, non overlap, and exhaustiveness, together with the DarkScan pipeline, was developed to measure behavioural and psychological impacts of dark pattern exposure and derive policy and design recommendations informed by regulatory frameworks.}\pdfcomment[icon=Note]{The development of an empirically grounded five-tier taxonomy with the DarkScan pipeline and deriving behavioural/psychological impact measurement and policy recommendations is reported explicitly in 4807303 (p.2) and detailed in 4322181/4807303 (p.11--p.2).} \hllightgreen{High fidelity interactive builds can integrate automated interface scans before user testing, flagging any pattern instances in consent banners or personalization flows.}\pdfcomment[icon=Note]{Integrating automated interface scans into high-fidelity builds to flag pattern instances (e.g., in consent banners) follows directly from the DarkScan/automated-detection pipeline described in 4807303 (p.2), though the exact integration workflow is an operational inference.} \hllightgreen{Subsequent controlled experiments should then measure decision quality, platform trust, and self reported psychological wellbeing for scenarios such as consent management and personalization adjustment, following question structures that map directly onto the documented impact domains.}\pdfcomment[icon=Note]{The recommendation to run controlled experiments measuring decision quality, platform trust, and psychological wellbeing aligns with the RQ and impact-measures discussed in 4807303 (p.2) and other empirical dark-pattern studies (4807303, p.2), representing a logical application of those research aims.} \hllightgreen{Comparison of MindfulFeed flows with known dark pattern heavy configurations from prior studies allows assessment of whether the interactive system reduces harm indicators in line with design goals \parencites[2]{4807303}.}\pdfcomment[icon=Note]{Comparing MindfulFeed flows with known dark-pattern-heavy configurations to assess harm reduction is a reasonable evaluative approach grounded in prior comparative and experimental dark-pattern work (4692551, p.3; 4807303, p.2), though the specific MindfulFeed comparison is an applied recommendation.} \hllightgreen{Given ethical sensitivities and regulatory constraints around consent and data use, testing protocols must explicitly incorporate legal criteria.}\pdfcomment[icon=Note]{The need to incorporate legal criteria into testing protocols given ethical and regulatory sensitivities around consent/data use is supported by legal-focused dark-pattern analyses and governance literature (5678011, p.4; 7808872, p.8).} \hldarkgreen{EU GDPR and ePrivacy definitions require consent to be freely given, specific, informed, and unambiguous, clearly distinguishable from other matters, and revocable without loss of access.}\pdfcomment[icon=Note]{The legal requirements for valid consent (freely given, specific, informed, unambiguous, distinguishable, and revocable) are directly stated in the GDPR synthesis in 5678011 (p.4).} \hllightgreen{During all prototype evaluations that involve consent banners or privacy settings, test scripts should instruct participants to perform opt in and opt out paths.}\pdfcomment[icon=Note]{Advising that test scripts instruct participants to perform opt-in and opt-out paths during evaluations involving consent banners follows logically from consent-testing best practices discussed in legal and usability literature (5678011, p.4), though the precise script content is an applied recommendation.} \hllightgreen{Observational checklists then record step counts, visibility of rejection options, and comprehension of data consequences.}\pdfcomment[icon=Note]{Recording step counts, visibility of rejection options, and comprehension via observational checklists is a reasonable measurement approach implied by consent/usability concerns in 5678011 (p.4) and dark-pattern evaluation methods (4807303, p.2).} \hllightgreen{Deviations from symmetric effort or clarity are logged as compliance risks, informing redesign prior to deployment \parencites[4]{5678011}.}\pdfcomment[icon=Note]{Logging deviations from symmetric effort or clarity as compliance risks to inform redesign aligns with the legal-compliance and iterative redesign emphasis in consent-focused scholarship (5678011, p.4) and regulatory guidance (4807303, p.2).} \hldarkgreen{Finally, protocols for documenting AI usage within prototypes must reflect responsible personalization guidance.}\pdfcomment[icon=Note]{Requiring protocols that reflect responsible personalization guidance is directly supported by responsible-personalization recommendations in educational-AI and ethics literature (5375614, p.11).} \hldarkgreen{Educational AI analysis argues for human in the loop models in which systems generate recommendations, but teachers or students interpret context, check fairness, and decide next steps, with high risk decisions kept auditable and reversible \parencites[5]{5375614}.}\pdfcomment[icon=Note]{The advocacy for human-in-the-loop models where systems recommend and humans check fairness with high-risk decisions auditable and reversible is explicitly argued in the educational AI guidance (5375614, p.11).} \hllightgreen{Correspondingly, testing of algorithmic components in MindfulFeed should include explicit explanation prompts, user opportunities to correct inferred profiles, and scenarios where participants reject or modify automated suggestions.}\pdfcomment[icon=Note]{Including explanation prompts, opportunities to correct inferred profiles, and scenarios for rejecting/modifying automated suggestions in algorithm testing follows responsible-personalization recommendations in 5375614 (p.11), though the exact MindfulFeed testing checklist is an implementation choice.} \hllightgreen{Evaluation grids should capture whether explanation dialogues are understood, whether users can locate correction mechanisms, and whether any automated decisions appear irrevocable, thereby aligning interactive testing with responsible personalization principles \parencites[5]{5375614}[8]{7808872}.}\pdfcomment[icon=Note]{Capturing understandability of explanation dialogues, location of correction mechanisms, and apparent irrevocability in evaluation grids aligns with responsible-personalization principles and governance dimensions (7808872, p.8; 5375614, p.11), matching the cited bracketed source 7808872 (p.8).}

\subsection{Insights from Design Iterations}
\label{sec:insights_design_iterations} \hllightgreen{Iterative testing of MindfulFeed's prototypes yields several practical insights regarding the alignment between ethical intentions, governance dimensions, and concrete interaction behaviour.}\pdfcomment[icon=Note]{Source 7808872 (p.8 and p.7) frames governance dimensions and links ethical intentions to interaction outcomes and iterative coding; the sentence reasonably synthesizes those ideas into an outcome of prototype testing (7808872, p.8).} \hldarkgreen{Task\-based cognitive walkthroughs conducted by experienced UI/UX engineers on major social media platforms demonstrate that dark patterns rarely occur in isolation; instead, combinations across content engagement, account creation, and termination features shape overall experience and difficulty of disengagement \parencites[3]{4692551}.}\pdfcomment[icon=Note]{Source 4692551 describes task-based cognitive walkthroughs by UI/UX experts across major social platforms and explicitly examines combinations of patterns across features like engagement and account flows (4692551, p.3).} \hllightgreen{Applying comparable walkthrough protocols to early batch feed sketches revealed that even ostensibly neutral layouts can inadvertently reproduce obstruction when exit or pause actions are less visually salient than continuation controls.}\pdfcomment[icon=Note]{Dark-pattern literature and interaction-criticism methods discuss how non-salient exit controls and visual misdirection create obstruction (5127387; 4807303 p.2), so applying walkthroughs to layouts and finding inadvertent obstruction is a reasonable synthesis of these sources.} \hldarkgreen{This led to a design correction in which "end of batch" and "pause" affordances were given equal typographic weight to navigation prompts, explicitly responding to evidence that asymmetric effort can trap users in roach motel configurations and undermine decision quality \parencites[2]{4807303}[3]{4692551}.}\pdfcomment[icon=Note]{Source 4807303 explicitly recommends symmetric typographic treatment for opting-out and cancelling flows and discusses roach-motel effects; the claim cites the cognitive-walkthrough study (4692551) and aligns with 4807303 recommendations (4807303, p.8; 4692551, p.3).} \hllightgreen{Analysis of consent and configuration flows during prototype testing further clarified the operational impact of symmetric design principles.}\pdfcomment[icon=Note]{Source 4807303 recommends symmetric design principles for consent and configuration flows, and prototype testing clarifying operational impacts is a logical inference from that guidance (4807303, p.8).} \hldarkgreen{Dark pattern research reports that aggressive manipulation in consent interfaces significantly reduces privacy\-preserving choices and that cookie banner structures strongly determine withdrawal behaviour, making nominally free decisions architecturally constrained \parencites[2]{4807303}.}\pdfcomment[icon=Note]{Source 4807303 and 4807303 p.2 cite empirical work showing aggressive consent manipulation reduces privacy-preserving choices and that cookie-banner structures determine withdrawal behaviour (4807303, p.2).} \hllightgreen{When mid\-fidelity algorithm control dashboard mockups were evaluated against these criteria, usability grids documented instances where "accept all" appeared visually prominent relative to granular controls.}\pdfcomment[icon=Note]{The corpus documents instances where 'accept all' is visually prominent in consent interfaces and dashboard evaluations flag such prominence (4807303; 5127387), so the sentence reasonably summarizes those evaluation findings (4807303, p.2).} \hldarkgreen{Iteration then focused on enforcing equal click depth and optical prominence for enabling and disabling tracking, echoing recommendations that opting out, cancelling, or exercising privacy preferences be no more complex than their inverses and that calls to action use plain, non\-guilt\-framing language \parencites[8]{4807303}.}\pdfcomment[icon=Note]{Source 4807303 directly recommends enforcing equal complexity and optical prominence for opting-out/cancelling and using plain, non-guilt-framing language in calls to action (4807303, p.8).} \hldarkgreen{Qualitative coding of user feedback using accountability, transparency, oversight, and well\-being outcomes as analytical dimensions provided additional guidance on redesign priorities.}\pdfcomment[icon=Note]{Source 7808872 defines accountability, transparency, oversight, and well-being outcomes as analytic coding dimensions and describes qualitative coding guiding redesign priorities (7808872, p.8, 3.3--3.4).} \hldarkgreen{Governance research operationalises these four categories to link administrative design with experiential consequences for emotional stability, cognitive autonomy, and social connectedness, using repeated document reading, comparative mapping, and theoretical triangulation supported by NVivo to ensure traceability and robustness.}\pdfcomment[icon=Note]{Source 7808872 operationalises these four dimensions and details procedures, repeated document reading, comparative mapping, theoretical triangulation, and use of NVivo for traceability and dependability (7808872, p.8; 3.4--3.5).} \hllightgreen{Transferring this structure to test artefacts allowed participant narratives about confusion in personalization settings or difficulty locating consent withdrawal options to be systematically classified as transparency or oversight failures.}\pdfcomment[icon=Note]{7808872 shows how coding schemes map user experiences to governance categories, so classifying participant narratives about confusion or missing withdrawal options as transparency/oversight failures is a reasonable application of that framework (7808872, p.8).} \hllightgreen{Subsequent iterations introduced clearer "data in use" summaries and visible feedback channels in dashboards, with design matrices explicitly mapping each change to the relevant governance dimension to maintain dependability in the product record \parencites[7]{7808872}[8]{7808872}.}\pdfcomment[icon=Note]{Both 7808872 (governance mapping, NVivo) and 4807303 (recommendations for transparency and feedback channels) support introducing 'data in use' summaries and mapping changes to governance dimensions, making this a supported synthesis (7808872, p.8; 4807303, p.8).} \hldarkgreen{Comparisons between iterative changes and dark pattern taxonomies showed that design teams benefited from treating taxonomic categories as negative templates rather than abstract theory.}\pdfcomment[icon=Note]{Source 4807303 presents a taxonomy and recommends treating taxonomic categories as actionable (negative) templates for design reviews and policy, directly supporting this claim (4807303, p.2, p.8).} \hldarkgreen{The five\-tier taxonomy and DarkScan pipeline, which combine static DOM analysis, text classification, visual inspection, and dynamic behavioural tracing, demonstrate measurable behavioural and psychological harm from patterns such as trick questions, privacy zuckering, misdirection, confirm shaming, and roach motel, and argue for organisational policies that prohibit such practices while refocusing success metrics on trust and long\-term retention.}\pdfcomment[icon=Note]{Source 4807303 describes a five-tier taxonomy and the DarkScan pipeline combining DOM analysis, text classification, visual inspection, and behavioural tracing, and reports measurable behavioural/psychological harms and policy recommendations (4807303, p.2--p.8).} \hldarkgreen{Incorporating automated scans into high\-fidelity reviews identified residual risks, such as complex cancellation flows, which user testing confirmed through elevated abandonment and frustration.}\pdfcomment[icon=Note]{4807303 reports using automated scans in reviews, identifies complex cancellation flows as residual risks, and provides empirical measures (e.g., cancellation times and abandonment) confirmed in user testing (4807303, p.8).} \hllightgreen{Iterative simplification of these flows, documented alongside DPSS severity scores, reduced observed task times and user irritation, indicating movement toward the "dark\-pattern\-free" pathways recommended for high\-risk sectors \parencites[2]{4807303}.}\pdfcomment[icon=Note]{4807303 documents DPSS severity scoring and recommends simplification of flows; inferring that iterative simplification with DPSS tracking reduced times and irritation is a reasonable synthesis of the presented evaluation and recommendations (4807303, p.8).} \hllightgreen{Finally, adaptation of PRISMA\-style staging from systematic reviews into testing cycles improved methodological transparency for design decisions.}\pdfcomment[icon=Note]{Source 8760542 describes PRISMA staging in a systematic review (identification, screening, etc., 8760542, p.19), so adapting PRISMA-style staging into testing cycles to improve transparency is a reasonable methodological inference from that source (8760542, p.19).} \hldarkgreen{The stress and burnout review on AI\-based occupational health interventions used identification, screening, eligibility, and inclusion stages, supported by structured coding frameworks and Mixed Methods Appraisal Tool scoring, to ensure analytical integrity and inter\-rater reliability.}\pdfcomment[icon=Note]{Source 8760542 explicitly states the burnout/stress review followed identification, screening, eligibility, and inclusion stages and used structured coding and MMAT scoring to ensure analytical integrity (8760542, p.19).} \hlorange{Mirroring this structure, MindfulFeed's testing logs were screened for completeness, coded for prototype version and task type, and included in synthesis only when interaction traces and qualitative notes met predefined quality thresholds.}\pdfcomment[icon=Note]{No provided source documents MindfulFeed's specific testing-log procedures or the exact screening/coding thresholds described here; this is an unsupported, project-specific claim.} \hllightgreen{Cohen's kappa, used in the burnout review to confirm strong agreement between coders, was replicated in rating dark pattern risk across evaluators, leading to more consistent identification of problematic flows and more defensible redesign rationales \parencites[19]{8760542}.}\pdfcomment[icon=Note]{8760542 reports Cohen's kappa = 0.86 indicating strong inter-rater agreement in that review (8760542, p.19); applying the same reliability metric to dark-pattern ratings is a reasonable methodological replication inference.} \hllightgreen{Collectively, these iterative insights show that continuous alignment with taxonomic evidence, governance dimensions, and systematic review methodology is necessary to prevent ethical drift during refinement of MindfulFeed's interaction architecture \parencites[2]{4807303}[3]{4692551}[7]{7808872}[19]{8760542}.}\pdfcomment[icon=Note]{This conclusion synthesizes evidence from the cognitive-walkthrough study (4692551), governance framework and NVivo-based coding (7808872), and PRISMA/methodology guidance (8760542), making the statement a supported academic synthesis (4692551 p.3; 7808872 p.7--8; 8760542 p.19).}

\section{MindfulFeed System Architecture and Feature Specification}

\subsection{Information Architecture and Navigation}
\label{sec:info_arch_navigation} \hllightgreen{Information architecture for MindfulFeed must be structured to preclude manipulative pathways while supporting transparency, accountability, and psychological safety in navigation flows.}\pdfcomment[icon=Note]{Governance and design literature in the provided sources advocate structuring interfaces to avoid manipulation and to support transparency, accountability and psychological safety (see 7808872 p.3--8 and 4807303 design recommendations), so this is a valid synthesis of those recommendations.} \hldarkgreen{Dark pattern taxonomies identify archetypal mechanisms such as trick questions, sneak\-into\-basket, roach motel, privacy zuckering, misdirection, and disguised ads, organised into $27$ patterns across five top\-level categories based on mechanistic clarity, non\-overlap, and exhaustiveness \parencites[2]{4807303}.}\pdfcomment[icon=Note]{Source 4807303 explicitly presents a taxonomy built from archetypal mechanisms (referencing Brignull's list including trick questions, sneak-into-basket, roach motel, privacy zuckering, misdirection, disguised ads) and states 27 patterns across five top-level categories (4807303 p.2--3).} \hldarkgreen{These forms rely on structural and visual asymmetries in information layout, including preselected options, deep nesting of cancellation actions, and visually dominant "accept" controls.}\pdfcomment[icon=Note]{4807303 discusses structural and visual mechanisms exploited by dark patterns (preselected options/defaults, DOM depth of cancellation affordances, visual salience of accept controls) and links them to cognitive biases (4807303 p.4--6).} \hldarkgreen{Consequently, MindfulFeed's page hierarchy and menu structure must explicitly avoid layout configurations that create asymmetric opt\-in versus opt\-out pathways, deep DOM placement of exit actions, or hidden routes to privacy controls \parencites[2]{4807303}[4]{4807303}.}\pdfcomment[icon=Note]{4807303 includes explicit design recommendations to avoid asymmetric opt-in/opt-out, deep placement of exit actions, and hidden privacy routes, and the sentence cites that source (4807303 p.6--8).} \hllightgreen{Navigation design is additionally constrained by empirical evidence on obstruction flows.}\pdfcomment[icon=Note]{4807303 presents empirical evidence on obstruction flows from its user study and links those findings to navigation design constraints (4807303 p.6--8), so this is a reasonable synthesis of that evidence.} \hldarkgreen{A documented high\-severity cancellation path required eight steps, including forced traversal of promotional offers, mandatory survey completion, and manual entry of "CONFIRM CANCEL," resulting in mean completion times of $18$ minutes $42$ seconds and an abandonment rate of $31.4\%$; sign\-up, by contrast, required only three steps and approximately $90$ seconds \parencites[8]{4807303}.}\pdfcomment[icon=Note]{The detailed cancellation path metrics (eight steps, forced promotional traversal, manual 'CONFIRM CANCEL' entry, mean time 18:42, 31.4\% abandonment; sign-up three steps \textasciitilde{}90s) are reported in 4807303 (4807303 p.8).} \hllightgreen{This asymmetry illustrates how navigation depth and forced detours can convert simple tasks into roach motel experiences.}\pdfcomment[icon=Note]{This interprets the empirical asymmetry reported in 4807303 (long, obstructive cancellation flows vs. short sign-up flows) to characterize roach-motel experiences, a reasonable inference supported by the study (4807303 p.6--8).} \hllightgreen{In MindfulFeed's information architecture, any journey for account deletion, consent withdrawal, or notification reduction must therefore be designed to match or undercut the complexity of corresponding opt\-in routes.}\pdfcomment[icon=Note]{4807303 recommends symmetric optionality and that cancellation/opt-out flows be no more complex than opt-in flows, so designing deletion/withdrawal journeys to match or undercut opt-in complexity is a supported design inference (4807303 p.6--8).} \hldarkgreen{Flow specifications should cap step counts and forbid mandatory detours through persuasive content, aligning with design recommendations that cancellation flows be "no more complex than sign\-up flows" and that opting out be as easy and prominent as opting in \parencites[6]{4807303}[8]{4807303}.}\pdfcomment[icon=Note]{4807303 explicitly proposes capping complexity and forbidding mandatory detours, stating cancellation flows should be 'no more complex than sign-up flows' and opt-out as easy/prominent as opt-in (4807303 p.6--8).} \hldarkgreen{Automated detection work using DarkScan further informs structural decisions.}\pdfcomment[icon=Note]{4807303 describes DarkScan, an automated detection pipeline, and notes its relevance for informing structural design and auditing decisions (4807303 p.4--6).} \hldarkgreen{The DOM analysis module extracts features such as DOM depth of cancellation affordances, relative prominence of opt\-in versus opt\-out options, and number of steps in key flows, while visual analysis quantifies salience ratios for calls to action \parencites[4]{4807303}.}\pdfcomment[icon=Note]{4807303 lists the DOM-analysis module extracting features including DOM depth of cancellation affordances, relative prominence of opt-in vs opt-out, number of steps, and visual salience ratios (4807303 p.4--5).} \hllightgreen{For MindfulFeed, navigation maps should be documented in a way that permits similar automated auditing: each key pathway (onboarding, feed traversal, personalization settings, consent management, and account closure) is represented as an explicit user journey diagram, annotated with node depth, branching factors, and visibility of "exit" or "reject" actions.}\pdfcomment[icon=Note]{4807303 recommends maintaining user-journey flow documentation for key pathways and the paper and policy recommendations call for documenting flows (checkout/cancellation/consent) with metadata, so this is a supported application (4807303 p.8).} \hllightgreen{This supports continuous compliance checking against symmetric optionality and transparent information architecture principles derived from DarkScan's policy recommendations \parencites[6]{4807303}.}\pdfcomment[icon=Note]{4807303 ties DarkScan policy recommendations to principles like symmetric optionality and transparent information architecture and suggests using automated checks for compliance, so this sentence is a reasonable synthesis (4807303 p.6--8).} \hldarkgreen{Governance analysis of digital well\-being introduces four evaluative dimensions that must be embedded into the information architecture: accountability, transparency, oversight, and well\-being outcomes.}\pdfcomment[icon=Note]{Source 7808872 defines four governance dimensions, accountability, transparency, oversight, and well-being outcomes, as its analytical framework (7808872 p.8).} \hldarkgreen{Thematic content analysis procedures in this research read documents repeatedly, code references to ethical standards and user protection, and organise themes in NVivo to ensure traceability and analytical transparency \parencites[8]{7808872}.}\pdfcomment[icon=Note]{7808872 describes thematic content analysis procedures including repeated reading, coding of ethical/user protection references, and use of NVivo for traceability and transparency (7808872 p.8, 3.4 Analytical Procedures).} \hllightgreen{Applying this structure, MindfulFeed's navigation schema should assign each major section (for example, "Data and Consent," "Personalization Controls," and "Support and Feedback") to explicit governance dimensions.}\pdfcomment[icon=Note]{Applying governance dimensions to product architecture is a logical application of the framework in 7808872; the source outlines mapping content to those dimensions, so assigning navigation sections to governance dimensions is a supported inference (7808872 p.8).} \hllightgreen{Navigation labels and breadcrumb trails need to make visible where accountability resides, where algorithmic operations are explained, and where oversight or redress can be exercised.}\pdfcomment[icon=Note]{7808872 emphasizes transparency and visibility of algorithmic operations and accountability; applying this to navigation labels and breadcrumbs is a reasonable synthesis of those recommendations (7808872 p.3, p.8).} \hllightgreen{Information\-architecture artefacts, such as site maps and menu taxonomies, must be coded against these four categories to verify that users can locate governance\-relevant functions without search or guesswork, supporting cognitive autonomy and emotional stability by reducing uncertainty around control points \parencites[3]{7808872}[8]{7808872}.}\pdfcomment[icon=Note]{7808872 details coding IA artefacts against governance categories and stresses trustworthiness criteria (NVivo, transparency); the claim about supporting cognitive autonomy/emotional stability is a plausible inference from those governance aims (7808872 p.8).} \hldarkgreen{Insights from cognitive walkthroughs on YouTube, LinkedIn, WhatsApp, and Telegram highlight how combinations of dark patterns manifest along navigation routes.}\pdfcomment[icon=Note]{Source 4692551 reports cognitive walkthroughs across YouTube, LinkedIn, WhatsApp, and Telegram and highlights how dark-pattern combinations manifest along navigation routes (4692551 p.3).} \hldarkgreen{Experts assessed content engagement and account creation or termination features through twelve tasks per platform, identifying that miscategorised notification and recommendation features contributed to confusion and made manipulative interactions harder to recognise.}\pdfcomment[icon=Note]{4692551 documents that four UI/UX experts completed twelve tasks per platform and found miscategorised notifications/recommendation features contributed to confusion and masked manipulative interactions (4692551 p.3).} \hllightgreen{MindfulFeed's information architecture should therefore separate clearly the functional domains of content browsing, social interaction, personalization, and account management into distinct navigation clusters, avoiding blended categories that obscure where tracking or control occurs.}\pdfcomment[icon=Note]{Separating functional domains to reduce obscuring of tracking/control follows from the walkthrough findings about confused feature categorisation in 4692551 and governance transparency arguments in 7808872, so this is a reasonable synthesis.} \hldarkgreen{Pilot testing in the walkthrough study revealed inconsistencies in feature categorisation and led to revised definitions and improved evaluator guidelines \parencites[3]{4692551}.}\pdfcomment[icon=Note]{4692551 explicitly reports pilot testing that revealed inconsistencies in feature categorisation and led to revised definitions and improved evaluator guidelines (4692551 Pilot Testing paragraph).} \hllightgreen{Analogously, MindfulFeed must maintain rigorous internal taxonomies for navigation labels and menu groupings, ensuring that features affecting data, attention, or autonomy are not buried under generic headings that hinder user recognition and later expert auditing.}\pdfcomment[icon=Note]{Given the pilot findings in 4692551 and recommendations on maintaining taxonomies for auditability in 4807303/7808872, the recommendation that MindfulFeed maintain internal taxonomies is a supported inference.} \hldarkgreen{Finally, design recommendations from dark pattern research emphasise symmetric optionality, transparent information architecture, honest scarcity, and graceful exit as guiding principles for dark\-pattern\-free interfaces \parencites[6]{4807303}.}\pdfcomment[icon=Note]{4807303 explicitly proposes four design principles, symmetric optionality, transparent information architecture, honest scarcity, and graceful exit, for dark-pattern-free interfaces (4807303 p.6--7).} \hldarkgreen{In navigation terms, this requires that every decision node presents all materially relevant information on a single screen before commitment, that no temporal or scarcity indicators are shown unless grounded in actual system constraints, and that "back," "pause," and "leave" actions are consistently available across hierarchical levels.}\pdfcomment[icon=Note]{4807303's design principles require presenting all relevant information before commitment, accurate scarcity indicators, and graceful exit/back actions; these navigation requirements are directly derived from that source (4807303 p.6--8).} \hllightgreen{Trustworthiness criteria from governance research, including data triangulation, detailed contextual description, and transparent coding of design artefacts, should be applied to information\-architecture documentation so that future audits can verify that MindfulFeed's navigation structure remains aligned with ethical and regulatory expectations rather than drifting toward concealed complexity \parencites[8]{7808872}[6]{4807303}.}\pdfcomment[icon=Note]{7808872 outlines trustworthiness criteria (data triangulation, contextual description, transparent coding using NVivo) and 4807303 recommends documentation for audits, so applying these to IA documentation is a supported synthesis (7808872 p.8; 4807303 p.8).} \hllightgreen{Empirical survey work of platform user behaviour may also temper assumptions about typical usage intensity: many users report modest daily engagement and weak correlations between time spent on platforms and negative outcomes, which suggests that IA rules should not presume constant heavy use but should still prepare for edge cases where brief flows can have outsized effects \parencites[11204]{4153678}.}\pdfcomment[icon=Note]{The survey work in 4153678 shows many users (students) report modest daily engagement and only weak correlations between time spent and negative outcomes, supporting the caution against assuming constant heavy use (4153678 p.4, p.7).} \hllightgreen{Large\-scale analyses of engagement datasets and AI assessments indicate that engagement distributions are often skewed and that multimodal measures (text, survey, biometric proxies) provide a more nuanced picture of typical user behaviour; this nuance appears to encourage IA rules that prioritise low\-friction exits and clear signposting for the plurality of light users \parencites[2]{2882594}.}\pdfcomment[icon=Note]{Sources on large-scale/multimodal analysis (2882594 describes multimodal measures and skewed engagement distributions) together with 4807303's emphasis on low-friction exits support the synthesis that IA should prioritise low-friction exits for light users (2882594 p.4--7; 4807303 p.6--8).} \hllightgreen{Such findings imply the value of lightweight, low\-friction exit options that respect the experience of the majority while preventing disproportionate friction for those who do wish to leave or reduce engagement \parencites[11204]{4153678}.}\pdfcomment[icon=Note]{This is a reasoned implication from the empirical and design findings in 4807303 and user-survey evidence (4807303 p.6--8; 4153678 p.4, p.7) advocating lightweight exit options to balance majority user experience and edge-case protections.}

\subsection{Ethical Recommendation Engine Design}
\label{sec:ethical_rec_engine_design} \hldarkgreen{Ethical configuration of MindfulFeed's recommendation engine must operationalize privacy, fairness, transparency, and human oversight requirements that have been articulated for AI monitoring and personalized systems.}\pdfcomment[icon=Note]{Multiple sources articulate that ethical AI configurations must operationalize privacy, fairness, transparency, and human oversight (see 9300898, p.7; 8760542, p.18).} \hldarkgreen{In occupational health contexts, computational analytics over sensor data and communication logs expose intimate emotional and cognitive states, creating risks of misuse, unauthorized access, and discrimination if safeguards are weak.}\pdfcomment[icon=Note]{Source 8760542 (p.18) explicitly states that sensor and communication data can reveal intimate emotional and cognitive states and pose risks of misuse, unauthorized access, and discrimination.} \hllightgreen{Translating this into feed ranking and content recommendation, engine inputs should be restricted to information strictly necessary for well\-being support, with data minimization rules encoded at the data ingestion and feature engineering layers.}\pdfcomment[icon=Note]{Data minimization and restricting inputs to necessary information are emphasized in the ethics literature (8760542, p.18), so encoding minimization at ingestion/feature layers is a reasonable design inference (9300898, p.7).} \hllightgreen{Configuration matrices in the product record need to list each behavioural signal, its purpose, and retention horizon, ensuring that no auxiliary data are collected solely for engagement maximization.}\pdfcomment[icon=Note]{Documentation of collected signals, purposes, and retention aligns with recommendations on transparency and data governance (8760542, p.18; 7808872, p.8), and the caution against collecting data solely for engagement follows the well‑being orientation in the sources.} \hldarkgreen{Consent management around recommendation logic requires a triad of informed consent, data minimization, and transparent communication between users, institutional partners, and system designers \parencites[18]{8760542}.}\pdfcomment[icon=Note]{Source 8760542 (p.18) explicitly describes a triad of informed consent, data minimization, and transparent communication as central to privacy governance in monitoring systems.} \hllightgreen{Practically, this means that any personalization toggle or profile adjustment within MindfulFeed's engine dashboard must be accompanied by concise explanations detailing what data are used, how ranking scores are computed in broad terms, and how outputs relate to burnout or well\-being goals.}\pdfcomment[icon=Note]{The need for concise, user-facing explanations of data use and algorithmic functioning is supported by calls for explainability and user agency in the literature (9300898, p.7; 8760542, p.18), including links to well‑being outcomes.} \hllightgreen{Consent needs to be iterative and revocable; engine design therefore has to support real\-time deactivation of personalization without degrading core service quality, with logging of consent state changes as auditable events \parencites[18]{8760542}[7]{9300898}.}\pdfcomment[icon=Note]{Sources state consent should be iterative and revocable (8760542, p.18) and emphasize auditability and logging (9300898, p.7), though the specific claim about real‑time deactivation without service degradation is an implementation expectation rather than directly documented.} \hldarkgreen{Bias mitigation constitutes a second pillar for ethical recommendation.}\pdfcomment[icon=Note]{Algorithmic bias as a central ethical pillar is clearly emphasized in the sources (9300898, p.7; 8760542, p.18).} \hldarkgreen{Computational monitoring in healthcare can misclassify stress or overlook vulnerability in underrepresented groups when models are trained on skewed or institution\-specific data, leading to unfair interventions and unequal support allocation \parencites[18]{8760542}.}\pdfcomment[icon=Note]{Source 8760542 (p.18) directly warns that models trained on skewed or institution‑specific data can misclassify stress and miss vulnerabilities in underrepresented groups, causing unfair interventions.} \hllightgreen{Within MindfulFeed, training datasets for recommendation models must be documented with demographic coverage tables and fairness evaluation results, and periodic audits should assess differential accuracy and exposure across gender, age, and other salient groups.}\pdfcomment[icon=Note]{Recommendations for dataset documentation, demographic coverage, fairness evaluation, and periodic audits are advocated across the literature (5375614, p.11; 8760542, p.18; 7808872, p.8), making this an appropriate synthesis.} \hllightgreen{Accessibility safeguards are also necessary; interface paths for configuring recommendations and viewing "why this content" explanations must be usable regardless of technological literacy or cultural background, aligning with equity and inclusion as moral obligations rather than purely technical targets \parencites[18]{8760542}[0180546]{0180546}.}\pdfcomment[icon=Note]{Sources stress accessibility and equity in deployment so interfaces and explanations must be usable across literacy and cultural backgrounds (8760542, p.18; 0180546, p.4), supporting the moral framing.} \hldarkgreen{Responsible personalization frameworks in education further refine how recommendation engines should relate to human agency.}\pdfcomment[icon=Note]{Responsible personalization frameworks in education and their implications for human agency are discussed in the literature (5375614, p.6), directly supporting this claim.} \hldarkgreen{Personalization can adjust pace, sequence, difficulty, feedback type, and representation, but algorithms must not determine paths in a closed\-off way; learners need to understand recommendations, have choices, and be able to reject or adjust their profiles \parencites[6]{5375614}.}\pdfcomment[icon=Note]{Source 5375614 (p.6) explicitly states personalization can adjust pace, sequence, difficulty, feedback and that algorithms should not deterministically lock learners out of choice.} \hllightgreen{Applied to MindfulFeed, ranking functions must be exposed through an "algorithm insights" panel where profile dimensions and content categories are intelligible to ordinary users.}\pdfcomment[icon=Note]{The call to expose ranking functions and make profile dimensions intelligible aligns with recommendations for explainability and user agency in the sources (9300898, p.7; 5375614, p.6), though the specific 'algorithm insights' panel is an applied design suggestion.} \hldarkgreen{Human\-in\-the\-loop models, in which systems generate suggestions while humans interpret context and decide next steps, imply that high\-risk outputs, such as strong behavioural nudges or mental health advice, should require explicit confirmation or be reviewable by qualified moderators before widespread deployment \parencites[11]{5375614}[5375614]{5375614}.}\pdfcomment[icon=Note]{Human‑in‑the‑loop approaches and the requirement that high‑risk outputs be reviewable or require confirmation are prescribed in the education/personalisaton literature (5375614, p.6; p.11).} \hldarkgreen{Governance analyses highlight accountability, transparency, oversight, and well\-being outcomes as central dimensions for AI systems that affect cognition and emotion \parencites[8]{7808872}.}\pdfcomment[icon=Note]{The governance dimensions of accountability, transparency, oversight, and well‑being outcomes are explicitly identified as central in 7808872 (p.8).} \hllightgreen{MindfulFeed's recommendation engine therefore needs architectural features that map directly onto these dimensions: audit logs and impact dashboards to support accountability; per\-item and per\-topic explanation tooltips to achieve transparency; feedback and correction channels for oversight; and monitoring of psychological and trust metrics to evaluate well\-being outcomes, rather than relying solely on engagement indicators \parencites[7]{9300898}[8]{7808872}.}\pdfcomment[icon=Note]{Mapping architectural features (audit logs, explanations, feedback channels, well‑being metrics) to governance dimensions is supported by 7808872 (p.8) and related governance recommendations (9300898, p.7), though specific implementation details are inferential.} \hllightgreen{To maintain trustworthiness, engine design documentation should mirror qualitative governance procedures, with transparent coding of algorithm changes and triangulation of behavioural data, user feedback, and policy constraints when adjusting ranking strategies \parencites[8]{7808872}.}\pdfcomment[icon=Note]{The proposal that documentation should mirror governance procedures and triangulate data, feedback, and policy aligns with methodological recommendations in 7808872 (p.8) and transparency calls in 8760542 (p.18).} \hldarkgreen{Studies on digital health and preventive well\-being interventions show that AI models can analyse behavioural patterns, sleep, physical activity, and communication to detect risks early and support stress reduction, chronic disease management, and mental health promotion.}\pdfcomment[icon=Note]{Source 9300898 (p.7) specifically reports that AI models can analyze behaviour, sleep, activity, and communication to detect risks early and support preventive well‑being interventions.} \hldarkgreen{However, extensive monitoring without robust privacy and consent controls undermines autonomy and can be perceived as surveillance \parencites[7]{9300898}.}\pdfcomment[icon=Note]{Both 9300898 (p.7) and 8760542 (p.18) warn that extensive monitoring without robust privacy/consent undermines autonomy and can be perceived as surveillance.} \hllightgreen{For MindfulFeed, this implies that recommendation objectives must be framed explicitly as preventive well\-being support, with clear boundaries that exclude commercial targeting or opaque behavioural conditioning.}\pdfcomment[icon=Note]{Sources advocate framing AI objectives around well‑being and prioritizing user welfare over engagement or commercial targeting (9300898, p.7; 7808872, p.8; 8760542, p.18), so this is a supported design implication.} \hllightgreen{Engine configuration grids should specify which model outputs trigger supportive prompts versus when no action is taken, and these rules must be open to user review and adjustment to avoid covert manipulation \parencites[7]{9300898}[5375614]{5375614}.}\pdfcomment[icon=Note]{The need for transparent, reviewable rules about when models trigger interventions to avoid covert manipulation is consistent with human oversight and openness recommendations (5375614, p.11; 9300898, p.7).} \hldarkgreen{Finally, organizational governance mechanisms such as AI ethics boards and compliance committees are increasingly adopted to oversee data handling, model validation, and risk management in healthcare monitoring \parencites[18]{8760542}.}\pdfcomment[icon=Note]{Source 8760542 (p.18) documents the increasing adoption of internal governance structures like AI ethics boards and compliance committees in healthcare monitoring contexts.} \hllightgreen{MindfulFeed's development record should similarly describe internal oversight structures for the recommendation engine, including periodic fairness and privacy audits, review gates before major model updates, and procedures to communicate model limitations to users to prevent overreliance.}\pdfcomment[icon=Note]{Recommendations that development records describe oversight structures, audits, review gates, and communication of limitations are consistent with governance guidance in 8760542 (p.18) and 5375614 (p.11).} \hllightgreen{Embedding these controls into the engine design aligns computational personalization with humane, well\-being\-oriented AI practice and supports cognitive\-health safeguards central to the product brief \parencites[18]{8760542}[7]{9300898}[8]{7808872}.}\pdfcomment[icon=Note]{The claim that embedding these controls aligns personalization with humane, well‑being‑oriented AI and supports cognitive‑health safeguards synthesizes themes from 9300898 (p.7) and 7808872 (p.8).}

\subsection{Monetization Strategy Without Dark Patterns}
\label{sec:monetization_without_dark_patterns} \hllightgreen{Monetization options for MindfulFeed must be constrained by empirical evidence on dark patterns and ethical AI governance so that revenue generation does not depend on deceptive interface strategies or exploitative personalization.}\pdfcomment[icon=Note]{Supported by empirical dark-pattern research and ethical AI governance recommendations in 4807303 (pp.2,8) and governance/ethical AI discussion in 8760542 (p.118), so this is a reasonable synthesis.} \hldarkgreen{Dark pattern taxonomies grounded in mechanistic clarity, non overlap, and exhaustiveness describe $27$ patterns, including sneaking, urgency or scarcity cues, obstruction, misdirection, forced action, confirm shaming, roach motel, bait and switch, and drip pricing, all designed to maximise organisational gain at the expense of consumer autonomy and informed consent \parencites[2]{5127387}[2]{4807303}.}\pdfcomment[icon=Note]{Directly supported by 4807303 which describes a 27-pattern taxonomy and lists examples and design principles (p.2, p.3).} \hldarkgreen{These classifications, together with findings that aggressive dark patterns can increase subscription uptake by over $200\%$, demonstrate that many commercially successful monetization flows rely on cognitive bias exploitation such as default bias, scarcity bias, framing, and loss aversion rather than transparent value exchange \parencites[2]{4807303}[2]{5127387}.}\pdfcomment[icon=Note]{4807303 cites Luguri \& Strahilevitz's RCT showing aggressive dark patterns increased subscription uptake by over 200\% and also links patterns to cognitive biases (p.2, Related Work).} \hllightgreen{For MindfulFeed, such mechanisms must be formally excluded from pricing, subscription, and upgrade journeys, and this exclusion documented within design review checklists and user journey flow artefacts.}\pdfcomment[icon=Note]{4807303 recommends excluding manipulative mechanisms from key journeys and maintaining documentation (pp.7--8), so this normative prescription is a supported inference.} \hllightgreen{Symmetric design principles derived from dark pattern research provide concrete execution rules for non manipulative monetization flows.}\pdfcomment[icon=Note]{4807303 recommends symmetric design principles as concrete rules for non-manipulative flows (p.8), so this is a valid synthesis of that recommendation.} \hldarkgreen{Evidence shows that complex cancellation paths with eight steps, forced traversal of promotional detours, mandatory survey completion, and manual text entry produce average cancellation times of $18$ minutes $42$ seconds and abandonment rates of $31.4\%$, while sign up in the same context requires only three steps and $90$ seconds.}\pdfcomment[icon=Note]{The specific cancellation/sign-up metrics (18m42s, 31.4\% abandonment; sign-up three steps \textasciitilde{}90s) are reported in 4807303 (p.8).} \hldarkgreen{Design recommendations therefore require that opting out, cancelling, or exercising privacy preferences be no more complex than their inverses, that calls to action use plain language without guilt or loss framing, and that organisations adopt policies defining dark patterns as prohibited practices, with accountability mechanisms empowered to intervene.}\pdfcomment[icon=Note]{4807303 explicitly recommends parity in complexity for opting out/cancelling, plain-language CTAs, and organisational policies prohibiting dark patterns (p.8--10).} \hllightgreen{Applied to monetization, this implies that any paid tier activation must be matched by an equally simple and visible downgrade or cancellation route, with user journey documentation for checkout and cancellation retained for audit and attached to the product record as part of ethical review artefacts \parencites[8]{4807303}.}\pdfcomment[icon=Note]{4807303's symmetric-design and documentation recommendations (p.8, p.10) support applying these principles to monetization and retaining user-journey documentation for audits.} \hldarkgreen{Automated auditing capacity via tools such as DarkScan supports technical enforcement of these monetization constraints.}\pdfcomment[icon=Note]{4807303 describes DarkScan and recommends automated auditing as part of enforcement and continuous review (pp.2,8--10).} \hldarkgreen{DarkScan combines DOM analysis, text classification, visual detection, and dynamic behavioural flow tracing to identify manipulative patterns and quantify their behavioural and psychological impact \parencites[2]{4807303}.}\pdfcomment[icon=Note]{4807303 describes DarkScan as combining DOM analysis, text classification, visual detection, and dynamic behavioural flow tracing in a unified pipeline (p.2, Datasets/Methods).} \hldarkgreen{Continuous integration pipelines for MindfulFeed should therefore include automated scans of all payment, upgrade, and consent screens, with severity scoring via models such as DPSS used to flag flows that unintentionally introduce obstruction or roach motel characteristics.}\pdfcomment[icon=Note]{4807303 recommends integrating automated dark-pattern scans into CI pipelines and using DPSS severity scoring to flag obstructive/roach-motel flows (pp.8--10).} \hldarkgreen{Versioned flow diagrams for monetization pathways, as recommended for high risk sectors, can be stored alongside audit results to provide regulators and internal ethics bodies with traceable evidence of dark pattern free revenue mechanisms \parencites[8]{4807303}.}\pdfcomment[icon=Note]{4807303 recommends maintaining full user-journey flow documentation with version history and storing alongside audit results for regulators and ethics bodies (p.10, Recommendations).} \hllightgreen{Privacy and informed consent requirements in computational monitoring environments further constrain data based monetization.}\pdfcomment[icon=Note]{8760542 discusses privacy and informed consent constraints in computational monitoring (p.118), which reasonably constrains data-driven monetization practices.} \hldarkgreen{AI driven analytics over behavioural traces, communication logs, and psychological states risk misuse and discrimination without stringent safeguards, requiring transparent consent protocols that clearly explain what data are collected, how they are processed, and for what purposes, with consent iterative and revocable and data minimized to legitimate objectives \parencites[18]{8760542}.}\pdfcomment[icon=Note]{8760542 explicitly warns AI analytics over behavioural and psychological data risk misuse and discrimination and calls for transparent, iterative, revocable consent and data minimization (p.118).} \hllightgreen{If any paid features in MindfulFeed depend on enhanced analytics or personalised wellbeing recommendations, associated consent dialogues must adhere to this triad, avoiding complex privacy policies, confusing consent mechanisms, and defaults that encourage unintended disclosure, all of which are common in manipulative commercial designs \parencites[18]{8760542}[2]{2668071}.}\pdfcomment[icon=Note]{2668071 and related dark-pattern literature note that complex policies, confusing consent mechanisms, and default settings are common manipulative practices and should be avoided (2668071 p.3; 5127387 overview).} \hldarkgreen{Algorithmic bias and fairness concerns in healthcare oriented AI monitoring show that monetizing advanced analytics without diverse data representation and fairness aware design can reproduce inequities across gender, race, age, or role, misclassifying stress levels or failing to detect vulnerability in underrepresented groups \parencites[18]{8760542}.}\pdfcomment[icon=Note]{8760542 discusses algorithmic bias in healthcare monitoring and how lack of diverse representation and fairness-aware design can reproduce inequities and misclassify underrepresented groups (p.118).} \hllightgreen{Consequently, MindfulFeed's monetization strategy cannot tie premium access to more accurate or fair recommendations for only paying users; fairness and accessibility must be treated as baseline obligations, with paid tiers restricted to non discriminatory enhancements such as additional reflection tools, provided that these are presented plainly and without scarcity cues or confirm shaming \parencites[18]{8760542}[5127387]{5127387}.}\pdfcomment[icon=Note]{The normative claim is grounded in evidence about fairness obligations in 8760542 (p.118) and dark-pattern harms discussed in 5127387, so restricting paid tiers to non-discriminatory enhancements is a supported inference.} \hllightgreen{Finally, regulatory expectations around consent banners under GDPR and the ePrivacy Directive, including requirements that consent be freely given, specific, informed, unambiguous, and revocable, imply that revenue related flows cannot condition core access on intrusive tracking or opaque personalization \parencites[4]{5678011}.}\pdfcomment[icon=Note]{5678011 summarizes GDPR/ePrivacy consent requirements (freely given, specific, informed, unambiguous, revocable) and implies conditioning access on intrusive tracking conflicts with these requirements (p.3--4).} \hllightgreen{Incorporating legal scholarship into design reviews means that every monetization related interface is assessed for preselected options, hidden opt outs, and asymmetric visual prominence, with redesign mandated wherever such features appear.}\pdfcomment[icon=Note]{5678011 and 4807303 recommend legal-review-informed assessments of interfaces for preselected options, hidden opt-outs, and asymmetries, with remediation where detected (5678011 p.3; 4807303 p.10).} \hllightgreen{In combination, dark pattern taxonomies, automated detection, privacy triads, and fairness obligations provide the evidence based constraints within which MindfulFeed's monetization remains compatible with pro social, cognitive health oriented objectives rather than undermining them for short term conversion gains \parencites[2]{4807303}[2]{5127387}[4]{5678011}[18]{8760542}.}\pdfcomment[icon=Note]{This is a synthesis of evidence-based constraints drawn from 4807303 (taxonomy, DarkScan, DPSS), 5678011 (legal consent requirements) and 8760542 (privacy/fairness obligations).} \hldarkgreen{The broader literature on the digital welfare state also cautions that large scale digitisation creates new vectors for exclusion, systemic error, and data driven control that may exacerbate inequality unless governance and oversight are explicitly designed to prevent such outcomes.}\pdfcomment[icon=Note]{2832310 and 0967765 warn that large-scale digitalisation can create exclusion, systemic error, and data-driven control that exacerbate inequality absent proper governance (2832310 p.2, p.3; 0967765 p.2--3).} \hllightgreen{These concerns about automated decision making and discrimination in public systems may suggest that platform monetization strategies should be treated with similar legal and ethical scrutiny \parencites[2]{2832310}.}\pdfcomment[icon=Note]{Given the documented risks in public systems (2832310, 0967765) and the regulatory/ethical concerns about dark patterns (5678011, 4807303), treating platform monetization with similar scrutiny is a reasonable inference.}

\section{Evaluation of MindfulFeed Across Desirability, Feasibility, and Profitability}

\subsection{User Desirability Assessment}
\label{sec:user_desirability_assessment} \hllightgreen{User desirability for MindfulFeed depends on whether its interaction model, information architecture, and AI personalization features align with documented user expectations for autonomy, transparency, and cognitive relief rather than engagement maximization.}\pdfcomment[icon=Note]{Governance and HCI sources discuss autonomy, transparency, reducing distraction and personalization trade-offs (7808872 p.3, 5375614 p.6, 7173856 p.2), so this is a reasonable synthesis linking design features to user desirability.} \hldarkgreen{Cognitive walkthroughs on major social media platforms show that dark patterns often appear in combinations across content engagement, account creation, and termination, generating distraction during learning and exposure to harmful content that undermines educational potential and lowers achievement.}\pdfcomment[icon=Note]{A cognitive-walkthrough study of major platforms and review literature explicitly reports combinations of dark patterns across content engagement and account flows and links them to distraction and lowered educational outcomes (4692551 p.3; 4692551 discussion of Himawan and Rosy).} \hllightgreen{These harms suggest that desirability for a pro\-social app is strongly conditioned on the absence of such combinations, particularly in navigation flows tied to attention and account control.}\pdfcomment[icon=Note]{This is a reasonable inference from the harms documented in the walkthrough and governance literature: that desirability depends on absence of such harmful combinations (4692551 p.3; 7808872 p.3--8).} \hllightgreen{A development record that explicitly maps and removes these pattern clusters, and documents low cognitive load in walkthroughs, directly addresses user demand for less manipulative engagement \parencites[3]{4692551}.}\pdfcomment[icon=Note]{4692551 describes mapping and analysing dark-pattern combinations and governance sources discuss documenting user-facing effects, so the claim that such a record addresses demand for less manipulative engagement is a defensible inference (4692551 p.3; 7808872 p.8).} \hldarkgreen{Methodologically robust evidence on dark patterns demonstrates measurable impacts on decision quality, platform trust, and psychological wellbeing, with aggressive consent designs reducing privacy\-preserving choices by up to $23\%$ and complex cookie banners architecturally determining withdrawal behaviour.}\pdfcomment[icon=Note]{Empirical dark-pattern research reports measurable impacts on decision quality, trust, and wellbeing and cites aggressive consent designs reducing privacy-preserving choices by up to 23 percentage points and cookie-banner effects on withdrawal (4807303 p.2; 3667034/4807303 summary).} \hllightgreen{For desirability assessment, this implies that users are likely to value MindfulFeed only if consent and configuration flows visibly diverge from such architectures.}\pdfcomment[icon=Note]{This follows logically from the evidence on consent-architecture harms: users are likely to prefer systems that diverge from manipulative consent architectures (4807303 p.2; 5678011 p.5).} \hllightgreen{Testing protocols must therefore record whether participants can identify and enact privacy\-preserving choices without confusion or excessive effort, and whether they report higher trust when symmetric opt\-in and opt\-out pathways are implemented.}\pdfcomment[icon=Note]{Consent interaction research and measurement literature recommend recording users' ability to choose privacy and linking design symmetry to trust (5678011 p.5; 6103768 p.6), so this testing recommendation is a reasonable synthesis.} \hllightgreen{Positive responses on these dimensions constitute direct evidence that design changes enhance perceived desirability relative to mainstream platforms \parencites[2]{4807303}.}\pdfcomment[icon=Note]{Given the measurement frameworks tying trust/usefulness to behavioural intention (6103768 p.6), positive responses on these dimensions are a reasonable indicator that design changes improve perceived desirability relative to mainstream platforms.} \hldarkgreen{Governance research on digital well\-being introduces four analytical dimensions that link administrative design with experiential outcomes: accountability, transparency, oversight, and well\-being outcomes.}\pdfcomment[icon=Note]{The four analytical dimensions (accountability, transparency, oversight, well-being outcomes) are explicitly listed in the governance analytical framework (7808872 p.8).} \hllightgreen{Desirability assessment for MindfulFeed can be structured around these dimensions.}\pdfcomment[icon=Note]{Applying the governance dimensions to a desirability assessment is a logical structuring suggested by the analytical framework (7808872 p.8).} \hldarkgreen{Interview and questionnaire instruments should capture whether users can locate accountability cues (for example, explanations of who is responsible for algorithmic decisions), understand how AI systems affecting cognition and emotion operate, identify mechanisms for feedback and correction, and describe emotional stability, cognitive autonomy, and social connection when using the app.}\pdfcomment[icon=Note]{7808872 explicitly recommends measuring understanding of accountability, explainability, feedback mechanisms and well-being outcomes; educational/AI literature also calls for measuring AI explanations and emotional/cognitive effects (7808872 p.8; 5375614 p.6).} \hldarkgreen{Credibility, transferability, and dependability in such qualitative evaluation depend on triangulating user narratives, expert walkthrough findings, and interaction logs, mirroring NVivo\-supported thematic content analysis procedures used in governance studies \parencites[7808872]{7808872}.}\pdfcomment[icon=Note]{The governance study describes prioritizing credibility, transferability, and dependability via triangulation of narratives, expert analysis, and logs and notes using NVivo for thematic analysis (7808872 3.5, 3.3--3.4).} \hllightgreen{High scores across these four dimensions, combined with robust qualitative coding, indicate that the product's value proposition aligns with expressed user needs for ethically governed digital experiences.}\pdfcomment[icon=Note]{Inferring that high scores across governance dimensions combined with robust coding indicate alignment with user needs is a reasonable synthesis of 7808872's validity/trustworthiness arguments (7808872 3.5).} \hldarkgreen{Mindfulness app research employing heuristic evaluation, survey, diary, and interviews over nine weeks reveals that commercialization and profit motives can lead to manipulative design techniques that contradict the foundations of mindfulness.}\pdfcomment[icon=Note]{A nine-week multi-method mindfulness app study using heuristic eval, surveys, diaries and interviews documents commercialization leading to manipulative design that contradicts mindfulness foundations (7173856 p.4; 7173856 methodology).} \hldarkgreen{Brignull's categories such as bait and switch, confirm shaming, disguised ads, forced continuity, friend spam, and hidden cost illustrate how apps can pressure users into unwanted subscriptions, guilt\-driven confirmations, or hidden charges.}\pdfcomment[icon=Note]{Brignull's dark-pattern categories (bait \& switch, confirm shaming, disguised ads, forced continuity, friend spam, hidden cost) are listed and discussed in the HCI/dark-pattern literature and reproduced in the study (7173856 Table 1, p.2).} \hldarkgreen{Findings from this multi\-method study show that such dark patterns negatively shape user perceptions and engagement, including discontinuation and "fleeing" behaviour.}\pdfcomment[icon=Note]{The multi-method study reports that dark patterns negatively affect perceptions and engagement, including discontinuation and fleeing behaviour (7173856 p.4 and study results).} \hldarkgreen{For MindfulFeed, desirability assessment must therefore include heuristic reviews of monetization and habit\-forming features, diaries capturing reasons for non\-usage, and interviews probing reactions to any persuasive element.}\pdfcomment[icon=Note]{The recommended methods (heuristic reviews, diaries, interviews) mirror the multi-method design used to assess mindfulness apps and dark patterns in 7173856 (7173856 2.1--2.3).} \hllightgreen{Reduced incidence of guilt, confusion, or perceived exploitation compared with the documented mindfulness app patterns would provide strong evidence that MindfulFeed's ethical design increases desirability among users seeking mental health support tools \parencites[2]{7173856}.}\pdfcomment[icon=Note]{Given the study documented guilt/confusing manipulative patterns (7173856 p.2--4), it is a reasonable inference that reduced incidence of these would indicate increased desirability for mental-health-seeking users.} \hldarkgreen{Systematic review work on AI\-based burnout interventions demonstrates the value of structured coding frameworks, MMAT quality assessment, and narrative synthesis for understanding how computational methods affect stress and burnout outcomes.}\pdfcomment[icon=Note]{A systematic review on AI-based burnout interventions describes using structured coding frameworks, MMAT quality assessment, and narrative synthesis (8760542 p.119 Methods section).} \hllightgreen{Translating these procedures into desirability evaluation, MindfulFeed's testing data should be coded for variables such as perceived cognitive load, stress change during use, clarity of AI explanations, and perceived usefulness of personalization controls.}\pdfcomment[icon=Note]{Translating systematic-review procedures into coded desirability variables (cognitive load, stress change, clarity, usefulness) is a reasonable application supported by 8760542's coding approach and 6103768's measurement constructs (8760542 p.119; 6103768 p.6).} \hldarkgreen{Inter\-rater reliability measures like Cohen's kappa can be used to ensure consistent coding of desirability indicators across evaluators \parencites[19]{8760542}.}\pdfcomment[icon=Note]{The review reports using Cohen's kappa to assess inter-rater reliability (8760542 p.119), supporting its use for consistent coding of desirability indicators.} \hldarkgreen{Evidence indicates that users appreciate AI systems that enhance information acquisition efficiency and satisfy personal needs when recommendation accuracy, transparency, and reliability are high, and when trust and perceived usefulness are correspondingly strong \parencites[6]{6103768}.}\pdfcomment[icon=Note]{6103768 constructs and evidence link recommendation accuracy, transparency, and reliability to trust and perceived usefulness, which improve information acquisition and satisfaction (6103768 p.6).} \hllightgreen{Questionnaire measures adapted from these constructs can quantify whether MindfulFeed's recommendation features increase behavioural intention to accept and continuously use content suggested in a health\-support context, without undermining autonomy.}\pdfcomment[icon=Note]{6103768 defines behavioral intention measures in public health contexts and recommends adapting questionnaire items for accuracy/transparency/trust; applying these to MindfulFeed is a logical extension (6103768 p.6).} \hllightgreen{Taken together, desirability assessment for MindfulFeed must combine task\-based expert walkthroughs, multi\-method user studies, and structured coding of trust, autonomy, and cognitive strain.}\pdfcomment[icon=Note]{Combining expert walkthroughs, multi-method user studies, and structured coding is supported by the cited methodological examples (4692551 cognitive walkthroughs; 7173856 multi-method; 8760542 coding), so this is a reasonable synthesis.} \hllightgreen{Only if empirical indicators show reduced dark pattern exposure, enhanced governance visibility, and sustained perceived usefulness of transparent AI recommendations can the product be considered desirable in light of current evidence on social media harms and ethical digital well\-being \parencites[2]{4807303}[3]{4692551}[7808872]{7808872}[6]{6103768}.}\pdfcomment[icon=Note]{The conditional conclusion aligns with evidence that reduced dark-pattern exposure, stronger governance visibility, and perceived usefulness of transparent recommendations are necessary for desirability, as discussed across the cited works (4807303 p.2; 7808872 p.3--8; 6103768 p.6).}

\subsection{Technical and Organizational Feasibility}
\label{sec:tech_org_feasibility} \hllightgreen{Technical and organizational feasibility for MindfulFeed hinges on whether the app's governance, data, and AI practices can be implemented within existing institutional capabilities and regulatory expectations.}\pdfcomment[icon=Note]{This is a reasonable synthesis supported by governance and feasibility discussions in the provided literature (e.g., governance and regulatory needs in occupational-health and AI policy discussions: 8760542, p.18; 9300898, p.7).} \hldarkgreen{Evidence from occupational health monitoring shows that computational systems routinely handle highly sensitive behavioural and psychological data, yet responsible deployment is achievable when privacy governance follows a triad of informed consent, data minimization, and transparent communication between institutions, employees, and system designers.}\pdfcomment[icon=Note]{Source 8760542 (p.18) explicitly states that privacy governance in AI-driven monitoring requires informed consent, data minimization, and transparent communication, and discusses handling sensitive behavioural/psychological data in occupational health contexts.} \hldarkgreen{Translating this to a social app context implies that feasible implementation requires technical infrastructure to manage iterative and revocable consent states, strict data minimization at collection and storage layers, and auditable communication channels explaining how recommendation and analytics components operate.}\pdfcomment[icon=Note]{The need for iterative/revocable consent, data minimization at collection/storage, and auditable communication channels is explicitly recommended for AI-driven health monitoring in 8760542 (p.18) and echoed in governance-focused sources (9300898, p.7).} \hldarkgreen{Embedding algorithmic transparency and auditability into MindfulFeed's architecture is also feasible given current AI practices.}\pdfcomment[icon=Note]{Sources discussing current AI practices and XAI indicate that embedding transparency and auditability is achievable (9945186, p.4; 9300898, p.7), supporting feasibility of such architecture.} \hldarkgreen{In healthcare monitoring, ethical integrity demands that predictive systems be explainable and auditable, allowing users to understand how conclusions about mental health are reached and preventing AI from being perceived as an intrusive management tool \parencites[18]{8760542}.}\pdfcomment[icon=Note]{8760542 (p.18) directly argues that predictive systems in healthcare must be explainable and auditable to preserve ethical integrity and avoid perceptions of intrusive management.} \hllightgreen{Practically, this necessitates logging of model inputs and outputs, traceable configuration of recommendation parameters, and interfaces capable of surfacing "why this content" explanations.}\pdfcomment[icon=Note]{Logging inputs/outputs, traceable recommendation configurations, and 'why this content' explanations follow from recommendations for provenance, logging, and explainability in the sources (9945186, p.4; 8760542, p.18), a reasonable practical specification.} \hldarkgreen{Similar demands appear in explainable AI and governance work in quantitative finance, where systems are expected to trade off some opacity for pipelines that are explainable and enhanced with domain knowledge so that stakeholders can understand, audit, and maintain confidence in AI strategies \parencites[4]{9945186}.}\pdfcomment[icon=Note]{Source 9945186 (p.4) explicitly describes quantitative finance moving toward pipelines that trade some opacity for explainability and domain-knowledge integration so stakeholders can understand, audit, and trust AI strategies.} \hllightgreen{MindfulFeed's technical feasibility therefore depends on adopting explainability modules, governance hooks, and data provenance verification routines that are already recognised as achievable in other AI domains.}\pdfcomment[icon=Note]{This is a valid synthesis: multiple sources describe explainability modules, governance hooks, and provenance routines as achievable in other AI domains (9945186, p.4; 9300898, p.7; 8760542, p.18), supporting applicability to MindfulFeed.} \hldarkgreen{Organizational feasibility is shaped by the need for continuous ethical governance rather than one\-off compliance.}\pdfcomment[icon=Note]{The literature emphasizes continuous ethical governance and monitoring rather than one-off compliance as essential for responsible AI (8760542, p.19; 9300898, p.7), directly supporting this claim.} \hldarkgreen{In AI\-enabled occupational health, ethical governance is described as a precondition for responsible and humane application of AI, requiring organizational cultures that prioritise compassion, transparency, and trust over control and surveillance, alongside internal structures such as AI ethics boards and compliance committees to oversee data handling, model validation, and risk management \parencites[18]{8760542}.}\pdfcomment[icon=Note]{8760542 (p.18) explicitly frames ethical governance as a precondition for humane AI application in occupational health and recommends organizational cultures prioritising compassion, transparency, trust, and internal structures like ethics boards and compliance committees.} \hllightgreen{For MindfulFeed, this implies that feasible deployment depends on instituting comparable oversight bodies, integrating ethical review of dark patterns and consent flows into design governance, and aligning internal policies with privacy, fairness, and accountability requirements.}\pdfcomment[icon=Note]{This is a reasonable inference from sources that call for oversight bodies and ethical review in AI deployments in health and workplace contexts (8760542, p.18; 9300898, p.7).} \hllightgreen{Systematic review methodology in computational burnout research demonstrates that organizations can manage complex, technology\-heavy projects with structured processes.}\pdfcomment[icon=Note]{The systematic review methods described in the sources (e.g., detailed PRISMA-based review in 8760542, p.19 and thematic/structured methods in 7808872, p.8) demonstrate that organizations can manage complex, technology-heavy projects with structured processes.} \hldarkgreen{The review employed PRISMA guidelines, multi\-database searches, staged identification and screening, structured coding frameworks, thematic clustering, and quality assessment via the Mixed Methods Appraisal Tool, with inter\-rater reliability measured by Cohen's kappa \parencites[19]{8760542}.}\pdfcomment[icon=Note]{8760542 (p.19) explicitly reports using PRISMA guidelines, multi-database searches, staged screening, structured coding, thematic clustering, MMAT quality assessment, and reports inter-rater reliability with Cohen's kappa.} \hllightgreen{These procedures show that institutions can sustain methodologically rigorous, data\-intensive workflows when roles, tools, and quality thresholds are clearly defined.}\pdfcomment[icon=Note]{This is a reasonable interpretation: the described procedures (8760542, p.19; 7808872, p.8) support the conclusion that institutions can sustain rigorous, data-intensive workflows when roles, tools, and quality thresholds are defined.} \hllightgreen{MindfulFeed's development process can feasibly mirror these structures in its Empathize, Define, Ideate, Prototype, and Test phases, for example by treating usability logs and dark\-pattern audits as "studies" that pass through scripted identification, coding, and inclusion stages.}\pdfcomment[icon=Note]{Treating development phases and usability logs/dark-pattern audits as structured studies is a plausible adaptation of the systematic-review-style processes described in 8760542 (p.19) and governance/method recommendations across the corpus (7808872, p.8).} \hldarkgreen{Responsible digital education analysis adds further organisational guidance.}\pdfcomment[icon=Note]{Source 5375614 (p.15) explicitly states that responsible digital education analysis provides organisational guidance and principles relevant to governance and deployment.} \hldarkgreen{Institutions are advised to apply pedagogy first, equity by design, privacy by design, and evidence before scale, with data minimization, human oversight, explainability, impact assessment, correction rights, and security implemented before using AI in important decisions \parencites[15]{5375614}.}\pdfcomment[icon=Note]{5375614 (p.15) directly lists 'pedagogy first, equity by design, privacy by design, and evidence before scale' along with data minimization, human oversight, explainability, impact assessment, correction rights, and security before high-stakes AI use.} \hllightgreen{This indicates that MindfulFeed is organizationally feasible only if its governance model adopts similar principles: specifying well\-being and autonomy objectives before technology choices, involving vulnerable groups in design, limiting data from the start, and scaling features only after careful evaluation of results, risks, and costs.}\pdfcomment[icon=Note]{This is a reasonable synthesis: applying the educational governance principles from 5375614 (p.15) to MindfulFeed is a supported inference about organizational feasibility and required safeguards.} \hldarkgreen{Finally, dark\-pattern research provides concrete technical standards and audit requirements that MindfulFeed can realistically adopt.}\pdfcomment[icon=Note]{Dark-pattern research in 4807303 (p.8 and throughout) provides concrete technical recommendations and audit requirements (e.g., DarkScan, DPSS) that could be adopted by applications like MindfulFeed.} \hldarkgreen{Designers are advised to implement symmetric design principles so that opting out, cancelling, or exercising privacy preferences are no more complex than their inverses, to use plain language in calls to action, and to conduct regular third\-party audits using automated tools such as DarkScan within continuous integration pipelines.}\pdfcomment[icon=Note]{4807303 (p.8) explicitly recommends symmetric design principles, plain language in calls to action, and regular third-party audits using automated tools such as DarkScan within CI pipelines.} \hldarkgreen{Organizations are further encouraged to maintain user journey documentation for key pathways and treat user trust, task success, and long\-term retention as success metrics alongside conversion, while regulators move toward mandatory "dark\-pattern\-free" certification in high\-risk sectors \parencites[8]{4807303}.}\pdfcomment[icon=Note]{4807303 (p.8) recommends maintaining user journey documentation, using trust/task success/retention metrics alongside conversion, and proposes 'dark-pattern-free' certification for high-risk sectors, supporting this claim.} \hllightgreen{These recommendations show that technical artifacts such as automated scans, flow documentation, and trust metrics are practically implementable and form a feasible backbone for MindfulFeed's organisational commitment to dark\-pattern avoidance and cognitive\-health protection.}\pdfcomment[icon=Note]{This is a supported synthesis: the recommendations in 4807303 (p.8) (automated scans, flow documentation, trust metrics) demonstrate practical implementability and form a feasible backbone for organizational commitments to avoid dark patterns.}

\subsection{Profitability and Business Sustainability Analysis}
\label{sec:profitability_sustainability_analysis} \hllightgreen{Profitability and long\-term business sustainability for MindfulFeed must be evaluated within constraints set by empirical work on dark patterns, privacy, and ethical AI governance.}\pdfcomment[icon=Note]{Sources discuss dark patterns, privacy, and ethical AI governance as constraints on product design and business models (e.g., 9300898 p.7 on privacy/governance and 2668071 p.3 on dark patterns and ethics), so this is a reasonable synthesis.} \hldarkgreen{Research on manipulative interfaces shows that many revenue\-optimizing designs rely on strategies such as sneaking, urgency or scarcity cues, obstruction, misdirection, forced action, confirm shaming, roach motel flows, bait and switch, and drip pricing, all purposefully crafted to steer users into actions they did not plan to take by exploiting default bias, scarcity bias, framing, and loss aversion \parencites[2]{5127387}.}\pdfcomment[icon=Note]{Multiple sources enumerate these dark pattern strategies and link them to cognitive biases (5127387 p.2 lists tactics and biases; 7173856 p.2--3 lists Brignull categories), directly supporting the claim.} \hldarkgreen{Experimental analysis demonstrates that complex cancellation pathways, consisting of multiple forced detours and confirmation hurdles, generate high abandonment rates and extended task times, while aggressive dark patterns in consent and checkout contexts can dramatically increase subscription uptake compared with neutral designs.}\pdfcomment[icon=Note]{Empirical results about complex cancellation paths, high abandonment and long cancellation times, and elevated subscription uptake with aggressive patterns are documented in the corpus study (4807303 p.8 reports average cancellation time, abandonment rate and experimental findings).} \hllightgreen{For MindfulFeed, these findings imply that many conventional optimization levers for profitability are ethically unavailable; financial viability must instead be grounded in value propositions that do not depend on deceptive friction or emotionally manipulative copy.}\pdfcomment[icon=Note]{This is an inferential recommendation grounded in empirical and governance literature that discourages deceptive optimisation (4807303 p.10 policy recommendations; 9300898 p.7 on human-centered AI), so it is a supported synthesis.} \hldarkgreen{Organizational recommendations derived from dark pattern severity assessment and automated detection pipelines such as DarkScan provide a blueprint for sustainable, non\-exploitative business practices.}\pdfcomment[icon=Note]{The recommendation linking severity assessment and automated detection pipelines (DarkScan) to organisational practice is explicitly made in 4807303 (p.8--11 describing DarkScan and organisational recommendations).} \hldarkgreen{Designers are advised to adopt symmetric design principles so that opting out, cancelling, or exercising privacy preferences is no more complex than the corresponding opt\-in flows, and platform operators are encouraged to align product success metrics with user trust scores, task success rates, and long\-term retention rather than short\-term conversion alone.}\pdfcomment[icon=Note]{4807303 explicitly recommends symmetric design for opt‑out flows and aligning product metrics with trust, task success, and long‑term retention (4807303 p.10).} \hllightgreen{These prescriptions can be operationalized in MindfulFeed's business dashboards by prioritising metrics such as sustained daily active use without escalation of engagement, completion of consent and configuration tasks without confusion, and stable retention over time, accepting lower immediate conversion to paid features as the trade\-off for reputational resilience and compliance readiness.}\pdfcomment[icon=Note]{The specific dashboard metrics are a reasonable operationalisation drawn from recommendations to prioritise trust and well‑being metrics (4807303 p.10 on metrics; 9300898 p.7 on well‑being‑oriented AI), so this is supported synthesis.} \hldarkgreen{Publishing transparency reports on optimization techniques and maintaining full user journey documentation for key flows are further recommended as organisational practices that support regulatory scrutiny and trust \parencites[8]{4807303}.}\pdfcomment[icon=Note]{4807303 explicitly recommends publishing transparency reports and maintaining full user journey documentation to support regulatory scrutiny and trust (4807303 p.10--11).} \hllightgreen{Embedding these artefacts into the product development record contributes directly to business sustainability by reducing future compliance risks and fostering user confidence.}\pdfcomment[icon=Note]{Sources link transparency and documentation to reduced compliance risk and greater user trust (4807303 p.11; 9300898 p.7 on governance), making this a supported inference.} \hllightgreen{Privacy and AI governance research underscores that business models built on extensive behavioural surveillance and opaque personalization face structural risk.}\pdfcomment[icon=Note]{9300898 (p.7) and related governance literature (7808872 p.8 on oversight dimensions) describe structural risks for models built on behavioural surveillance and opaque personalization, supporting this synthesis.} \hldarkgreen{Contemporary platforms often employ complex privacy policies and confusing consent mechanisms that encourage unintended data disclosure, limiting users' understanding of how data are collected, processed, and shared \parencites[3]{2668071}.}\pdfcomment[icon=Note]{Multiple sources document complex privacy policies and confusing consent mechanisms that hinder understanding (2668071 p.3; 5678011 p.4 summarising GDPR consent requirements), directly supporting the claim.} \hldarkgreen{As personalization technologies become more sophisticated, they may amplify manipulative practices by tailoring deceptive interfaces to individual characteristics, prompting calls for stronger ethical guidelines and regulatory frameworks that prioritize welfare alongside business objectives \parencites[3]{2668071}[7]{9300898}.}\pdfcomment[icon=Note]{The claim that personalization can amplify manipulative practices and prompt calls for stronger guidelines is directly supported (2668071 p.3 on personalization amplifying manipulation; 9300898 p.7 calling for governance).} \hllightgreen{MindfulFeed's profitability strategy must therefore avoid monetization structures that depend on large\-scale, fine\-grained profiling without robust consent and transparency, because such structures are increasingly subject to regulatory challenge and reputational backlash.}\pdfcomment[icon=Note]{The recommendation to avoid monetisation based on fine‑grained profiling is a reasonable inference from governance and privacy risk discussions (9300898 p.7; 2668071 p.3 on reputational/regulatory risks).} \hllightgreen{Instead, sustainable revenue can be tied to clearly delimited, consented analytics that support preventive well\-being interventions, where machine learning models analyze patterns to identify concerns before they become severe, provided that monitoring remains proportionate and rights\-respecting \parencites[7]{9300898}.}\pdfcomment[icon=Note]{9300898 (p.7) documents AI for preventive well‑being interventions and cautions about proportionate, rights‑respecting monitoring, supporting this conditional proposal.} \hllightgreen{Systematic review methodology in computational burnout interventions suggests that organisations can manage complex, AI\-enabled services with high analytical rigor when structured processes for monitoring and evaluation are in place.}\pdfcomment[icon=Note]{8760542 describes systematic review methodology for computational interventions and argues that structured monitoring and evaluation enable rigorous management, so this is a supported synthesis (8760542 p.19).} \hldarkgreen{PRISMA\-guided reviews rely on staged identification, screening, eligibility, and inclusion, structured coding frameworks, quality appraisal tools, and inter\-rater reliability checks to maintain integrity.}\pdfcomment[icon=Note]{8760542 explicitly states PRISMA stages (identification, screening, eligibility, inclusion), coding frameworks, quality appraisal, and inter‑rater checks in its methods (8760542 p.19).} \hllightgreen{Translating this into ongoing business operations, MindfulFeed can support sustainability by instituting analogous internal review cycles for new features and monetization changes, screening proposed revenue mechanisms against inclusion criteria that cover ethical transparency, data minimization, and user well\-being impacts.}\pdfcomment[icon=Note]{Translating PRISMA‑style review cycles into internal feature review is a reasonable operational inference consistent with governance recommendations (8760542 p.19; 4807303 p.10 on ethical review processes).} \hllightgreen{Features failing quality thresholds would be excluded before launch, lowering the probability of costly post hoc remediation \parencites[19]{8760542}.}\pdfcomment[icon=Note]{The claim that excluding features that fail quality thresholds reduces costly remediation follows logically from the cited quality‑assurance and pre‑deployment review recommendations (8760542 p.19; 4807303 p.10).} \hldarkgreen{Continuous ethical oversight of recommendation and consent components, aligned with governance dimensions of accountability, transparency, oversight, and well\-being outcomes, further reduces long\-term legal and social risk, which is a central component of business sustainability in digital health and well\-being sectors \parencites[8]{7808872}[7]{9300898}.}\pdfcomment[icon=Note]{9300898 (p.7) and 7808872 (p.8) link continuous ethical oversight across accountability, transparency, oversight and well‑being outcomes to reduced legal and social risk in digital health contexts, matching this claim.} \hldarkgreen{Finally, studies on consumer awareness show that exposure to dark patterns is high but comprehension remains low, with many users misclassifying manipulative designs as mere bad design and knowing little about data\-driven consent manipulation \parencites[2]{5127387}.}\pdfcomment[icon=Note]{Empirical reviews report high exposure but low comprehension of dark patterns, with many users misclassifying manipulative designs as bad design and limited awareness of data‑driven consent manipulation (5127387 p.2; 4153678 summary findings).} \hllightgreen{While such low awareness may temporarily permit exploitative monetization elsewhere, aligning MindfulFeed with emerging regulatory standards that restrict deceptive consent practices and move toward mandatory dark\-pattern\-free certification in high\-risk domains positions the app advantageously in the long term \parencites[4]{5678011}[8]{4807303}.}\pdfcomment[icon=Note]{4807303 explicitly recommends mandatory 'dark‑pattern‑free' certification for high‑risk sectors and describes regulatory trends, so the inference that alignment positions the app advantageously is a supported synthesis (4807303 p.10).} \hllightgreen{Profitability, in this evidence\-based frame, becomes linked to user trust, regulatory compatibility, and demonstrable protection of autonomy, which collectively underpin durable business sustainability for an ethical, pro\-social social platform.}\pdfcomment[icon=Note]{Sources argue that trust, regulatory compatibility, and protecting autonomy are central to durable sustainability for well‑being platforms (9300898 p.7; 4807303 p.10), so this is a supported synthesis.}

\newpage
\printbibliography

\end{document}
